\chapter{Real Analysis}\label{real-analysis}
\chapterstatus{complete}

\section{Introduction}\label{real-analysis:section-introduction}

This chapter develops calculus in one and several real variables with complete proofs, starting slowly from limits of sequences and of functions: the sequential characterisation of limits,
continuous functions and the theorems of Bolzano, Weierstrass and Heine--Cantor, the mean value theorems, de l'H\^opital's rule, Taylor's formula and convexity;
then series and power series, the exponential
and trigonometric functions (including a definition of $\pi$), the Riemann integral, uniform convergence of sequences of functions (with a continuous nowhere differentiable function), derivatives of maps between Euclidean spaces, curves, their length and the line integrals of the first and second kind,
the inverse and implicit function theorems, extrema in several variables with Lagrange multipliers and the Karush--Kuhn--Tucker conditions of optimisation, ordinary differential equations, and the approximation and compactness theorems of
Stone--Weierstrass and Arzel\`a--Ascoli. The topological background (metric spaces, compactness, completeness, uniform convergence,
the Banach fixed point theorem) is in Chapter~\ref{general-topology}. References: \cite{rudin-principles}, \cite{dieudonne-foundations},
\cite{spivak-calculus-manifolds}, \cite{lang-real-functional} and \cite{hartman-ode}.

\noindent\textbf{Prerequisites.} Chapter~\ref{general-topology} (General Topology), Chapter~\ref{number-systems} (Number Systems) and
Chapter~\ref{linear-algebra} (Linear Algebra); Theorem~\ref{real-analysis:theorem-pi} also uses Chapter~\ref{homotopy} (Homotopy and Covering Spaces).

\noindent\textbf{On the way to the Hodge conjecture.} Needed: limits, continuity, series, differentiation, the exponential function, the Riemann integral, sequences of functions, derivatives in
several variables and the inverse and implicit function theorems; the length of curves and line integrals (Section~\ref{real-analysis:section-curves}) are used for paths and periods; ordinary differential equations (Section~\ref{real-analysis:section-ode}) give the flows of vector fields. Extrema and Lagrange multipliers
(Section~\ref{real-analysis:section-extrema}) and the Stone--Weierstrass and Arzel\`a--Ascoli theorems (Section~\ref{real-analysis:section-approximation}) can be postponed.

Throughout, $\mathbb{K}$ is $\RR$ or $\CC$, and $\RR^n$ and $\CC^n$ carry the Euclidean norm unless otherwise stated.

\section{Limits of sequences}\label{real-analysis:section-sequences}

We recall that a real sequence $(a_n)_{n \in \NN}$ \emph{converges} to $a \in \RR$ if for every $\varepsilon > 0$ there is $N$ with $|a_n - a| < \varepsilon$ for all $n \geq N$
(Definition~\ref{number-systems:definition-cauchy}). This is convergence in the metric space $\RR$ (Chapter~\ref{general-topology}). In this section we develop the theory of limits of sequences directly
from this definition.

\begin{definition}\label{real-analysis:definition-sequence-terms}
Let $(a_n)$ be a real sequence.
\begin{enumerate}
\item If $(a_n)$ converges to $a$ we write $a_n \to a$ or $\lim_na_n = a$, and call $a$ the \emph{limit}.
\item $(a_n)$ \emph{diverges to $+\infty$}, written $a_n \to +\infty$, if for every $M \in \RR$ there is $N$ with $a_n > M$ for all $n \geq N$; similarly $a_n \to -\infty$ if for every $M$ eventually $a_n < M$.
\item $(a_n)$ is \emph{regular} if it converges or diverges to $\pm\infty$, and \emph{oscillating} otherwise.
\item $(a_n)$ is \emph{bounded} if there is $C$ with $|a_n| \leq C$ for all $n$; \emph{bounded above} or \emph{below} are defined similarly.
\item $(a_n)$ is \emph{increasing} if $a_n \leq a_{n+1}$ for all $n$, \emph{strictly increasing} if $a_n < a_{n+1}$, and similarly \emph{decreasing}; it is \emph{monotone} if it is increasing or decreasing.
\item A \emph{subsequence} of $(a_n)$ is a sequence $(a_{n_k})_{k \in \NN}$ with $n_0 < n_1 < n_2 < \cdots$.
\end{enumerate}
A property holds \emph{eventually}, or \emph{for large $n$}, if there is $N$ such that it holds for all $n \geq N$.
\end{definition}

\begin{proposition}\label{real-analysis:proposition-sequence-basic}
Let $(a_n)$ be a real sequence.
\begin{enumerate}
\item A sequence has at most one limit.
\item A convergent sequence is bounded.
\item (Permanence of sign) If $a_n \to a$ and $a > 0$, then $a_n > a/2 > 0$ for large $n$.
\item Changing, adding or removing finitely many terms does not affect convergence or the limit.
\item If $a_n \to a$, every subsequence converges to $a$; the same holds for divergence to $\pm\infty$.
\end{enumerate}
\end{proposition}

\begin{proof}
(1) Suppose $a_n \to a$ and $a_n \to b$ with $a \neq b$, and let $\varepsilon = |a - b|/2$. For $n$ larger than both thresholds, $|a - b| \leq |a - a_n| + |a_n - b| < 2\varepsilon = |a - b|$, a contradiction.
(2) Take $\varepsilon = 1$: for $n \geq N$, $|a_n| \leq |a_n - a| + |a| < |a| + 1$, so $|a_n| \leq \max(|a_0|, \ldots, |a_{N-1}|, |a| + 1)$ for all $n$. (3) Take $\varepsilon = a/2$: then
$a_n > a - a/2 = a/2$ for large $n$. (4) The definitions only involve the terms with $n \geq N$, and $N$ may be enlarged. (5) By induction $n_k \geq k$, so $k \geq N$ implies $n_k \geq N$.
\end{proof}

\begin{lemma}[Bernoulli's inequality]\label{real-analysis:lemma-bernoulli}
For $h \geq -1$ and $n \in \NN$, $(1 + h)^n \geq 1 + nh$.
\end{lemma}

\begin{proof}
For $n = 0$ both sides are $1$. If $(1 + h)^n \geq 1 + nh$, then, as $1 + h \geq 0$, $(1 + h)^{n+1} \geq (1 + nh)(1 + h) = 1 + (n + 1)h + nh^2 \geq 1 + (n + 1)h$.
\end{proof}

\begin{proposition}[Algebra of limits]\label{real-analysis:proposition-algebra-limits}
Let $a_n \to a$ and $b_n \to b$ be real sequences and $c \in \RR$. Then $a_n + b_n \to a + b$, $ca_n \to ca$, $a_nb_n \to ab$ and $|a_n| \to |a|$. If $b \neq 0$, then $b_n \neq 0$ for large $n$ and
$a_n/b_n \to a/b$.
\end{proposition}

\begin{proof}
Given $\varepsilon > 0$, for large $n$ we have $|a_n - a| < \varepsilon/2$ and $|b_n - b| < \varepsilon/2$, so $|(a_n + b_n) - (a + b)| < \varepsilon$. The reverse triangle inequality
$||a_n| - |a|| \leq |a_n - a|$ gives $|a_n| \to |a|$, and $|ca_n - ca| = |c|\,|a_n - a|$ gives $ca_n \to ca$. For the product, let $C$ bound $|a_n|$
(Proposition~\ref{real-analysis:proposition-sequence-basic}(2)); then
\[
|a_nb_n - ab| \leq |a_n|\,|b_n - b| + |b|\,|a_n - a| \leq C|b_n - b| + |b|\,|a_n - a|,
\]
which is less than $\varepsilon$ for large $n$. For the quotient it suffices to show $1/b_n \to 1/b$. As $|b_n| \to |b| > 0$, permanence of sign gives $|b_n| > |b|/2$ for large $n$, and then
$|1/b_n - 1/b| = |b - b_n|/(|b_n|\,|b|) \leq 2|b_n - b|/|b|^2$, which tends to $0$.
\end{proof}

\begin{proposition}[Comparison and squeeze]\label{real-analysis:proposition-squeeze}
Let $(a_n)$, $(b_n)$, $(c_n)$ be real sequences.
\begin{enumerate}
\item If $a_n \leq b_n$ for large $n$, $a_n \to a$ and $b_n \to b$, then $a \leq b$.
\item (Squeeze theorem) If $a_n \leq c_n \leq b_n$ for large $n$ and $a_n \to L$, $b_n \to L$, then $c_n \to L$.
\item If $a_n \leq b_n$ for large $n$ and $a_n \to +\infty$, then $b_n \to +\infty$.
\item If $|a_n| \leq b_n$ for large $n$ and $b_n \to 0$, then $a_n \to 0$. In particular, if $(a_n)$ is bounded and $b_n \to 0$, then $a_nb_n \to 0$.
\end{enumerate}
\end{proposition}

\begin{proof}
(1) If $a > b$, then $a_n - b_n \to a - b > 0$ by Proposition~\ref{real-analysis:proposition-algebra-limits}, so $a_n - b_n > 0$ for large $n$ by permanence of sign, a contradiction.
(2) Given $\varepsilon > 0$, for large $n$ we have $L - \varepsilon < a_n \leq c_n \leq b_n < L + \varepsilon$. (3) Eventually $b_n \geq a_n > M$. (4) Apply (2) to $-b_n \leq a_n \leq b_n$; for the last
statement, $|a_nb_n| \leq C|b_n|$.
\end{proof}

\begin{counterexample}\label{real-analysis:counterexample-strict-inequality}
Strict inequalities do not pass to the limit: $1/(n + 1) > 0$ for all $n$, but the limit is $0$. The convergence $|a_n| \to |a|$ does not imply $a_n \to a$ unless $a = 0$: the sequence $a_n = (-1)^n$
oscillates although $|a_n| \to 1$.
\end{counterexample}

\begin{proposition}[Infinite limits]\label{real-analysis:proposition-infinite-limits}
Let $a_n \to +\infty$.
\begin{enumerate}
\item If $(b_n)$ is bounded below, then $a_n + b_n \to +\infty$.
\item If $b_n \geq c$ for large $n$, for some $c > 0$, then $a_nb_n \to +\infty$.
\item $a_n > 0$ for large $n$, and $1/a_n \to 0$. Conversely, if $b_n > 0$ and $b_n \to 0$, then $1/b_n \to +\infty$.
\end{enumerate}
\end{proposition}

\begin{proof}
(1) If $b_n \geq -C$, then for any $M$, eventually $a_n > M + C$, so $a_n + b_n > M$. (2) Eventually $a_n > M/c$ with $M > 0$, so $a_nb_n > M$. (3) Eventually $a_n > 1/\varepsilon > 0$, so
$0 < 1/a_n < \varepsilon$. Conversely, eventually $0 < b_n < 1/M$ for $M > 0$, so $1/b_n > M$.
\end{proof}

\begin{counterexample}[Indeterminate forms]\label{real-analysis:counterexample-indeterminate}
If $a_n \to +\infty$ and $b_n \to -\infty$, the limits alone do not determine the behaviour of $a_n + b_n$: with $a_n = n$, the choices $b_n = -n + c$, $b_n = -n/2$, $b_n = -2n$ and $b_n = -n + (-1)^n$ give
the limit $c$, the limit $+\infty$, the limit $-\infty$, and an oscillating sequence. Similarly, if $a_n \to 0$ and $b_n \to +\infty$, then with $a_n = 1/(n+1)$ the choices $b_n = c(n + 1)$, $(n + 1)^2$ and
$(-1)^n(n + 1)$ show that $a_nb_n$ can converge to any $c$, diverge, or oscillate. The same happens for the forms $\infty/\infty$ and $0/0$.
\end{counterexample}

\begin{proposition}[Fundamental limits]\label{real-analysis:proposition-fundamental-limits}
\begin{enumerate}
\item For $q \in \RR$: $q^n \to 0$ if $|q| < 1$, $q^n \to 1$ if $q = 1$, $q^n \to +\infty$ if $q > 1$, and $(q^n)$ oscillates if $q \leq -1$.
\item For $a > 0$, $a^{1/n} \to 1$, and $n^{1/n} \to 1$.
\item For $k \in \NN$ and $a > 1$, $n^k/a^n \to 0$: exponentials beat powers.
\item For every $a \in \RR$, $a^n/n! \to 0$, and $n!/n^n \to 0$.
\end{enumerate}
\end{proposition}

\begin{proof}
(1) If $q > 1$, write $q = 1 + h$ with $h > 0$; by Lemma~\ref{real-analysis:lemma-bernoulli}, $q^n \geq 1 + nh \to +\infty$. If $0 < |q| < 1$, then $|q|^{-n} \to +\infty$, so $|q^n| \to 0$ by
Proposition~\ref{real-analysis:proposition-infinite-limits}(3), and $q^n \to 0$ by Proposition~\ref{real-analysis:proposition-squeeze}(4). If $q = -1$ the terms alternate between $\pm1$; if $q < -1$, $|q^n| \to +\infty$
and the signs alternate, so the sequence is unbounded above and below and has no limit.
(2) Let $a \geq 1$ and $a^{1/n} = 1 + h_n$ with $h_n \geq 0$ (the $n$-th root exists by Proposition~\ref{number-systems:proposition-roots}). By Bernoulli, $a \geq 1 + nh_n$, so $0 \leq h_n \leq (a - 1)/n \to 0$. For
$a < 1$ apply this to $1/a$ and use Proposition~\ref{real-analysis:proposition-algebra-limits}. For $n^{1/n} = 1 + h_n$ with $n \geq 2$, the binomial formula gives $n \geq \binom{n}{2}h_n^2$, so
$0 \leq h_n \leq \sqrt{2/(n - 1)} \to 0$.
(3) Write $a = 1 + h$. For $n \geq 2k + 2$ the binomial formula gives $a^n \geq \binom{n}{k+1}h^{k+1} \geq \frac{(n/2)^{k+1}}{(k+1)!}h^{k+1}$, since each of the $k + 1$ factors $n, n - 1, \ldots, n - k$ is at least $n/2$.
Hence $0 \leq n^k/a^n \leq 2^{k+1}(k + 1)!/(h^{k+1}n) \to 0$.
(4) Choose an integer $N > 2|a|$. For $n \geq N$, $|a|^n/n! \leq \frac{|a|^N}{N!}2^{N-n}$, which tends to $0$ by (1). Finally $n!/n^n = \frac1n\cdot\frac2n\cdots\frac nn \leq \frac1n$.
\end{proof}

\begin{proposition}\label{real-analysis:proposition-monotone-regular}
Every monotone sequence is regular. An increasing sequence converges to its supremum if it is bounded above, and diverges to $+\infty$ otherwise; a decreasing sequence converges to its infimum if it is
bounded below, and diverges to $-\infty$ otherwise.
\end{proposition}

\begin{proof}
Let $(a_n)$ be increasing. If it is bounded above, let $s = \sup_na_n$ (Definition~\ref{number-systems:definition-complete}). For $\varepsilon > 0$, $s - \varepsilon$ is not an upper bound, so $a_N > s - \varepsilon$ for
some $N$, and then $s - \varepsilon < a_N \leq a_n \leq s$ for $n \geq N$. If it is not bounded above, for every $M$ some $a_N > M$, and then $a_n \geq a_N > M$ for $n \geq N$. The decreasing case follows by
passing to $(-a_n)$.
\end{proof}

\begin{lemma}\label{real-analysis:lemma-monotone-subsequence}
Every real sequence has a monotone subsequence.
\end{lemma}

\begin{proof}
Call $n$ a \emph{peak} if $a_n \geq a_m$ for all $m > n$. If there are infinitely many peaks $n_0 < n_1 < \cdots$, then $a_{n_0} \geq a_{n_1} \geq \cdots$ is a decreasing subsequence. Otherwise let $n_0$ be larger than
all peaks. As $n_0$ is not a peak, there is $n_1 > n_0$ with $a_{n_1} > a_{n_0}$; as $n_1$ is not a peak, there is $n_2 > n_1$ with $a_{n_2} > a_{n_1}$; continuing (by recursion), we obtain a strictly increasing
subsequence.
\end{proof}

\begin{theorem}[Bolzano--Weierstrass]\label{real-analysis:theorem-bolzano-weierstrass-real}
Every bounded real sequence has a convergent subsequence. Every real sequence has a regular subsequence: a sequence that is unbounded above has a subsequence diverging to $+\infty$.
\end{theorem}

\begin{proof}
By Lemma~\ref{real-analysis:lemma-monotone-subsequence} there is a monotone subsequence, which is regular by Proposition~\ref{real-analysis:proposition-monotone-regular}, and convergent if the sequence is
bounded. If $(a_n)$ is unbounded above, choose recursively $n_{k+1} > n_k$ with $a_{n_{k+1}} > k + 1$, which is possible because only finitely many terms have index $\leq n_k$.
\end{proof}

\begin{theorem}[Cauchy criterion]\label{real-analysis:theorem-cauchy-criterion}
A real sequence converges if and only if it is a Cauchy sequence.
\end{theorem}

\begin{proof}
If $a_n \to a$, then $|a_n - a_m| \leq |a_n - a| + |a - a_m| < \varepsilon$ when both are within $\varepsilon/2$ of $a$. Conversely, let $(a_n)$ be Cauchy. Taking $\varepsilon = 1$ shows that $(a_n)$ is bounded,
so by Theorem~\ref{real-analysis:theorem-bolzano-weierstrass-real} a subsequence converges to some $a$. Given $\varepsilon > 0$, choose $N$ with $|a_n - a_m| < \varepsilon/2$ for $m, n \geq N$, and $k$ with $n_k \geq N$
and $|a_{n_k} - a| < \varepsilon/2$. Then $|a_n - a| \leq |a_n - a_{n_k}| + |a_{n_k} - a| < \varepsilon$ for $n \geq N$. (This is a second proof of
Proposition~\ref{number-systems:proposition-cauchy-complete} for $\RR$.)
\end{proof}

\begin{counterexample}\label{real-analysis:counterexample-consecutive-differences}
It is not enough that consecutive terms get close. The partial sums $H_n = \sum_{k=1}^n1/k$ of the harmonic series satisfy $H_{n+1} - H_n = 1/(n + 1) \to 0$, but
$H_{2n} - H_n = \sum_{k=n+1}^{2n}1/k \geq n \cdot \frac{1}{2n} = \frac12$, so $(H_n)$ is not Cauchy; being increasing, it diverges to $+\infty$. Likewise $\sqrt{n + 1} - \sqrt n = 1/(\sqrt{n + 1} + \sqrt n) \to 0$
while $\sqrt n \to +\infty$.
\end{counterexample}

\begin{proposition}[Subsequence criterion]\label{real-analysis:proposition-subsequence-criterion}
\begin{enumerate}
\item $a_n \to a$ if and only if every subsequence of $(a_n)$ has a further subsequence converging to $a$.
\item A bounded sequence all of whose convergent subsequences have the same limit $a$ converges to $a$.
\end{enumerate}
\end{proposition}

\begin{proof}
(1) One direction is Proposition~\ref{real-analysis:proposition-sequence-basic}(5). Conversely, if $a_n \not\to a$, there is $\varepsilon > 0$ such that $|a_n - a| \geq \varepsilon$ for infinitely many $n$; these terms form a
subsequence none of whose subsequences converges to $a$. (2) By Theorem~\ref{real-analysis:theorem-bolzano-weierstrass-real}, every subsequence is bounded and has a convergent subsequence, whose limit is $a$ by
hypothesis; apply (1).
\end{proof}

\begin{proposition}[Stolz--Ces\`aro]\label{real-analysis:proposition-stolz-cesaro}
Let $(b_n)$ be strictly increasing with $b_n \to +\infty$, and let $(a_n)$ be a real sequence with $\frac{a_{n+1} - a_n}{b_{n+1} - b_n} \to L \in \RR$. Then $a_n/b_n \to L$. In particular (Ces\`aro), if $a_n \to L$, then
the averages $\frac{1}{n}(a_1 + \cdots + a_n) \to L$.
\end{proposition}

\begin{proof}
Let $\varepsilon > 0$ and choose $N$ such that $b_N > 0$ and $|(a_{k+1} - a_k) - L(b_{k+1} - b_k)| \leq \varepsilon(b_{k+1} - b_k)$ for $k \geq N$. Summing over $N \leq k < n$ gives
$|a_n - a_N - L(b_n - b_N)| \leq \varepsilon(b_n - b_N) \leq \varepsilon b_n$. Dividing by $b_n > 0$,
\[
\Big|\frac{a_n}{b_n} - L\Big| \leq \varepsilon + \frac{|a_N - Lb_N|}{b_n},
\]
and the last term is less than $\varepsilon$ for large $n$, as $b_n \to +\infty$. For Ces\`aro, apply this to $a'_n = a_1 + \cdots + a_n$ and $b_n = n$.
\end{proof}

\begin{counterexample}\label{real-analysis:counterexample-cesaro-converse}
The converse of Ces\`aro's theorem fails: $a_n = (-1)^n$ does not converge, but its averages tend to $0$, since the partial sums are $0$ or $-1$.
\end{counterexample}

\section{Limits of functions}\label{real-analysis:section-limits-functions}

To treat finite and infinite limits at once we use the \emph{extended real line} $\bar{\RR} = \RR \cup \{-\infty, +\infty\}$, ordered so that $-\infty < x < +\infty$ for $x \in \RR$.

\begin{definition}\label{real-analysis:definition-neighbourhoods}
For $a \in \bar{\RR}$ and $\delta > 0$, the \emph{basic neighbourhood} $U_\delta(a)$ is $(a - \delta, a + \delta)$ if $a \in \RR$, $(1/\delta, +\infty)$ if $a = +\infty$, and $(-\infty, -1/\delta)$ if $a = -\infty$. Let
$D \subseteq \RR$. A point $a \in \bar{\RR}$ is an \emph{accumulation point} of $D$ if every $U_\delta(a)$ contains a point of $D \setminus \{a\}$. A point of $D$ that is not an accumulation point is \emph{isolated}.
\end{definition}

A sequence $(x_n)$ tends to $a \in \bar{\RR}$ if and only if for every $\delta > 0$, eventually $x_n \in U_\delta(a)$: for $a \in \RR$ this is convergence, and for $a = \pm\infty$ divergence to $\pm\infty$. Every $a \in \bar{\RR}$
that is an accumulation point of $D$ is the limit of a sequence in $D \setminus \{a\}$: choose $x_n \in U_{1/(n+1)}(a) \cap D \setminus \{a\}$.

\begin{definition}\label{real-analysis:definition-function-limit}
Let $f \colon D \to \RR$, let $a \in \bar{\RR}$ be an accumulation point of $D$, and $L \in \bar{\RR}$. We write $\lim_{x \to a}f(x) = L$, or $f(x) \to L$ as $x \to a$, if for every $\varepsilon > 0$ there is $\delta > 0$ such that
$f(x) \in U_\varepsilon(L)$ for all $x \in D \cap U_\delta(a)$ with $x \neq a$. For $a, L \in \RR$ this reads: for every $\varepsilon > 0$ there is $\delta > 0$ with $|f(x) - L| < \varepsilon$ whenever $x \in D$ and
$0 < |x - a| < \delta$. The \emph{right limit} $\lim_{x \to a^+}f(x)$ is the limit of the restriction of $f$ to $D \cap (a, +\infty)$, when $a$ is an accumulation point of that set, and the \emph{left limit}
$\lim_{x \to a^-}f(x)$ is defined similarly. The same definitions apply to $\CC$- or $\RR^m$-valued functions and finite limits, with $|\cdot|$ the modulus or the Euclidean norm.
\end{definition}

The value of $f$ at $a$, if $a \in D$, plays no role in the definition.

\begin{proposition}\label{real-analysis:proposition-function-limit-basic}
Let $f \colon D \to \RR$ and let $a \in \bar{\RR}$ be an accumulation point of $D$.
\begin{enumerate}
\item The limit of $f$ at $a$, if it exists, is unique.
\item If $\lim_{x \to a}f(x) = L \in \RR$, then $f$ is bounded on $D \cap U_\delta(a) \setminus \{a\}$ for some $\delta > 0$; if $L > 0$, then $f(x) > L/2$ for $x \in D \cap U_\delta(a) \setminus \{a\}$ for some $\delta > 0$.
\item If $a \in \RR$ is an accumulation point of both $D \cap (a, +\infty)$ and $D \cap (-\infty, a)$, then $\lim_{x \to a}f(x) = L$ if and only if both one-sided limits exist and equal $L$.
\end{enumerate}
\end{proposition}

\begin{proof}
(1) If $L \neq L'$ are both limits, choose $\varepsilon$ with $U_\varepsilon(L) \cap U_\varepsilon(L') = \emptyset$, and $\delta$ working for both; since $a$ is an accumulation point, some $x \in D \cap U_\delta(a) \setminus \{a\}$
exists, and $f(x)$ lies in both neighbourhoods. (2) Take $\varepsilon = 1$, respectively $\varepsilon = L/2$. (3) If both one-sided limits equal $L$, take the smaller of the two $\delta$; the converse is clear.
\end{proof}

\begin{theorem}[Sequential characterisation of limits]\label{real-analysis:theorem-sequential-limits}
Let $f \colon D \to \RR$, let $a \in \bar{\RR}$ be an accumulation point of $D$, and $L \in \bar{\RR}$. The following are equivalent:
\begin{enumerate}
\item $\lim_{x \to a}f(x) = L$;
\item for every sequence $(x_n)$ in $D \setminus \{a\}$ with $x_n \to a$, we have $f(x_n) \to L$.
\end{enumerate}
\end{theorem}

\begin{proof}
(1) $\Rightarrow$ (2). Let $\varepsilon > 0$ and choose $\delta$ as in Definition~\ref{real-analysis:definition-function-limit}. Eventually $x_n \in U_\delta(a)$, and $x_n \in D \setminus \{a\}$, so eventually
$f(x_n) \in U_\varepsilon(L)$.

(2) $\Rightarrow$ (1). Suppose (1) fails. Then there is $\varepsilon > 0$ such that for every $\delta > 0$ some $x \in D \cap U_\delta(a) \setminus \{a\}$ has $f(x) \notin U_\varepsilon(L)$. Choosing such a point $x_n$ for
$\delta = 1/(n + 1)$, for each $n$ (this uses the axiom of countable choice), we obtain a sequence in $D \setminus \{a\}$ with $x_n \to a$ and $f(x_n) \notin U_\varepsilon(L)$ for all $n$, so $f(x_n) \not\to L$,
contradicting (2).
\end{proof}

Theorem~\ref{real-analysis:theorem-sequential-limits}, sometimes called the \emph{bridge theorem}, transfers every result about limits of sequences to limits of functions, and it is the standard way to
show that a limit does \emph{not} exist: exhibit two sequences tending to $a$ along which $f$ has different limits.

\begin{corollary}[Rules for limits of functions]\label{real-analysis:corollary-limit-rules}
Let $f, g, h \colon D \to \RR$ and let $a \in \bar{\RR}$ be an accumulation point of $D$.
\begin{enumerate}
\item If $f(x) \to L$ and $g(x) \to M$ with $L, M \in \RR$, then $f + g \to L + M$, $fg \to LM$, and, if $M \neq 0$, $f/g \to L/M$ (where $f/g$ is defined near $a$).
\item If $f \leq g$ near $a$ and $f \to L$, $g \to M$, then $L \leq M$.
\item (Squeeze) If $f \leq h \leq g$ near $a$ and $f(x) \to L$, $g(x) \to L$, then $h(x) \to L$.
\item The rules for infinite limits of Proposition~\ref{real-analysis:proposition-infinite-limits} hold for functions.
\end{enumerate}
\end{corollary}

\begin{proof}
Apply Theorem~\ref{real-analysis:theorem-sequential-limits} and Propositions~\ref{real-analysis:proposition-algebra-limits}, \ref{real-analysis:proposition-squeeze} and
\ref{real-analysis:proposition-infinite-limits} to $f(x_n)$, $g(x_n)$, $h(x_n)$ for sequences $x_n \to a$ in $D \setminus \{a\}$. For the quotient, $g(x) \neq 0$ near $a$ by
Proposition~\ref{real-analysis:proposition-function-limit-basic}(2) applied to $|g|$.
\end{proof}

\begin{proposition}[Limits of composite functions]\label{real-analysis:proposition-composite-limit}
Let $f \colon D \to \RR$ and $g \colon E \to \RR$ with $f(D) \subseteq E$. Suppose $\lim_{x \to a}f(x) = b$ and $\lim_{y \to b}g(y) = L$, and that one of the following holds:
\begin{enumerate}
\item $f(x) \neq b$ for all $x \in D \setminus \{a\}$ near $a$;
\item $b \in E$ and $g(b) = L$.
\end{enumerate}
Then $\lim_{x \to a}g(f(x)) = L$.
\end{proposition}

\begin{proof}
Let $x_n \to a$ in $D \setminus \{a\}$. Then $f(x_n) \to b$ by Theorem~\ref{real-analysis:theorem-sequential-limits}. In case (1), eventually $f(x_n) \in E \setminus \{b\}$, so $g(f(x_n)) \to L$ by the same theorem. In
case (2), given $\varepsilon > 0$ choose $\eta$ with $g(y) \in U_\varepsilon(L)$ for $y \in E \cap U_\eta(b)$, including $y = b$; eventually $f(x_n) \in U_\eta(b)$, so $g(f(x_n)) \in U_\varepsilon(L)$. Apply the theorem again.
\end{proof}

\begin{counterexample}\label{real-analysis:counterexample-composite}
Neither hypothesis can be dropped. Let $f(x) = 0$ for all $x$, and $g(y) = 0$ for $y \neq 0$, $g(0) = 1$. Then $\lim_{x \to 0}f(x) = 0$ and $\lim_{y \to 0}g(y) = 0$, but $g(f(x)) = 1$ for all $x$.
\end{counterexample}

\begin{proposition}[Cauchy criterion for functions]\label{real-analysis:proposition-cauchy-functions}
Let $f \colon D \to \RR$ and let $a \in \bar{\RR}$ be an accumulation point of $D$. Then $\lim_{x \to a}f(x)$ exists in $\RR$ if and only if for every $\varepsilon > 0$ there is $\delta > 0$ with
$|f(x) - f(y)| < \varepsilon$ for all $x, y \in D \cap U_\delta(a) \setminus \{a\}$.
\end{proposition}

\begin{proof}
If $f(x) \to L$, take $\delta$ for $\varepsilon/2$ and use the triangle inequality. Conversely, for every sequence $x_n \to a$ in $D \setminus \{a\}$, the sequence $(f(x_n))$ is Cauchy, hence converges
(Theorem~\ref{real-analysis:theorem-cauchy-criterion}). If $x_n \to a$ and $y_n \to a$ are two such sequences, the interleaved sequence $x_0, y_0, x_1, y_1, \ldots$ also tends to $a$, so $(f(x_n))$ and $(f(y_n))$ have the
same limit, being subsequences of a convergent sequence. So all these limits equal a common $L$, and $f(x) \to L$ by Theorem~\ref{real-analysis:theorem-sequential-limits}.
\end{proof}

\begin{proposition}[Monotone functions]\label{real-analysis:proposition-monotone-limits}
Let $f \colon (a, b) \to \RR$ be increasing, with $-\infty \leq a < b \leq +\infty$. For every $c \in (a, b)$ the one-sided limits exist and
\[
\lim_{x \to c^-}f(x) = \sup_{a < x < c}f(x) \leq f(c) \leq \inf_{c < x < b}f(x) = \lim_{x \to c^+}f(x),
\]
and $\lim_{x \to b^-}f(x) = \sup f \in \RR \cup \{+\infty\}$, $\lim_{x \to a^+}f(x) = \inf f \in \RR \cup \{-\infty\}$. The analogous statements hold for decreasing functions.
\end{proposition}

\begin{proof}
The set $\{f(x) : a < x < c\}$ is bounded above by $f(c)$; let $s$ be its supremum, so $s \leq f(c)$. Given $\varepsilon > 0$ there is $x_0 < c$ with $f(x_0) > s - \varepsilon$, and for $x_0 < x < c$,
$s - \varepsilon < f(x_0) \leq f(x) \leq s$. So $\lim_{x \to c^-}f(x) = s$. The other statements are proved in the same way (for $b = +\infty$ or an unbounded supremum, use neighbourhoods of $+\infty$), and the decreasing
case follows by passing to $-f$.
\end{proof}

\begin{counterexample}\label{real-analysis:counterexample-no-limit}
\begin{enumerate}
\item The sign function $\operatorname{sgn}(x) = x/|x|$ on $\RR \setminus \{0\}$ has one-sided limits $\pm1$ at $0$, and no limit there.
\item $1/x$ has no limit at $0$: its one-sided limits are $+\infty$ and $-\infty$.
\item The \emph{Dirichlet function} $\mathbf{1}_\QQ$, equal to $1$ on $\QQ$ and $0$ on $\RR \setminus \QQ$, has no limit at any $a \in \RR$. Both $\QQ$ and $\RR \setminus \QQ$ are dense in $\RR$: $\QQ$ by
Proposition~\ref{number-systems:proposition-complete-archimedean}, and $\RR \setminus \QQ$ because $q\sqrt2$ is irrational for rational $q \neq 0$ (Proposition~\ref{number-systems:proposition-sqrt-two}) and
$x/\sqrt2 < q < y/\sqrt2$ is possible for $x < y$. So there are rational $x_n \to a$ and irrational $y_n \to a$, all $\neq a$, and $\mathbf{1}_\QQ(x_n) \to 1$, $\mathbf{1}_\QQ(y_n) \to 0$.
\end{enumerate}
\end{counterexample}

\section{Continuous functions of a real variable}\label{real-analysis:section-continuity-real}

\begin{definition}\label{real-analysis:definition-continuity-real}
Let $D \subseteq \RR$ and $f \colon D \to \RR$. Then $f$ is \emph{continuous at} $a \in D$ if for every $\varepsilon > 0$ there is $\delta > 0$ with $|f(x) - f(a)| < \varepsilon$ for all $x \in D$ with
$|x - a| < \delta$, and \emph{continuous} if it is continuous at every point of $D$.
\end{definition}

This is continuity for the metric of $\RR$ restricted to $D$ (Proposition~\ref{general-topology:proposition-metric-continuity}). If $a \in D$ is an accumulation point of $D$, then $f$ is continuous at $a$ if and
only if $\lim_{x \to a}f(x) = f(a)$; if $a$ is isolated, every $f$ is continuous at $a$.

\begin{proposition}\label{real-analysis:proposition-continuity-sequential}
Let $f, g \colon D \to \RR$ and $a \in D$.
\begin{enumerate}
\item $f$ is continuous at $a$ if and only if $f(x_n) \to f(a)$ for every sequence $(x_n)$ in $D$ with $x_n \to a$.
\item If $f$ and $g$ are continuous at $a$, so are $f + g$, $fg$, $|f|$, $\max(f, g)$, $\min(f, g)$, and $f/g$ if $g(a) \neq 0$.
\item If $f$ is continuous at $a$ and $h \colon E \to \RR$, with $f(D) \subseteq E$, is continuous at $f(a)$, then $h \circ f$ is continuous at $a$.
\item Polynomials are continuous on $\RR$, and rational functions are continuous where the denominator does not vanish.
\end{enumerate}
\end{proposition}

\begin{proof}
(1) The proof of Theorem~\ref{real-analysis:theorem-sequential-limits} applies verbatim, now allowing $x_n = a$. (2) follows from (1) and Proposition~\ref{real-analysis:proposition-algebra-limits}, since
$\max(f, g) = \frac12(f + g + |f - g|)$ and $\min(f, g) = \frac12(f + g - |f - g|)$. (3) If $x_n \to a$, then $f(x_n) \to f(a)$ and so $h(f(x_n)) \to h(f(a))$. (4) Constants and the identity are continuous;
apply (2).
\end{proof}

\begin{definition}\label{real-analysis:definition-discontinuities}
Let $f$ be defined on an open interval containing $a$, and discontinuous at $a$. The discontinuity is \emph{removable} if $\lim_{x \to a}f(x)$ exists in $\RR$ (and differs from $f(a)$); it is a \emph{jump} if both
one-sided limits exist in $\RR$ and differ; otherwise it is \emph{essential}.
\end{definition}

\begin{example}\label{real-analysis:example-discontinuities}
The function equal to $1$ for $x \neq 0$ and $0$ at $0$ has a removable discontinuity at $0$. The function equal to $\operatorname{sgn}(x)$ for $x \neq 0$ and $0$ at $0$ has a jump. The function equal to $1/x$ for
$x \neq 0$ and $0$ at $0$ has an essential discontinuity, and the Dirichlet function $\mathbf{1}_\QQ$ has an essential discontinuity at every point (Counterexample~\ref{real-analysis:counterexample-no-limit}).
\end{example}

\begin{theorem}\label{real-analysis:theorem-monotone-discontinuities}
Let $f \colon I \to \RR$ be monotone on an interval $I$. Then every discontinuity of $f$ at an interior point of $I$ is a jump, and the set of points of discontinuity of $f$ is countable.
\end{theorem}

\begin{proof}
We may assume $f$ increasing. At an interior point $c$ the one-sided limits $f(c^-) \leq f(c) \leq f(c^+)$ exist by Proposition~\ref{real-analysis:proposition-monotone-limits}, and $f$ is continuous at $c$ if and only if
$f(c^-) = f(c^+)$. So every interior discontinuity is a jump, with nonempty open interval $J_c = (f(c^-), f(c^+))$. If $c < c'$ are interior discontinuity points and $c < x < c'$, then
$f(c^+) \leq f(x) \leq f(c'^-)$, so $J_c \cap J_{c'} = \emptyset$. The set $\QQ$ is countable (see the proof of Proposition~\ref{number-systems:proposition-reals-cardinality}); fix an injection
$\iota \colon \QQ \to \NN$ and send $c$ to the rational number $q \in J_c$ with the smallest $\iota(q)$, which exists as $J_c$ contains rational numbers. This map is injective, since the $J_c$ are disjoint, so
the interior discontinuities form a countable set. Adding the at most two endpoints of $I$ keeps it countable.
\end{proof}

\begin{example}[Thomae's function]\label{real-analysis:example-thomae}
Let $T(x) = 1/q$ if $x = p/q$ with $p \in \ZZ$, $q \geq 1$ coprime, and $T(x) = 0$ for irrational $x$. Then $T$ is continuous at every irrational number and discontinuous at every rational number.

Indeed, at a rational $a$, $T(a) > 0$, but there are irrational $y_n \to a$ (Counterexample~\ref{real-analysis:counterexample-no-limit}) with $T(y_n) = 0$. At an irrational $a$, let $\varepsilon > 0$. The rationals
$p/q$ in $(a - 1, a + 1)$ with $q \leq 1/\varepsilon$ are finitely many, and none equals $a$; let $\delta$ be smaller than $1$ and than their distances to $a$. For $|x - a| < \delta$, either $x$ is irrational and $T(x) = 0$,
or $x = p/q$ with $q > 1/\varepsilon$, so $|T(x) - T(a)| = 1/q < \varepsilon$.
\end{example}

\begin{remark}\label{real-analysis:remark-continuity-set}
There is no function $\RR \to \RR$ that is continuous exactly at the rational numbers: the set of points of continuity of any function is a countable intersection of open sets, and $\QQ$ is not, by the
Baire category theorem (Theorem~\ref{general-topology:theorem-baire}). By contrast, $x\mathbf{1}_\QQ(x)$ is continuous exactly at $0$ (Exercise~\ref{real-analysis:exercise-continuous-one-point}).
\end{remark}

\begin{theorem}[Bolzano's theorem on zeros]\label{real-analysis:theorem-bolzano-zeros}
Let $f \colon [a, b] \to \RR$ be continuous with $f(a) < 0 < f(b)$. Then $f(c) = 0$ for some $c \in (a, b)$.
\end{theorem}

\begin{proof}
We construct intervals by bisection. Let $a_0 = a$, $b_0 = b$. Given $[a_n, b_n]$ with $f(a_n) < 0 \leq f(b_n)$, let $m = (a_n + b_n)/2$; if $f(m) < 0$ put $a_{n+1} = m$, $b_{n+1} = b_n$, and otherwise $a_{n+1} = a_n$,
$b_{n+1} = m$. Then $f(a_{n+1}) < 0 \leq f(b_{n+1})$, $(a_n)$ is increasing, $(b_n)$ is decreasing, and $b_n - a_n = (b - a)/2^n \to 0$. By Proposition~\ref{real-analysis:proposition-monotone-regular} both converge,
to the same limit $c \in [a, b]$. By continuity, $f(c) = \lim f(a_n) \leq 0$ and $f(c) = \lim f(b_n) \geq 0$ (Proposition~\ref{real-analysis:proposition-squeeze}(1)), so $f(c) = 0$, and $c \neq a, b$.
\end{proof}

\begin{corollary}\label{real-analysis:corollary-ivt-real}
Let $I$ be an interval and $f \colon I \to \RR$ continuous.
\begin{enumerate}
\item (Intermediate value theorem) $f$ takes every value between any two of its values, so $f(I)$ is an interval.
\item If $f$ is injective, it is strictly monotone.
\item Every continuous map $f \colon [a, b] \to [a, b]$ has a fixed point.
\item Every real polynomial of odd degree has a real root.
\end{enumerate}
\end{corollary}

\begin{proof}
(1) If $f(x) < y < f(z)$, apply Theorem~\ref{real-analysis:theorem-bolzano-zeros} to $f - y$ on the interval between $x$ and $z$ (to $y - f$ if $x > z$). This is also the case $X = I$ of
Corollary~\ref{general-topology:corollary-ivt}.
(2) Let $a < b$ in $I$, and assume $f(a) < f(b)$ (otherwise replace $f$ by $-f$). Let $x < y$ in $I$ and put $h(t) = f((1 - t)b + ty) - f((1 - t)a + tx)$ for $t \in [0, 1]$. Since $(1 - t)a + tx < (1 - t)b + ty$ and $f$ is injective,
$h(t) \neq 0$ for all $t$. As $h$ is continuous and $h(0) > 0$, (1) gives $h(1) = f(y) - f(x) > 0$. So $f$ is strictly increasing.
(3) $g(x) = f(x) - x$ satisfies $g(a) \geq 0 \geq g(b)$; if neither is $0$, apply (1).
(4) We may take $p(x) = x^{2k+1} + c_{2k}x^{2k} + \cdots + c_0$. For $|x| \geq 1 + \sum|c_i|$ we have $|c_{2k}x^{2k} + \cdots + c_0| \leq \sum|c_i|\,|x|^{2k} < |x|^{2k+1}$, so $p(x)$ has the sign of $x$ there. Apply (1).
\end{proof}

\begin{theorem}[Weierstrass]\label{real-analysis:theorem-weierstrass-extreme}
A continuous function $f \colon [a, b] \to \RR$ is bounded and attains its maximum and its minimum.
\end{theorem}

\begin{proof}
If $f$ were unbounded, there would be $x_n \in [a, b]$ with $|f(x_n)| > n$. By Theorem~\ref{real-analysis:theorem-bolzano-weierstrass-real} a subsequence converges to some $c$, and $c \in [a, b]$ by
Proposition~\ref{real-analysis:proposition-squeeze}(1). Then $f(x_{n_k}) \to f(c)$, so $(f(x_{n_k}))$ is bounded, a contradiction. Let $s = \sup f$ and choose $x_n$ with $f(x_n) > s - 1/(n + 1)$; a subsequence converges to
some $c \in [a, b]$, and $f(c) = \lim f(x_{n_k}) = s$ by the squeeze theorem. The minimum is the maximum of $-f$. (This is the case $X = [a, b]$ of Corollary~\ref{general-topology:corollary-extreme-value}.)
\end{proof}

\begin{counterexample}\label{real-analysis:counterexample-extreme-value}
Each hypothesis of Theorem~\ref{real-analysis:theorem-weierstrass-extreme} is needed: $x$ on $(0, 1)$ attains neither its supremum nor its infimum; $1/x$ on $(0, 1]$ is unbounded; $x$ on $[0, +\infty)$ is
unbounded; and the function on $[0, 1]$ equal to $x$ for $x < 1$ and to $0$ at $1$ has supremum $1$, which is not attained.
\end{counterexample}

\begin{definition}\label{real-analysis:definition-uniform-continuity-real}
A function $f \colon D \to \RR$ is \emph{uniformly continuous} if for every $\varepsilon > 0$ there is $\delta > 0$ with $|f(x) - f(y)| < \varepsilon$ for all $x, y \in D$ with $|x - y| < \delta$. It is
\emph{Lipschitz} with constant $K \geq 0$ if $|f(x) - f(y)| \leq K|x - y|$ for all $x, y \in D$.
\end{definition}

\begin{proposition}\label{real-analysis:proposition-uniform-continuity-real}
Let $f \colon D \to \RR$.
\begin{enumerate}
\item Lipschitz functions are uniformly continuous, and uniformly continuous functions are continuous.
\item $f$ is uniformly continuous if and only if $f(x_n) - f(y_n) \to 0$ for all sequences $(x_n)$, $(y_n)$ in $D$ with $x_n - y_n \to 0$.
\item (Heine--Cantor) If $D = [a, b]$ and $f$ is continuous, then $f$ is uniformly continuous.
\item A uniformly continuous function maps Cauchy sequences to Cauchy sequences. A uniformly continuous $f \colon (a, b) \to \RR$ on a bounded interval extends to a continuous function on $[a, b]$.
\end{enumerate}
\end{proposition}

\begin{proof}
(1) Take $\delta = \varepsilon/K$ (any $\delta$ if $K = 0$); the second statement is clear. (2) If $f$ is uniformly continuous and $x_n - y_n \to 0$, eventually $|x_n - y_n| < \delta$ and so $|f(x_n) - f(y_n)| < \varepsilon$.
If $f$ is not uniformly continuous, there is $\varepsilon > 0$ and, for each $n$, points $x_n, y_n$ with $|x_n - y_n| < 1/(n + 1)$ and $|f(x_n) - f(y_n)| \geq \varepsilon$. (3) If not, take $x_n$, $y_n$ as in the last
sentence. By Theorem~\ref{real-analysis:theorem-bolzano-weierstrass-real}, $x_{n_k} \to c \in [a, b]$, and then $y_{n_k} \to c$. By continuity $f(x_{n_k}) - f(y_{n_k}) \to f(c) - f(c) = 0$, a contradiction. (4) If
$(x_n)$ is Cauchy and $\delta$ corresponds to $\varepsilon$, eventually $|x_n - x_m| < \delta$, so $|f(x_n) - f(x_m)| < \varepsilon$. For the extension to $a$: for any sequence $x_n \to a$ in $(a, b)$, $(f(x_n))$ is Cauchy,
hence converges; by (2), two such sequences give the same limit (as $x_n - y_n \to 0$), say $L$. By Theorem~\ref{real-analysis:theorem-sequential-limits}, $\lim_{x \to a^+}f(x) = L$, and setting $f(a) = L$ gives a function
continuous at $a$; similarly at $b$.
\end{proof}

\begin{counterexample}\label{real-analysis:counterexample-uniform-continuity}
\begin{enumerate}
\item $f(x) = x^2$ on $\RR$ is continuous but not uniformly continuous: with $x_n = n + 1$ and $y_n = n + 1 + 1/(n+1)$, $x_n - y_n \to 0$ but $f(y_n) - f(x_n) \geq 2$.
\item $f(x) = 1/x$ on $(0, 1)$ is continuous but not uniformly continuous: with $x_n = 1/(n + 2)$ and $y_n = 1/(n + 3)$, $f(y_n) - f(x_n) = 1$; it has no continuous extension to $[0, 1)$, in accordance with
Proposition~\ref{real-analysis:proposition-uniform-continuity-real}(4).
\item $f(x) = \sqrt x$ on $[0, +\infty)$ is uniformly continuous, since $|\sqrt x - \sqrt y| \leq \sqrt{|x - y|}$, but not Lipschitz, since $\sqrt x/x \to +\infty$ as $x \to 0^+$. So the implications in (1) are strict.
\end{enumerate}
\end{counterexample}

\section{Upper and lower limits, and series}\label{real-analysis:section-series}

\begin{definition}\label{real-analysis:definition-limsup}
For a real sequence $(a_n)$ let $\limsup_n a_n = \inf_n\sup_{k \geq n}a_k$ and $\liminf_n a_n = \sup_n\inf_{k \geq n}a_k$, with values in the extended real line.
\end{definition}

\begin{proposition}\label{real-analysis:proposition-monotone}
\begin{enumerate}
\item A bounded monotone real sequence converges, to its supremum or infimum.
\item $(a_n)$ converges to $a \in \RR$ if and only if $\limsup a_n = \liminf a_n = a$.
\item (Bolzano--Weierstrass) Every bounded sequence in $\RR^n$ has a convergent subsequence.
\end{enumerate}
\end{proposition}

\begin{proof}
(1) is Proposition~\ref{real-analysis:proposition-monotone-regular}. (2) The sequences $\sup_{k \geq n}a_k$ and $\inf_{k \geq n}a_k$ are monotone and bracket $a_n$. If both converge to $a$, so does
$a_n$; conversely, if $|a_k - a| < \varepsilon$ for $k \geq N$, both lie within $\varepsilon$ of $a$ for $n \geq N$. (3) A bounded set lies in a closed cube, which is compact
(Theorem~\ref{general-topology:theorem-heine-borel}) and hence sequentially compact (Theorem~\ref{general-topology:theorem-compact-metric}).
\end{proof}

\begin{definition}\label{real-analysis:definition-series}
Let $V$ be a normed space. A series $\sum_n v_n$ \emph{converges} to $s$ if the partial sums $s_N = \sum_{n \leq N}v_n$ converge to $s$, and it \emph{converges absolutely}
if $\sum_n\|v_n\| < \infty$.
\end{definition}

\begin{proposition}\label{real-analysis:proposition-series}
Let $V$ be a Banach space.
\begin{enumerate}
\item An absolutely convergent series converges, and every rearrangement $\sum_n v_{\sigma(n)}$ ($\sigma$ a bijection of $\NN$) converges to the same sum.
\item (Comparison) If $\|v_n\| \leq c_n$ and $\sum c_n < \infty$, then $\sum v_n$ converges absolutely.
\item (Root and ratio tests) If $\limsup\|v_n\|^{1/n} < 1$, or $v_n \neq 0$ and $\limsup\|v_{n+1}\|/\|v_n\| < 1$, then $\sum v_n$ converges absolutely. If $\limsup\|v_n\|^{1/n} > 1$, it
diverges.
\item (Geometric series) For $z \in \CC$ with $|z| < 1$, $\sum_n z^n = 1/(1 - z)$.
\item (Weierstrass $M$-test) If $f_n \colon X \to V$ satisfy $\sup_x\|f_n(x)\| \leq M_n$ with $\sum M_n < \infty$, then $\sum f_n$ converges uniformly.
\end{enumerate}
\end{proposition}

\begin{proof}
(1) $\|s_M - s_N\| \leq \sum_{N < n \leq M}\|v_n\|$, so the partial sums are Cauchy. For a rearrangement, given $\varepsilon$ choose $N$ with $\sum_{n \geq N}\|v_n\| < \varepsilon$ and then $M$
with $\{0, \ldots, N-1\} \subseteq \sigma(\{0, \ldots, M\})$; for $m \geq M$, $\|\sum_{k \leq m}v_{\sigma(k)} - s_{N-1}\| \leq \sum_{n \geq N}\|v_n\| < \varepsilon$, and $\|s_{N-1} - s\| \leq \varepsilon$.
(2) follows from (1). (3) If $\limsup\|v_n\|^{1/n} < q < 1$, then $\|v_n\| \leq q^n$ for large $n$; compare with $\sum q^n$ using (4). If the ratio limsup is $< q < 1$, then
$\|v_n\| \leq Cq^n$ for large $n$. If $\limsup\|v_n\|^{1/n} > 1$, then $\|v_n\| > 1$ infinitely often, so $v_n \not\to 0$ and the partial sums are not Cauchy. (4)
$(1 - z)\sum_{n \leq N}z^n = 1 - z^{N+1}$, and $|z|^N \to 0$: if $0 < |z| < 1$ write $1/|z| = 1 + h$, so $|z|^{-N} \geq 1 + Nh$. (5) The partial sums are bounded and uniformly Cauchy, and the space
$B(X, V)$ of bounded maps is complete (Proposition~\ref{general-topology:proposition-complete-examples}).
\end{proof}

\begin{proposition}[More convergence tests]\label{real-analysis:proposition-series-tests}
Let $(a_n)$, $(b_n)$ be sequences of real numbers.
\begin{enumerate}
\item If $\sum a_n$ converges, then $a_n \to 0$. The converse fails: the harmonic series $\sum_{n \geq 1}1/n$ diverges.
\item (Limit comparison) If $a_n, b_n > 0$ and $a_n/b_n \to L \in (0, \infty)$, then $\sum a_n$ and $\sum b_n$ both converge or both diverge. In particular this holds if $a_n \sim b_n$, that is,
$a_n/b_n \to 1$.
\item (Cauchy condensation) If $a_1 \geq a_2 \geq \cdots \geq 0$, then $\sum_{n \geq 1}a_n$ converges if and only if $\sum_{k \geq 0}2^ka_{2^k}$ converges. Consequently $\sum_{n \geq 1}n^{-p}$
converges if and only if $p > 1$, and $\sum_{n \geq 2}1/(n(\log n)^p)$ converges if and only if $p > 1$.
\item (Ratio and root) If $a_n > 0$, then $\liminf\frac{a_{n+1}}{a_n} \leq \liminf a_n^{1/n} \leq \limsup a_n^{1/n} \leq \limsup\frac{a_{n+1}}{a_n}$; so whenever the ratio test applies, the root test
does too. If $a_{n+1} \geq a_n > 0$ for large $n$, the series diverges. If the limit of the ratios or of the roots is $1$, no conclusion is possible: $\sum 1/n$ diverges and $\sum 1/n^2$
converges, and both have limit $1$.
\item (Dirichlet) If $b_n \in \CC$ have bounded partial sums $B_n = \sum_{k \leq n}b_k$ and $a_n \geq a_{n+1} \to 0$, then $\sum a_nb_n$ converges. For $b_n = (-1)^n$ this is the alternating
series test (Lemma~\ref{real-analysis:lemma-alternating} below).
\end{enumerate}
\end{proposition}

\begin{proof}
(1) $a_n = s_n - s_{n-1} \to s - s = 0$. For the harmonic series, $s_{2n} - s_n = \sum_{n < k \leq 2n}\frac1k \geq n \cdot \frac1{2n} = \frac12$, so the partial sums are not Cauchy.

(2) A series of positive terms converges if and only if its partial sums are bounded (Proposition~\ref{real-analysis:proposition-monotone}). For large $n$, $\frac L2b_n \leq a_n \leq 2Lb_n$, so
the partial sums are bounded for both series or for neither.

(3) For $2^k \leq n < 2^{k+1}$ there are $2^k$ terms, each between $a_{2^{k+1}}$ and $a_{2^k}$, so
$\frac12 2^{k+1}a_{2^{k+1}} \leq \sum_{n = 2^k}^{2^{k+1} - 1}a_n \leq 2^ka_{2^k}$. Summing over $k \leq K$ shows that the partial sums of the two series bound each other, up to a factor
$2$ and the term $a_1$. For $a_n = n^{-p}$ with $p > 0$, $2^ka_{2^k} = (2^{1-p})^k$, a geometric series converging if and only if $p > 1$; for $p \leq 0$ the terms do not tend to $0$. For
$a_n = 1/(n(\log n)^p)$, $2^ka_{2^k} = (k\log 2)^{-p}$, which reduces to the previous case.

(4) If $\limsup a_{n+1}/a_n < q$, then $a_{n+1} \leq qa_n$ for $n \geq N$, so $a_n \leq Cq^n$ and $a_n^{1/n} \leq C^{1/n}q \to q$; hence $\limsup a_n^{1/n} \leq q$. The inequality for the $\liminf$
is proved in the same way. If $a_n$ is eventually nondecreasing and positive, it does not tend to $0$.

(5) Summation by parts: for $m \leq N$, $\sum_{n=m}^Na_nb_n = a_NB_N - a_mB_{m-1} + \sum_{n=m}^{N-1}(a_n - a_{n+1})B_n$. With $|B_n| \leq M$ and $a_n - a_{n+1} \geq 0$, the absolute value is
at most $M(a_N + a_m + (a_m - a_N)) = 2Ma_m$, which tends to $0$; so the partial sums are Cauchy.
\end{proof}

\begin{example}\label{real-analysis:example-series-tests}
\emph{Combining tests.} The tests are sufficient conditions, and in practice one combines them: first replace $a_n$ by a simpler sequence with the same behaviour (limit comparison with an
asymptotic equivalent), then apply the ratio, root or condensation test to it.
\begin{enumerate}
\item $\sum n!/n^n$: the ratio is $(n/(n + 1))^n = (1 + 1/n)^{-n} \to 1/e < 1$, so the series converges.
\item $\sum\sin(1/n)$ diverges and $\sum(1 - \cos(1/n))$ converges, by limit comparison with $\sum 1/n$ and $\sum 1/(2n^2)$, since $\sin t/t \to 1$ and $(1 - \cos t)/t^2 \to \frac12$ as $t \to 0$.
\item $\sqrt{n + 1} - \sqrt n = 1/(\sqrt{n + 1} + \sqrt n) \sim 1/(2\sqrt n)$; so $\sum(\sqrt{n + 1} - \sqrt n)$ diverges (its partial sums are $\sqrt{N + 1} - 1$) and
$\sum(\sqrt{n + 1} - \sqrt n)/n$ converges, by comparison with $\sum n^{-3/2}$.
\item $\sum_{n \geq 1}\sin(n)/n$ converges by Dirichlet's test: from $2\sin\frac12\sin k = \cos(k - \frac12) - \cos(k + \frac12)$ the partial sums of $\sum\sin k$ telescope and are bounded by
$1/\sin\frac12$. It does not converge absolutely, since $|\sin n| \geq \sin^2n = \frac12(1 - \cos 2n)$ and $\sum\cos(2n)/n$ converges by the same argument while $\sum 1/n$ diverges.
\item $\sum n^2/2^n$ converges by the root test ($n^{2/n} \to 1$), and $\sum 2^n/n^2$ diverges since its terms tend to $\infty$.
\end{enumerate}
\end{example}

\begin{lemma}[Alternating series]\label{real-analysis:lemma-alternating}
If $t_0 \geq t_1 \geq \cdots \geq 0$ and $t_k \to 0$, then $S = \sum_k(-1)^kt_k$ converges and $S_{2m+1} \leq S \leq S_{2m}$ for all $m$, where $S_N$ are the partial sums.
\end{lemma}

\begin{proof}
$S_{2m+2} - S_{2m} = -t_{2m+1} + t_{2m+2} \leq 0$ and $S_{2m+3} - S_{2m+1} = t_{2m+2} - t_{2m+3} \geq 0$, and $S_{2m+1} = S_{2m} - t_{2m+1} \leq S_{2m}$. So the even partial sums decrease, the odd ones
increase, every odd one is below every even one, and their difference tends to $0$.
\end{proof}

\emph{Rearranging series.} A finite sum does not depend on the order of its terms, and by Proposition~\ref{real-analysis:proposition-series}(1) neither does the sum of an absolutely convergent
series. For series that converge but not absolutely the order matters, and Riemann discovered that it matters as much as it possibly could: by reordering the terms one can make the sum anything
one likes. A series $\sum_na_n$ of real numbers is \emph{conditionally convergent} if it converges but $\sum_n|a_n| = \infty$; the alternating harmonic series $\sum_{n \geq 1}(-1)^{n+1}/n$ is the basic
example (Lemma~\ref{real-analysis:lemma-alternating} and Proposition~\ref{real-analysis:proposition-series-tests}(1)). Write $a^+ = \max(a, 0)$ and $a^- = \max(-a, 0)$ for the positive and negative parts
of a real number, so that $a = a^+ - a^-$ and $|a| = a^+ + a^-$.

\begin{lemma}\label{real-analysis:lemma-conditional-parts}
If $\sum_na_n$ is a conditionally convergent series of real numbers, then $a_n \to 0$ and $\sum_na_n^+ = \sum_na_n^- = \infty$. In particular infinitely many $a_n$ are positive and infinitely
many are negative.
\end{lemma}

\begin{proof}
$a_n \to 0$ by Proposition~\ref{real-analysis:proposition-series-tests}(1). If $\sum_na_n^+$ were finite, then $\sum_na_n^- = \sum_n(a_n^+ - a_n)$ would converge, being the difference of two
convergent series, and so would $\sum_n|a_n| = \sum_n(a_n^+ + a_n^-)$, a contradiction. The same argument applies to $\sum_na_n^-$. A series of non-negative terms with only finitely many nonzero
terms is finite, so infinitely many $a_n^+$ and infinitely many $a_n^-$ are nonzero.
\end{proof}

So the positive terms alone have infinite sum, and so do the negative terms alone; the series converges only because they are interleaved so as to cancel. This is what a rearrangement can
exploit.

\begin{theorem}[Riemann's rearrangement theorem]\label{real-analysis:theorem-riemann-rearrangement}
Let $\sum_na_n$ be a conditionally convergent series of real numbers, and let $-\infty \leq \alpha \leq \beta \leq \infty$. There is a bijection $\sigma \colon \NN \to \NN$ such that the partial sums
$T_N = \sum_{n \leq N}a_{\sigma(n)}$ satisfy
\[
\liminf_{N \to \infty}T_N = \alpha \qquad\text{and}\qquad \limsup_{N \to \infty}T_N = \beta .
\]
In particular, for every $c \in \RR$ some rearrangement converges to $c$, and some rearrangements diverge to $\infty$ and to $-\infty$.
\end{theorem}

\begin{proof}
\emph{Targets.} Choose real numbers $\alpha_k \to \alpha$ and $\beta_k \to \beta$ with $\alpha_k < \beta_k$ and $\alpha_k < \beta_{k+1}$ for all $k \geq k_0$: put $\beta_k = \beta + \frac1k$, $k$ or $-k$
according as $\beta$ is finite, $+\infty$ or $-\infty$, and $\alpha_k = \alpha - \frac1k$, $k - 1$ or $-k - 2$ according as $\alpha$ is finite, $+\infty$ or $-\infty$. One checks the inequalities case by
case for $k$ large (recall $\alpha \leq \beta$; for instance if $\alpha = \beta = -\infty$, then $\alpha_k = -k - 2 < -k - 1 = \beta_{k+1}$).

\emph{Construction.} Let $p_1 < p_2 < \cdots$ be the indices $n$ with $a_n \geq 0$ and $q_1 < q_2 < \cdots$ those with $a_n < 0$; both lists are infinite by
Lemma~\ref{real-analysis:lemma-conditional-parts}, and together they contain every index once. We build the rearrangement in stages $k = k_0, k_0 + 1, \ldots$, each consisting of an \emph{up step} and
a \emph{down step}. In the up step of stage $k$ we append the next unused terms $a_{p_j}$, in order, as long as the running sum is $\leq \beta_k$, stopping as soon as it exceeds $\beta_k$. In the down step
we append the next unused terms $a_{q_j}$, in order, until the running sum is $< \alpha_k$. Each step stops after finitely many terms: the unused positive terms are all but finitely many of the $a_n^+$,
so their sum is still infinite by Lemma~\ref{real-analysis:lemma-conditional-parts}, and likewise for the negative terms.

Every index is used. A down step adds at least one term, because it starts with a sum $> \beta_k > \alpha_k$; an up step of stage $k + 1$ adds at least one term, because it starts with a sum
$< \alpha_k < \beta_{k+1}$. So infinitely many $p_j$ and infinitely many $q_j$ are used, and since each list is used in order, all of them are. Hence the order in which the terms are appended defines a
bijection $\sigma \colon \NN \to \NN$.

\emph{The partial sums.} Let $U_k$ be the running sum at the end of the up step of stage $k$ and $L_k$ the one at the end of the down step. If the up step added terms, then just before its last term
$a_{u_k}$ the sum was $\leq \beta_k$, so
\[
\beta_k < U_k \leq \beta_k + a_{u_k}, \qquad\text{and similarly}\qquad \alpha_k + a_{l_k} \leq L_k < \alpha_k,
\]
with $a_{l_k} < 0$ the last term of the down step (for $k > k_0$ the up step always adds terms, as shown above). The indices $u_k$ and $l_k$ are distinct for distinct $k$, so they tend to $\infty$, and
$a_{u_k}, a_{l_k} \to 0$ because $a_n \to 0$. Hence $U_k - \beta_k \to 0$ and $L_k - \alpha_k \to 0$, so $U_k \to \beta$ and $L_k \to \alpha$ (in the extended sense when $\alpha$ or $\beta$ is
infinite). Within an up step the running sum increases, and within a down step it decreases; so for $N$ in the up step of stage $k$ we have $L_{k-1} \leq T_N \leq U_k$, and for $N$ in the down step
$L_k \leq T_N \leq U_k$. Therefore $\limsup_NT_N \leq \limsup_kU_k = \beta$ and $\liminf_NT_N \geq \liminf_kL_k = \alpha$; and the values $T_N = U_k$ and $T_N = L_k$ occur for every $k$, so
equality holds in both. The last statement is the case $\alpha = \beta$.
\end{proof}

\begin{corollary}\label{real-analysis:corollary-unconditional}
A series $\sum_na_n$ of real numbers is absolutely convergent if and only if every rearrangement of it converges, and if and only if all its rearrangements converge to the same sum.
\end{corollary}

\begin{proof}
If $\sum_n|a_n| < \infty$, every rearrangement converges to the same sum by Proposition~\ref{real-analysis:proposition-series}(1). If $\sum_n|a_n| = \infty$, either the series diverges, and then the
identity rearrangement does not converge, or it is conditionally convergent, and then Theorem~\ref{real-analysis:theorem-riemann-rearrangement} with $\alpha = \beta = \infty$ gives a rearrangement that
diverges.
\end{proof}

\begin{example}\label{real-analysis:example-rearrangement}
\begin{enumerate}
\item Let $S = 1 - \frac12 + \frac13 - \frac14 + \cdots$. By Lemma~\ref{real-analysis:lemma-alternating}, $\frac12 = S_1 \leq S \leq S_0 = 1$, so $S \neq 0$. Take one positive term followed by two
negative ones:
\[
1 - \frac12 - \frac14 + \frac13 - \frac16 - \frac18 + \frac15 - \frac1{10} - \frac1{12} + \cdots
\]
Every term of $S$ occurs exactly once: the odd denominators in order, and the even denominators in order, two at a time. The $k$-th group of three terms is
$\frac{1}{2k - 1} - \frac{1}{4k - 2} - \frac{1}{4k} = \frac{1}{4k - 2} - \frac{1}{4k} = \frac12\bigl(\frac{1}{2k - 1} - \frac{1}{2k}\bigr)$. So the partial sums after $3m$ terms are half the partial
sums of $S$ after $2m$ terms, and tend to $\frac12S$; the partial sums in between differ from these by at most two terms, which tend to $0$. So this rearrangement converges to $\frac12S \neq S$.
\item The proof of the theorem is an algorithm. To rearrange $S$ so that it converges to $c = 2$: take $1 + \frac13 + \frac15 + \cdots$ until the sum exceeds $2$, which happens at
$1 + \frac13 + \cdots + \frac1{15} \approx 2.02$ (eight terms), then subtract $\frac12$ to get $\approx 1.52 < 2$, then add $\frac1{17} + \frac1{19} + \cdots$ until the sum exceeds $2$ again, and so on.
The positive terms are used much faster than the negative ones. This is how it must be: a rearrangement with a larger sum takes more positive terms per negative term, just as in (1)
taking fewer positive terms per negative term gave a smaller sum.
\item For complex series the theorem fails as stated: the rearrangements of $\sum_ni(-1)^{n+1}/n$ have as sums exactly the points of the imaginary axis $i\RR$, by the theorem applied to the real series
$\sum_n(-1)^{n+1}/n$. For series in $\RR^m$ the possible sums of rearrangements form a more complicated set, which can be a line in the plane even when the series does not lie on a line, as the example
$\sum_n\bigl((-1)^{n+1}/n, (-1)^{n+1}/n + 1/n^2\bigr)$ shows: the second coordinate differs from the first by an absolutely convergent series, whose sum does not change, so every sum lies on the line
$y = x + \sum_n1/n^2$, and every point of it occurs.
\end{enumerate}
\end{example}

\begin{proposition}[Cauchy product]\label{real-analysis:proposition-cauchy-product}
If $\sum a_n$ and $\sum b_n$ are absolutely convergent series in $\CC$, then $c_n = \sum_{k=0}^na_kb_{n-k}$ defines an absolutely convergent series with
$\sum c_n = (\sum a_n)(\sum b_n)$.
\end{proposition}

\begin{proof}
Let $A = \sum|a_n|$ and $B = \sum|b_n|$. Then $\sum_{n \leq N}|c_n| \leq AB$, so $\sum c_n$ converges absolutely. The difference $\sum_{n \leq 2N}c_n - (\sum_{k \leq N}a_k)(\sum_{l \leq N}b_l)$ is the sum
of $a_kb_l$ over $k + l \leq 2N$ with $k > N$ or $l > N$, bounded by $B\sum_{k > N}|a_k| + A\sum_{l > N}|b_l| \to 0$.
\end{proof}

\begin{theorem}\label{real-analysis:theorem-power-series}
Let $a_n \in \CC$ and $R = 1/\limsup_n|a_n|^{1/n} \in [0, \infty]$, the \emph{radius of convergence}. The power series $\sum_n a_nz^n$ converges absolutely for $|z| < R$, uniformly on
$|z| \leq r$ for every $r < R$, and diverges for $|z| > R$. The series $\sum_n na_nz^{n-1}$ and $\sum_n n^2a_nz^n$ have the same radius of convergence.
\end{theorem}

\begin{proof}
For $|z| \leq r < R$ choose $r < q < R$; then $|a_n| \leq q^{-n}$ for large $n$, so $|a_nz^n| \leq (r/q)^n$, and the $M$-test applies. For $|z| > R$, $\limsup|a_nz^n|^{1/n} > 1$. Finally
$n^{1/n} \to 1$: writing $n^{1/n} = 1 + h_n$, the binomial formula gives $n \geq \binom{n}{2}h_n^2$, so $h_n^2 \leq 2/(n-1)$. Hence $\limsup|n^ka_n|^{1/n} = \limsup|a_n|^{1/n}$.
\end{proof}

\section{Differentiation in one variable}\label{real-analysis:section-one-variable}

\begin{definition}\label{real-analysis:definition-derivative-one}
Let $I \subseteq \RR$ be an interval with more than one point and $f \colon I \to \mathbb{K}^m$. Then $f$ is \emph{differentiable} at $a \in I$ with derivative $f'(a)$ if
$(f(x) - f(a))/(x - a) \to f'(a)$ as $x \to a$ in $I \setminus \{a\}$. Equivalently, there is a function $\phi \colon I \to \mathbb{K}^m$, continuous at $a$, with $f(x) - f(a) = \phi(x)(x - a)$;
then $f'(a) = \phi(a)$. The same definition with $I$ replaced by an open subset of $\CC$ and $x$ complex defines the \emph{complex derivative}.
\end{definition}

\begin{proposition}\label{real-analysis:proposition-derivative-rules}
Differentiable functions are continuous. Sums, products and quotients (where the denominator does not vanish) of differentiable scalar functions are differentiable with the usual
formulas, and if $f$ is differentiable at $a$ and $g$ at $f(a)$, then $(g \circ f)'(a) = g'(f(a))f'(a)$.
\end{proposition}

\begin{proof}
Write $f(x) - f(a) = \phi(x)(x - a)$ and $g(y) - g(b) = \psi(y)(y - b)$ with $\phi$, $\psi$ continuous at $a$, $b = f(a)$. Continuity is clear. For products,
$fg(x) - fg(a) = (\phi(x)g(x) + f(a)\chi(x))(x - a)$ if $g(x) - g(a) = \chi(x)(x - a)$. For $1/f$, $1/f(x) - 1/f(a) = -\phi(x)(x - a)/(f(x)f(a))$. For composites,
$g(f(x)) - g(f(a)) = \psi(f(x))\phi(x)(x - a)$ and $\psi \circ f$ is continuous at $a$. Evaluate the continuous factors at $a$.
\end{proof}

\begin{definition}\label{real-analysis:definition-local-extremum}
Let $f \colon I \to \RR$ on an interval $I$. Then $f$ has a \emph{local maximum} at $c \in I$ if there is $\delta > 0$ with $f(x) \leq f(c)$ for all $x \in I$ with $|x - c| < \delta$, a \emph{strict} local maximum if
$f(x) < f(c)$ for these $x \neq c$, and similarly for (strict) local minima. A \emph{local extremum} is a local maximum or minimum.
\end{definition}

\begin{lemma}[Fermat]\label{real-analysis:lemma-fermat}
If $f \colon I \to \RR$ has a local extremum at an interior point $c$ of $I$ and is differentiable at $c$, then $f'(c) = 0$.
\end{lemma}

\begin{proof}
Say $c$ is a local maximum. For $x > c$ close to $c$, $(f(x) - f(c))/(x - c) \leq 0$, so $f'(c) \leq 0$ by Corollary~\ref{real-analysis:corollary-limit-rules}(2); for $x < c$ the quotient is $\geq 0$, so $f'(c) \geq 0$.
\end{proof}

\begin{counterexample}\label{real-analysis:counterexample-fermat}
The converse fails: $f(x) = x^3$ has $f'(0) = 0$ and no extremum at $0$. The hypothesis that $c$ is interior is needed: $f(x) = x$ on $[0, 1]$ has a maximum at $1$ with $f'(1) = 1$. And $|x|$ has a minimum at $0$,
where it is not differentiable.
\end{counterexample}

\begin{theorem}[Mean value theorem]\label{real-analysis:theorem-mean-value}
Let $f, g \colon [a, b] \to \RR$ be continuous and differentiable on $(a, b)$.
\begin{enumerate}
\item (Rolle) If $f(a) = f(b)$, then $f'(\xi) = 0$ for some $\xi \in (a, b)$.
\item There is $\xi \in (a, b)$ with $f(b) - f(a) = f'(\xi)(b - a)$, and (Cauchy) some $\eta$ with $(f(b) - f(a))g'(\eta) = (g(b) - g(a))f'(\eta)$.
\item If $f' = 0$ on $(a, b)$, $f$ is constant; if $f' \geq 0$ ($> 0$), $f$ is increasing (strictly increasing).
\end{enumerate}
\end{theorem}

\begin{proof}
(1) By Corollary~\ref{general-topology:corollary-extreme-value}, $f$ attains a maximum and a minimum; if both are at endpoints, $f$ is constant. Otherwise one is at some
$\xi \in (a, b)$, and $f'(\xi) = 0$ by Lemma~\ref{real-analysis:lemma-fermat}. (2) Apply (1) to $f(x) - \frac{f(b) - f(a)}{b - a}(x - a)$, and to
$(f(b) - f(a))g(x) - (g(b) - g(a))f(x)$. (3) Apply (2) on subintervals.
\end{proof}

\begin{remark}\label{real-analysis:remark-derivative-zero}
Both hypotheses of (3) matter. On a set that is not an interval a function with zero derivative need not be constant: $f = 0$ on $(0, 1)$ and $f = 1$ on $(2, 3)$. And ``$f^{\prime} = 0$'' must hold at
every point: the Cantor function (Chapter~\ref{measure-theory}) is continuous, nondecreasing and nonconstant on $[0, 1]$, with $f^{\prime} = 0$ outside a set of measure zero. On the other hand finitely
many exceptional points do not matter: if $f$ is continuous on $[a, b]$ and $f^{\prime} = 0$ except at $c_1 < \cdots < c_k$, applying (3) on each interval between consecutive points and using continuity
shows that $f$ is constant.
\end{remark}

\begin{proposition}\label{real-analysis:proposition-inverse-one}
Let $f \colon I \to \RR$ be continuous and strictly monotone. Then $J = f(I)$ is an interval, $f^{-1} \colon J \to I$ is continuous, and if $f$ is differentiable at $a$ with
$f'(a) \neq 0$, then $f^{-1}$ is differentiable at $f(a)$ with derivative $1/f'(a)$.
\end{proposition}

\begin{proof}
$J$ is an interval by Corollary~\ref{general-topology:corollary-ivt}. A strictly monotone bijection of intervals maps subintervals onto subintervals, so $f^{-1}$ maps open subintervals of $J$
(relative to $J$) to open subintervals of $I$, and is continuous. Write $f(x) - f(a) = \phi(x)(x - a)$; strict monotonicity gives $\phi(x) \neq 0$ for $x \neq a$, and $\phi(a) = f'(a) \neq 0$. So
$f^{-1}(y) - a = (y - f(a))/\phi(f^{-1}(y))$ with $1/(\phi \circ f^{-1})$ continuous at $f(a)$.
\end{proof}

\begin{proposition}[Limits of derivatives]\label{real-analysis:proposition-derivative-limit}
Let $f \colon [a, b) \to \RR$ be continuous, differentiable on $(a, b)$, and suppose $\lim_{x \to a^+}f'(x) = L \in \RR$. Then $f$ is differentiable at $a$ (from the right) with $f'(a) = L$. Consequently, if $f$ is
differentiable on an open interval containing $c$ and $f'$ has one-sided limits at $c$, both equal $f'(c)$: a derivative has neither removable nor jump discontinuities.
\end{proposition}

\begin{proof}
For $a < x < b$, the mean value theorem gives $\xi_x \in (a, x)$ with $(f(x) - f(a))/(x - a) = f'(\xi_x)$. If $x_n \to a^+$, then $\xi_{x_n} \to a^+$ by the squeeze theorem, so $f'(\xi_{x_n}) \to L$; by
Theorem~\ref{real-analysis:theorem-sequential-limits}, the difference quotient tends to $L$. For the second statement, apply this on each side of $c$.
\end{proof}

\begin{theorem}[de l'H\^opital]\label{real-analysis:theorem-lhopital}
Let $f, g \colon (a, b) \to \RR$ be differentiable, with $-\infty \leq a < b \leq +\infty$, and $g'(x) \neq 0$ for all $x$. Suppose that $\lim_{x \to a^+}f'(x)/g'(x) = L \in \bar{\RR}$ and that either
\begin{enumerate}
\item $\lim_{x \to a^+}f(x) = \lim_{x \to a^+}g(x) = 0$, or
\item $\lim_{x \to a^+}|g(x)| = +\infty$ and $L \in \RR$.
\end{enumerate}
Then $\lim_{x \to a^+}f(x)/g(x) = L$. The same holds for left limits at $b$, and hence for two-sided limits.
\end{theorem}

\begin{proof}
Since $g' \neq 0$, $g$ is injective on $(a, b)$ by Rolle's theorem, so $g(x) \neq g(y)$ for $x \neq y$, and by the Cauchy mean value theorem (Theorem~\ref{real-analysis:theorem-mean-value}(2)), for $x < y$ in $(a, b)$ there is
$\xi \in (x, y)$ with
\begin{equation}\label{real-analysis:equation-cauchy-quotient}
\frac{f(x) - f(y)}{g(x) - g(y)} = \frac{f^{\prime}(\xi)}{g^{\prime}(\xi)}.
\end{equation}
Case (1) with $a \in \RR$: extend $f$ and $g$ by $f(a) = g(a) = 0$; they become continuous on $[a, b)$, and \eqref{real-analysis:equation-cauchy-quotient} holds for $x = a$ as well, with $g(y) \neq 0$ for $y > a$ by Rolle's theorem. Let $y_n \to a^+$ and choose
$\xi_n \in (a, y_n)$ with $f(y_n)/g(y_n) = f'(\xi_n)/g'(\xi_n)$. Then $\xi_n \to a^+$, so $f(y_n)/g(y_n) \to L$, and Theorem~\ref{real-analysis:theorem-sequential-limits} gives the claim. If $a = -\infty$, apply the case just
proved to $F(t) = f(-1/t)$ and $G(t) = g(-1/t)$ on an interval $(0, \eta)$, noting $F'(t)/G'(t) = f'(-1/t)/g'(-1/t)$ and using Proposition~\ref{real-analysis:proposition-composite-limit}.

Case (2). Let $\varepsilon > 0$ and choose $c \in (a, b)$ with $|f'(\xi)/g'(\xi) - L| < \varepsilon$ for $a < \xi < c$. Fix $y \in (a, c)$. For $a < x < y$, \eqref{real-analysis:equation-cauchy-quotient} gives $Q(x) = (f(x) - f(y))/(g(x) - g(y))$ with $|Q(x) - L| < \varepsilon$. Since
$|g(x)| \to +\infty$, $g(x) \neq 0$ for $x$ close to $a$, and
\[
\frac{f(x)}{g(x)} - L = \big(Q(x) - L\big)\Big(1 - \frac{g(y)}{g(x)}\Big) - L\frac{g(y)}{g(x)} + \frac{f(y)}{g(x)}.
\]
As $x \to a^+$ with $y$ fixed, the last two terms tend to $0$ and $1 - g(y)/g(x) \to 1$, so $|f(x)/g(x) - L| < 2\varepsilon$ for $x$ close to $a$. Left limits at $b$ follow by applying the result to $x \mapsto f(-x)$, $g(-x)$.
\end{proof}

\begin{theorem}[Darboux]\label{real-analysis:theorem-darboux}
Let $f \colon [a, b] \to \RR$ be differentiable, with one-sided derivatives at the endpoints. Then $f'$ takes every value between $f'(a)$ and $f'(b)$.
\end{theorem}

\begin{proof}
Let $f'(a) < y < f'(b)$ (the other case follows by passing to $-f$) and $g(x) = f(x) - yx$. Then $g'(a) < 0 < g'(b)$. By Theorem~\ref{real-analysis:theorem-weierstrass-extreme}, $g$ attains a minimum at some
$c \in [a, b]$. As $g'(a) < 0$, $(g(x) - g(a))/(x - a) < 0$ for $x > a$ close to $a$ (Proposition~\ref{real-analysis:proposition-function-limit-basic}(2)), so $g(x) < g(a)$ and $c \neq a$; similarly $c \neq b$. By
Lemma~\ref{real-analysis:lemma-fermat}, $g'(c) = 0$, that is, $f'(c) = y$.
\end{proof}

\begin{theorem}[Taylor]\label{real-analysis:theorem-taylor-one}
Let $f \colon [a, x] \to \RR$ (or $[x, a]$) have $n$ continuous derivatives, with $f^{(n+1)}$ existing on the open interval. Then there is $\xi$ strictly between $a$ and $x$ with
\[
f(x) = \sum_{k=0}^n\frac{f^{(k)}(a)}{k!}(x - a)^k + \frac{f^{(n+1)}(\xi)}{(n+1)!}(x - a)^{n+1}.
\]
\end{theorem}

\begin{proof}
Let $g(t) = f(x) - \sum_{k \leq n}f^{(k)}(t)(x - t)^k/k! - M(x - t)^{n+1}/(n+1)!$, with $M$ chosen so that $g(a) = 0$. Also $g(x) = 0$. In the derivative the sum telescopes:
$g'(t) = -f^{(n+1)}(t)(x - t)^n/n! + M(x - t)^n/n!$. By Rolle's theorem $g'(\xi) = 0$ for some $\xi$ between $a$ and $x$, so $M = f^{(n+1)}(\xi)$.
\end{proof}

\begin{proposition}[Taylor's formula with Peano remainder]\label{real-analysis:proposition-taylor-peano}
Let $n \geq 1$, let $f \colon I \to \RR$ be $n - 1$ times differentiable on an interval $I$, and suppose $f^{(n)}(a)$ exists at a point $a \in I$. Let $T_n(x) = \sum_{k=0}^nf^{(k)}(a)(x - a)^k/k!$. Then
\[
\lim_{x \to a}\frac{f(x) - T_n(x)}{(x - a)^n} = 0,
\]
that is, $f(x) = T_n(x) + o((x - a)^n)$ as $x \to a$.
\end{proposition}

\begin{proof}
Let $R = f - T_n$. Then $R$ is $n - 1$ times differentiable, $R^{(k)}(a) = 0$ for $0 \leq k \leq n$, and each $R^{(k)}$ with $k \leq n - 1$ is continuous at $a$. For $n = 1$ the claim is
$(R(x) - R(a))/(x - a) \to R'(a) = 0$. For $n \geq 2$, apply Theorem~\ref{real-analysis:theorem-lhopital}(1) $n - 1$ times, on each side of $a$, to the quotients $R^{(j)}(x)/(x - a)^{n-j}$, $0 \leq j \leq n - 2$, whose
numerators and denominators tend to $0$ and whose denominators have nonvanishing derivatives for $x \neq a$. It remains to see that the last quotient tends to $0$:
\[
\frac{R^{(n-1)}(x)}{n!\,(x - a)} = \frac{1}{n!}\cdot\frac{R^{(n-1)}(x) - R^{(n-1)}(a)}{x - a} \longrightarrow \frac{R^{(n)}(a)}{n!} = 0.
\]
\end{proof}

\begin{corollary}[Second derivative test]\label{real-analysis:corollary-second-derivative-test}
Let $f$ be differentiable on an open interval containing $c$, with $f'(c) = 0$, and suppose $f''(c)$ exists. If $f''(c) > 0$, then $f$ has a strict local minimum at $c$; if $f''(c) < 0$, a strict local maximum.
\end{corollary}

\begin{proof}
By Proposition~\ref{real-analysis:proposition-taylor-peano} with $n = 2$, $(f(x) - f(c))/(x - c)^2 \to f''(c)/2 > 0$, so $f(x) - f(c) > 0$ for $x \neq c$ close to $c$ (Proposition~\ref{real-analysis:proposition-function-limit-basic}(2)).
\end{proof}

\begin{counterexample}\label{real-analysis:counterexample-second-derivative}
If $f''(c) = 0$ the test says nothing: $x^4$, $-x^4$ and $x^3$ all have vanishing first and second derivatives at $0$, and have a minimum, a maximum and no extremum there.
\end{counterexample}

\begin{definition}\label{real-analysis:definition-convex}
A function $f \colon I \to \RR$ on an interval is \emph{convex} if $f((1 - t)x + ty) \leq (1 - t)f(x) + tf(y)$ for all $x, y \in I$ and $t \in [0, 1]$, and \emph{strictly convex} if the inequality is strict whenever $x \neq y$ and
$0 < t < 1$. It is \emph{(strictly) concave} if $-f$ is (strictly) convex.
\end{definition}

\begin{lemma}[Three slopes]\label{real-analysis:lemma-three-slopes}
A function $f \colon I \to \RR$ is convex if and only if for all $x < y < z$ in $I$,
\[
\frac{f(y) - f(x)}{y - x} \leq \frac{f(z) - f(x)}{z - x} \leq \frac{f(z) - f(y)}{z - y}.
\]
\end{lemma}

\begin{proof}
Write $y = (1 - t)x + tz$ with $t = (y - x)/(z - x) \in (0, 1)$. Convexity at this triple says $(z - x)f(y) \leq (z - y)f(x) + (y - x)f(z)$. Multiplying out, each of the two displayed inequalities is equivalent to
this one. Since every $t \in (0, 1)$ and $x < z$ arise, and the cases $t \in \{0, 1\}$ or $x = y$ are trivial, the lemma follows (for $x > y$ exchange the roles).
\end{proof}

\begin{proposition}\label{real-analysis:proposition-convex}
Let $f \colon I \to \RR$ on an open interval $I$.
\begin{enumerate}
\item If $f$ is convex, it is continuous, it has one-sided derivatives $f'_-(x) \leq f'_+(x)$ at every point, and $f'_+(x) \leq \frac{f(y) - f(x)}{y - x} \leq f'_-(y)$ for $x < y$.
\item If $f$ is differentiable, the following are equivalent: $f$ is convex; $f'$ is increasing; $f(y) \geq f(x) + f'(x)(y - x)$ for all $x, y \in I$.
\item If $f$ is twice differentiable, $f$ is convex if and only if $f'' \geq 0$; if $f'' > 0$, $f$ is strictly convex.
\item (Jensen) If $f$ is convex, $x_1, \ldots, x_k \in I$, $t_i \geq 0$ and $\sum t_i = 1$, then $f(\sum t_ix_i) \leq \sum t_if(x_i)$.
\end{enumerate}
\end{proposition}

\begin{proof}
(1) By Lemma~\ref{real-analysis:lemma-three-slopes}, the slope $s(x, y) = (f(y) - f(x))/(y - x)$ is increasing in each variable. For fixed $x$, $s(x, y)$ decreases as $y \downarrow x$ and is bounded below by $s(w, x)$ for any
$w < x$ in $I$ (which exists, $I$ being open); so $f'_+(x) = \lim_{y \to x^+}s(x, y)$ exists by Proposition~\ref{real-analysis:proposition-monotone-limits}, and $f'_+(x) \geq s(w, x)$. Letting $w \uparrow x$ gives $f'_-(x)$, with
$f'_-(x) \leq f'_+(x)$. The inequality for $x < y$ follows from the monotonicity of slopes. The existence of one-sided derivatives implies continuity, as $f(y) - f(x) = s(x, y)(y - x) \to 0$.
(2) If $f$ is convex and differentiable, (1) gives $f'(x) = f'_+(x) \leq f'_-(y) = f'(y)$ for $x < y$. If $f'$ is increasing and $y > x$, the mean value theorem gives $f(y) - f(x) = f'(\xi)(y - x) \geq f'(x)(y - x)$; for $y < x$,
$f'(\xi) \leq f'(x)$ and $y - x < 0$ give the same inequality. Finally, if the tangent inequality holds, let $z = (1 - t)x + ty$; adding $(1 - t)$ times $f(x) \geq f(z) + f'(z)(x - z)$ and $t$ times $f(y) \geq f(z) + f'(z)(y - z)$ gives
$(1 - t)f(x) + tf(y) \geq f(z)$.
(3) $f'' \geq 0$ if and only if $f'$ is increasing (Theorem~\ref{real-analysis:theorem-mean-value}(3) and, conversely, $f'' = \lim$ of nonnegative quotients); apply (2). If $f'' > 0$, $f'$ is strictly increasing and the
inequalities in the proof of (2) become strict for $y \neq x$, hence so does the final one for $x \neq y$, $0 < t < 1$.
(4) Induction on $k$: for $t_k < 1$, $\sum t_ix_i = (1 - t_k)\sum_{i<k}\frac{t_i}{1 - t_k}x_i + t_kx_k$, and the inner point lies in $I$ because $I$ is an interval.
\end{proof}

\begin{counterexample}\label{real-analysis:counterexample-convex}
A convex function need not be differentiable: $|x|$ has $f'_-(0) = -1 < 1 = f'_+(0)$. On a closed interval a convex function need not be continuous at the endpoints: the function equal to $0$ on $[0, 1)$ and to $1$
at $1$ is convex. The converse of the last statement of (3) fails: $x^4$ is strictly convex, but its second derivative vanishes at $0$.
\end{counterexample}

\begin{theorem}\label{real-analysis:theorem-power-series-derivative}
Let $f(z) = \sum_na_nz^n$ have radius of convergence $R > 0$. Then $f$ is complex differentiable on $\{|z| < R\}$ with $f'(z) = \sum_nna_nz^{n-1}$. Hence $f$ has complex derivatives of all
orders there, and $a_n = f^{(n)}(0)/n!$. The same holds for real $z$ in $(-R, R)$.
\end{theorem}

\begin{proof}
Let $g(z) = \sum na_nz^{n-1}$, which has radius $R$ by Theorem~\ref{real-analysis:theorem-power-series}. Fix $|z| < \rho < R$ and $h \neq 0$ with $|z| + |h| \leq \rho$. By the binomial formula,
$(z + h)^n - z^n - nz^{n-1}h = \sum_{k \geq 2}\binom{n}{k}z^{n-k}h^k$, and $\binom{n}{k} \leq \binom{n}{2}\binom{n-2}{k-2}$ for $k \geq 2$, so its absolute value is at most
$\binom{n}{2}|h|^2(|z| + |h|)^{n-2} \leq n^2|h|^2\rho^{n-2}$. Hence
$|(f(z + h) - f(z))/h - g(z)| \leq |h|\sum_n n^2|a_n|\rho^{n-2}$, and the last series converges. Iterating gives all derivatives, and evaluating at $0$ gives $a_n$.
\end{proof}

\section{The exponential and trigonometric functions}\label{real-analysis:section-exponential}

\begin{definition}\label{real-analysis:definition-exponential}
For $z \in \CC$ let $\exp(z) = e^z = \sum_{n \geq 0}z^n/n!$ (radius of convergence $\infty$, by the ratio test), and for $x \in \RR$ let $\cos x = \operatorname{Re} e^{ix}$ and
$\sin x = \operatorname{Im} e^{ix}$, so that $\cos x = \sum_k(-1)^kx^{2k}/(2k)!$ and $\sin x = \sum_k(-1)^kx^{2k+1}/(2k+1)!$.
\end{definition}

\begin{proposition}\label{real-analysis:proposition-exponential}
\begin{enumerate}
\item $e^{z+w} = e^ze^w$, $e^z \neq 0$, $\overline{e^z} = e^{\bar z}$, and $\exp' = \exp$.
\item $\exp$ restricts to a strictly increasing bijection $\RR \to (0, \infty)$, whose inverse $\log$ is differentiable with $\log'(x) = 1/x$ and $\log(xy) = \log x + \log y$.
\item $\cos^2 x + \sin^2 x = 1$, $\cos' = -\sin$, $\sin' = \cos$, and $\cos(x + y) = \cos x\cos y - \sin x\sin y$, $\sin(x + y) = \sin x\cos y + \cos x\sin y$.
\end{enumerate}
\end{proposition}

\begin{proof}
(1) By Proposition~\ref{real-analysis:proposition-cauchy-product} and the binomial formula, $e^ze^w = \sum_n\sum_k\frac{z^kw^{n-k}}{k!(n-k)!} = \sum_n(z + w)^n/n!$. So $e^ze^{-z} = 1$. Conjugation is
continuous and commutes with the partial sums. The derivative is Theorem~\ref{real-analysis:theorem-power-series-derivative}. (2) For real $x$, $e^x = (e^{x/2})^2 > 0$, so $\exp$ is strictly
increasing by Theorem~\ref{real-analysis:theorem-mean-value}. For $x \geq 0$, $e^x \geq 1 + x$, and $e^{-x} = 1/e^x$; so $\exp$ is unbounded above and takes values arbitrarily close to $0$, and
by the intermediate value theorem its image is $(0, \infty)$. Apply Proposition~\ref{real-analysis:proposition-inverse-one}. (3) $|e^{ix}|^2 = e^{ix}\overline{e^{ix}} = e^{ix}e^{-ix} = 1$;
$(e^{ix})' = ie^{ix}$ by Theorem~\ref{real-analysis:theorem-power-series-derivative} applied to the real power series $\sum_n i^nx^n/n!$; the addition formulas are the real and imaginary parts of $e^{i(x+y)} = e^{ix}e^{iy}$.
\end{proof}

\begin{theorem}\label{real-analysis:theorem-pi}
The function $\cos$ has a unique zero in $(0, 2)$; call it $\pi/2$. Then $e^{i\pi/2} = i$, the map $t \mapsto e^{it}$ is a surjective continuous homomorphism $\RR \to S^1 = \{|z| = 1\}$ with kernel
$2\pi\ZZ$, and $p(t) = e^{2\pi it}$ is a covering map $\RR \to S^1$. Every $z \neq 0$ is $re^{i\theta}$ with $r = |z|$ and $\theta$ unique modulo $2\pi$.
\end{theorem}

\begin{proof}
For $0 < x \leq 2$ the terms $x^{2k+1}/(2k+1)!$ decrease, since consecutive ratios are $x^2/((2k+2)(2k+3)) < 1$; by Lemma~\ref{real-analysis:lemma-alternating},
$\sin x \geq x - x^3/6 > 0$. So $\cos$ is strictly decreasing on $[0, 2]$. The terms $4^k/(2k)!$ of $\cos 2$ decrease from $k = 1$ on, so $\cos 2 \leq 1 - 2 + 16/24 < 0$, while $\cos 0 = 1$; by
the intermediate value theorem there is a unique zero $\pi/2 \in (0, 2)$. Then $\sin(\pi/2) = 1$, so $e^{i\pi/2} = i$ and $e^{2\pi i} = i^4 = 1$.

Kernel: if $0 < t < 2\pi$, write $e^{it/4} = u + iv$ with $u, v > 0$ (as $0 < t/4 < \pi/2$). If $e^{it} = (u + iv)^4 = 1$, then $\operatorname{Im}(u + iv)^4 = 4uv(u^2 - v^2) = 0$ forces $u^2 = v^2 = 1/2$, and then
$\operatorname{Re}(u + iv)^4 = u^4 - 6u^2v^2 + v^4 = -1$, a contradiction. With periodicity, the kernel is $2\pi\ZZ$.

Surjectivity: for $a + ib \in S^1$ with $a, b \geq 0$, choose $t \in [0, \pi/2]$ with $\cos t = a$ (intermediate value theorem); then $\sin t = \sqrt{1 - a^2} = b$. The other quadrants are obtained by
multiplying by powers of $i = e^{i\pi/2}$. For the covering property, $p$ is injective on $[a, a + 1)$ and identifies only $a$ and $a + 1$ in $[a, a + 1]$, so the argument of
Proposition~\ref{homotopy:proposition-circle-covering} applies verbatim. The polar form follows by applying this to $z/|z|$.
\end{proof}

\begin{remark}\label{real-analysis:remark-exponential-covering}
The proof of Theorem~\ref{homotopy:theorem-circle} uses only that $\RR$ is simply connected and that the covering map has period $4$ with fibres the translates of $4\ZZ$; it applies verbatim to
$t \mapsto e^{2\pi it}$, with period $1$, so $\pi_1(S^1) \cong \ZZ$ is generated by the loop $t \mapsto e^{2\pi it}$. For $x > 0$ and $s \in \CC$ we define $x^s = e^{s\log x}$; then $\frac{d}{dx}x^s = sx^{s-1}$.
\end{remark}

\begin{proposition}[Notable limits]\label{real-analysis:proposition-notable-limits}
\begin{enumerate}
\item $\lim_{x \to 0}\frac{e^x - 1}{x} = 1$, $\lim_{x \to 0}\frac{\log(1 + x)}{x} = 1$, $\lim_{x \to 0}\frac{\sin x}{x} = 1$ and $\lim_{x \to 0}\frac{1 - \cos x}{x^2} = \frac12$.
\item For every $x \in \RR$, $(1 + x/n)^n \to e^x$; in particular $(1 + 1/n)^n \to e = \exp(1)$.
\item For $k \in \NN$ and $a > 0$: $x^ke^{-ax} \to 0$ and $(\log x)/x^a \to 0$ as $x \to +\infty$, and $x^a\log x \to 0$ as $x \to 0^+$.
\end{enumerate}
\end{proposition}

\begin{proof}
(1) The first three are the derivatives of $\exp$ at $0$, of $\log$ at $1$ and of $\sin$ at $0$ (Proposition~\ref{real-analysis:proposition-exponential}), which are $e^0 = 1$, $1/1$ and $\cos 0 = 1$. For the last, apply
Theorem~\ref{real-analysis:theorem-lhopital}(1): $\frac{1 - \cos x}{x^2}$ has the same limit as $\frac{\sin x}{2x} \to \frac12$.
(2) For $x = 0$ this is clear. Otherwise, for $n > |x|$, $1 + x/n > 0$ and $(1 + x/n)^n = \exp(n\log(1 + x/n))$, with $n\log(1 + x/n) = x\cdot\frac{\log(1 + x/n)}{x/n} \to x$ by (1) and
Theorem~\ref{real-analysis:theorem-sequential-limits} (as $x/n \to 0$ and $x/n \neq 0$). By continuity of $\exp$, $(1 + x/n)^n \to e^x$.
(3) For $x > 0$, $e^{ax} \geq (ax)^{k+1}/(k+1)!$, one term of the exponential series, so $0 < x^ke^{-ax} \leq (k+1)!/(a^{k+1}x) \to 0$. For $x = e^t$, $(\log x)/x^a = te^{-at}$, which tends to $0$ as $t \to +\infty$ by the
case $k = 1$; since $\log x \to +\infty$ as $x \to +\infty$, Proposition~\ref{real-analysis:proposition-composite-limit} applies. Similarly $x^a\log x = -te^{-at}$ for $x = e^{-t}$, and $t \to +\infty$ as $x \to 0^+$.
\end{proof}

\begin{counterexample}\label{real-analysis:counterexample-sin-one-over-x}
\begin{enumerate}
\item $\sin(1/x)$ has no limit at $0$: along $x_n = 1/(2\pi n + \pi/2)$ it is $1$, along $y_n = 1/(\pi n)$ it is $0$ ($n \geq 1$), and Theorem~\ref{real-analysis:theorem-sequential-limits} applies. Since
$x_n - y_n \to 0$, $\sin(1/x)$ is continuous but not uniformly continuous on $(0, 1)$. On the other hand $x\sin(1/x) \to 0$ by the squeeze theorem, as $|x\sin(1/x)| \leq |x|$.
\item A derivative need not be continuous. Let $f(x) = x^2\sin(1/x)$ for $x \neq 0$ and $f(0) = 0$. Then $f'(0) = \lim_{x \to 0}x\sin(1/x) = 0$, and $f'(x) = 2x\sin(1/x) - \cos(1/x)$ for $x \neq 0$, which has no limit at $0$.
In accordance with Proposition~\ref{real-analysis:proposition-derivative-limit} and Theorem~\ref{real-analysis:theorem-darboux}, the discontinuity of $f'$ at $0$ is essential.
\item The converse of de l'H\^opital's rule fails. With $f$ as in (2) and $g(x) = x$, $f(x)/g(x) = x\sin(1/x) \to 0$, while $f'(x)/g'(x) = 2x\sin(1/x) - \cos(1/x)$ has no limit as $x \to 0$.
\end{enumerate}
\end{counterexample}

\begin{corollary}[Inequality of arithmetic and geometric means]\label{real-analysis:corollary-am-gm}
For $x_1, \ldots, x_k > 0$ and weights $t_i \geq 0$ with $\sum t_i = 1$, $\prod_ix_i^{t_i} \leq \sum_it_ix_i$. In particular $\sqrt[k]{x_1\cdots x_k} \leq \frac{1}{k}(x_1 + \cdots + x_k)$, with equality if and only if all $x_i$ are equal.
\end{corollary}

\begin{proof}
$\log'' = -1/x^2 < 0$, so $\log$ is concave (Proposition~\ref{real-analysis:proposition-convex}(3)), and by Jensen's inequality $\sum t_i\log x_i \leq \log\sum t_ix_i$. Apply $\exp$, which is increasing. For
the equality statement, let $m = \frac1k\sum x_i$. As $\log'$ is strictly decreasing, the mean value theorem gives $\log x < \log m + (x - m)/m$ for every $x > 0$ with $x \neq m$ (the tangent inequality of
Proposition~\ref{real-analysis:proposition-convex}(2) for $-\log$, made strict). Averaging over $x = x_1, \ldots, x_k$ gives $\frac1k\sum\log x_i < \log m$ unless all $x_i = m$.
\end{proof}

\begin{lemma}\label{real-analysis:lemma-bump}
The function $\beta(t) = e^{-1/t}$ for $t > 0$ and $\beta(t) = 0$ for $t \leq 0$ has derivatives of all orders on $\RR$, all vanishing at $0$. The function
$\gamma(t) = \beta(t)/(\beta(t) + \beta(1 - t))$ has derivatives of all orders, takes values in $[0, 1]$, and equals $0$ for $t \leq 0$ and $1$ for $t \geq 1$.
\end{lemma}

\begin{proof}
By induction, on $t > 0$, $\beta^{(k)}(t) = p_k(1/t)e^{-1/t}$ for a polynomial $p_k$. Since $e^{s} \geq s^{m}/m!$ for $s > 0$, $s^je^{-s} \to 0$ as $s \to \infty$; so $\beta^{(k)}(t) \to 0$ as $t \to 0^+$, and
the difference quotients of $\beta^{(k)}$ at $0$ tend to $0$, which by induction shows $\beta^{(k+1)}(0) = 0$. The denominator of $\gamma$ never vanishes (one of $t$, $1 - t$ is positive), so $\gamma$ has derivatives of all orders by
Proposition~\ref{real-analysis:proposition-derivative-rules}, and the other properties are clear.
\end{proof}

\begin{definition}\label{real-analysis:definition-analytic}
A function $f \colon U \to \RR$ on an open set $U \subseteq \RR$ is \emph{(real) analytic} if every $a \in U$ has a neighbourhood on which $f(x) = \sum_nc_n(x - a)^n$ for a convergent power series.
\end{definition}

\begin{proposition}\label{real-analysis:proposition-analytic}
\begin{enumerate}
\item If $f(x) = \sum_na_nx^n$ has radius of convergence $R > 0$ and $|b| < R$, then $f(b + h) = \sum_k\frac{f^{(k)}(b)}{k!}h^k$ for $|h| < R - |b|$. So $f$ is analytic on $(-R, R)$.
\item An analytic function is smooth, and near each point it equals its Taylor series there.
\item (Identity theorem) Let $f$ be analytic on an open interval $I$. If all derivatives of $f$ vanish at one point, or if the zeros of $f$ have an accumulation point in $I$, then $f = 0$ on $I$.
\end{enumerate}
\end{proposition}

\begin{proof}
(1) $\sum_n\sum_{k \leq n}|a_n|\binom nk|b|^{n-k}|h|^k = \sum_n|a_n|(|b| + |h|)^n < \infty$, so the double series $\sum_{n,k}a_n\binom nkb^{n-k}h^k$ converges absolutely and may be summed in any order
(Proposition~\ref{real-analysis:proposition-series}(1)): summing over $k$ first gives $f(b + h)$, and summing over $n$ first gives $\sum_kc_kh^k$ with $c_k = \sum_{n \geq k}a_n\binom nkb^{n-k}$, which is
$f^{(k)}(b)/k!$ by Theorem~\ref{real-analysis:theorem-power-series-derivative}. (2) follows from Theorem~\ref{real-analysis:theorem-power-series-derivative} applied to the local expansions.

(3) Let $S$ be the set of points of $I$ where all derivatives of $f$ vanish. It is closed in $I$, as the derivatives are continuous, and open: near $x \in S$, $f$ equals its Taylor series at $x$, which is
$0$, so $f = 0$ near $x$. As $I$ is connected (Theorem~\ref{general-topology:theorem-intervals-connected}), $S = \emptyset$ or $S = I$. If the zeros accumulate at $a \in I$ and $f(x) = \sum c_n(x - a)^n$
near $a$ with some $c_n \neq 0$, let $m$ be least with $c_m \neq 0$; then $f(x) = (x - a)^mg(x)$ with $g$ continuous and $g(a) = c_m \neq 0$, so $f$ has no zeros near $a$ other than $a$, a contradiction.
So all derivatives vanish at $a$, and $S = I$.
\end{proof}

Polynomials, $\exp$, $\sin$ and $\cos$ are analytic on $\RR$, by (1), with radius $\infty$. The function $1/(1 + x^2)$ is smooth on $\RR$, but its Taylor series at $0$, $\sum(-1)^nx^{2n}$, has radius only
$1$: the obstruction is invisible on $\RR$, and comes from the complex points $\pm i$ where $1/(1 + z^2)$ is not defined; it is nevertheless analytic on all of $\RR$, as holomorphic functions are analytic
(Chapter~\ref{complex-analysis}). The identity theorem fails for smooth functions:

\begin{counterexample}\label{real-analysis:counterexample-not-analytic}
The function $\beta$ of Lemma~\ref{real-analysis:lemma-bump} has all derivatives zero at $0$, so its Taylor series at $0$ is zero, but $\beta(t) \neq 0$ for $t > 0$. Smooth functions need not be
given by their Taylor series; complex differentiable functions are (Chapter~\ref{complex-analysis}).
\end{counterexample}

\section{The Riemann integral}\label{real-analysis:section-riemann}

\begin{definition}\label{real-analysis:definition-riemann}
Let $f \colon [a, b] \to \RR$ be bounded. For a partition $P \colon a = t_0 < \cdots < t_k = b$ let $U(f, P) = \sum_i(t_i - t_{i-1})\sup_{[t_{i-1}, t_i]}f$ and $L(f, P)$ the same with $\inf$. The function is
\emph{Riemann integrable} if $\sup_PL(f, P) = \inf_PU(f, P)$, and this common value is $\int_a^bf = \int_a^bf(t)\,dt$. A $\CC^m$-valued function is integrable if its real and imaginary components are, with
the integral taken componentwise. We set $\int_b^af = -\int_a^bf$.
\end{definition}

\begin{proposition}\label{real-analysis:proposition-riemann}
\begin{enumerate}
\item If $P'$ refines $P$, then $L(f, P) \leq L(f, P') \leq U(f, P') \leq U(f, P)$; and $L(f, P) \leq U(f, Q)$ for all $P, Q$. So $f$ is integrable if and only if for every $\varepsilon > 0$ there is $P$ with
$U(f, P) - L(f, P) < \varepsilon$.
\item Continuous functions are integrable, and $\sum_i(t_i - t_{i-1})f(s_i) \to \int_a^bf$ for any choice of $s_i \in [t_{i-1}, t_i]$ as $\max_i(t_i - t_{i-1}) \to 0$.
\item For continuous $f, g$: the integral is linear, $f \leq g$ implies $\int f \leq \int g$, $\int_a^cf = \int_a^bf + \int_b^cf$, $|\int_a^bf| \leq \int_a^b|f|$ for $\CC^m$-valued $f$ with the Euclidean
norm, and if $f_n \to f$ uniformly, $\int f_n \to \int f$.
\end{enumerate}
\end{proposition}

\begin{proof}
(1) Adding a point to a partition can only increase infima and decrease suprema on the pieces; for $P, Q$ compare both with $P \cup Q$. (2) By uniform continuity
(Corollary~\ref{general-topology:corollary-uniform-continuity}), for $\varepsilon > 0$ there is $\delta$ such that $f$ varies by less than $\varepsilon$ on intervals of length $< \delta$; for partitions of mesh
$< \delta$, $U - L \leq \varepsilon(b - a)$, and Riemann sums lie between $L$ and $U$. (3) Linearity and additivity follow from (2), monotonicity from $L(f, P) \leq L(g, P)$, and uniform limits from
$|\int(f_n - f)| \leq (b - a)\sup|f_n - f|$. For the norm inequality let $v = \int f \in \CC^m$; then
$\|v\|^2 = \operatorname{Re}\langle \int f, v \rangle = \int\operatorname{Re}\langle f(t), v \rangle\,dt \leq \int\|f\|\,\|v\|$ by Cauchy--Schwarz (Proposition~\ref{linear-algebra:proposition-cauchy-schwarz}).
\end{proof}

\begin{theorem}[Fundamental theorem of calculus]\label{real-analysis:theorem-ftc}
Let $f \colon [a, b] \to \CC^m$ be continuous. Then $F(x) = \int_a^xf$ is differentiable with $F' = f$. If $G \colon [a, b] \to \CC^m$ is continuously differentiable, then $\int_a^bG' = G(b) - G(a)$.
\end{theorem}

\begin{proof}
$|(F(x + h) - F(x))/h - f(x)| = |\frac1h\int_x^{x+h}(f - f(x))| \leq \sup_{|t - x| \leq |h|}|f(t) - f(x)| \to 0$. For the second statement, $\int_a^xG' - G(x)$ has derivative $0$, hence is constant by
Theorem~\ref{real-analysis:theorem-mean-value} applied to real and imaginary components.
\end{proof}

\begin{corollary}\label{real-analysis:corollary-integration-rules}
For continuously differentiable $f, g$ and $\varphi \colon [c, d] \to [a, b]$: $\int_a^bf'g = [fg]_a^b - \int_a^bfg'$ and $\int_c^df(\varphi(s))\varphi'(s)\,ds = \int_{\varphi(c)}^{\varphi(d)}f$ for continuous $f$.
Moreover $\|G(b) - G(a)\| \leq (b - a)\sup\|G'\|$ for $G$ continuously differentiable.
\end{corollary}

\begin{proof}
Apply Theorem~\ref{real-analysis:theorem-ftc} to $fg$, to $F \circ \varphi$ with $F(x) = \int_a^xf$, and use the norm inequality of Proposition~\ref{real-analysis:proposition-riemann}.
\end{proof}

The first formula of the corollary is \emph{integration by parts}: it moves the derivative from one factor to the other, at the price of a boundary term. It is useful whenever differentiating
one factor simplifies it while integrating the other does not make it worse.

\begin{example}\label{real-analysis:example-integration-by-parts}
\begin{enumerate}
\item With $f(x) = e^x$ and $g(x) = x$, so that $f^{\prime}g = xe^x$: $\int_0^1xe^x\,dx = [xe^x]_0^1 - \int_0^1e^x\,dx = e - (e - 1) = 1$.
\item With $f(x) = x$ and $g = \log$, so that $f^{\prime}g = \log x$ and $fg^{\prime} = 1$: $\int_1^e\log x\,dx = [x\log x]_1^e - \int_1^e1\,dx = e - (e - 1) = 1$.
\item \emph{Wallis' integrals.} Let $I_n = \int_0^{\pi/2}\sin^nx\,dx$, so $I_0 = \frac\pi2$ and $I_1 = 1$. For $n \geq 2$ take $f(x) = -\cos x$ and $g(x) = \sin^{n-1}x$: the boundary term
$[-\cos x\sin^{n-1}x]_0^{\pi/2}$ vanishes, and
\[
I_n = (n - 1)\int_0^{\pi/2}\cos^2x\sin^{n-2}x\,dx = (n - 1)(I_{n-2} - I_n), \qquad\text{so}\qquad I_n = \frac{n - 1}{n}I_{n-2} .
\]
Hence $I_{2m} = \frac{1 \cdot 3 \cdots (2m-1)}{2 \cdot 4 \cdots (2m)}\cdot\frac\pi2$ and $I_{2m+1} = \frac{2 \cdot 4 \cdots (2m)}{3 \cdot 5 \cdots (2m+1)}$.
\item \emph{Periodic functions.} If $f$ and $g$ are continuously differentiable and $2\pi$-periodic, the boundary term $[fg]_0^{2\pi}$ vanishes, and $\int_0^{2\pi}f^{\prime}g = -\int_0^{2\pi}fg^{\prime}$:
without boundary, the derivative can be moved from one factor to the other at the cost of a sign. In several variables this becomes $\int(\partial_if)g = -\int f\,\partial_ig$ for compactly supported
functions (Chapter~\ref{measure-theory}), and on a region with boundary the boundary term becomes an integral over the boundary (Green's identities, Chapter~\ref{riemannian-geometry}).
\end{enumerate}
\end{example}

\begin{proposition}[Mean value theorem for integrals]\label{real-analysis:proposition-mean-value-integrals}
Let $f, g \colon [a, b] \to \RR$ be continuous with $g \geq 0$. Then there is $c \in [a, b]$ with $\int_a^bfg = f(c)\int_a^bg$. In particular $\int_a^bf = f(c)(b - a)$ for some $c$: the average
value of a continuous function is attained.
\end{proposition}

\begin{proof}
Let $m$, $M$ be the minimum and maximum of $f$ (Corollary~\ref{general-topology:corollary-extreme-value}). Then $mg \leq fg \leq Mg$, so $m\int g \leq \int fg \leq M\int g$. If $\int g = 0$ any $c$ works;
otherwise $\int fg/\int g \in [m, M]$ is a value of $f$ by the intermediate value theorem (Corollary~\ref{general-topology:corollary-ivt}).
\end{proof}

The hypotheses matter: for $g(x) = x$ on $[-1, 1]$ and $f = g$, $\int fg = \frac23$ while $\int g = 0$; and for vector-valued $f$ the statement fails, since $f(t) = e^{it}$ on $[0, 2\pi]$ has integral $0$
but never vanishes.

\begin{example}\label{real-analysis:example-computing-integrals}
\emph{Computing integrals.} The fundamental theorem of calculus reduces integration to finding antiderivatives, and the rules of Corollary~\ref{real-analysis:corollary-integration-rules} are the
main tools.
\begin{enumerate}
\item \emph{Substitution.} With $u = 1 + x^2$, $du = 2x\,dx$: $\int_0^1\frac{x}{1 + x^2}\,dx = \frac12\int_1^2\frac{du}{u} = \frac12\log 2$.
\item \emph{Partial fractions.} $\frac{1}{x^2 - 1} = \frac12\bigl(\frac{1}{x - 1} - \frac{1}{x + 1}\bigr)$, so $\int_2^3\frac{dx}{x^2 - 1} = \frac12\bigl[\log\frac{x - 1}{x + 1}\bigr]_2^3 = \frac12\log\frac32$.
In general every rational function is a polynomial plus a combination of terms $(x - a)^{-k}$ and $(Ax + B)((x - p)^2 + q^2)^{-k}$, whose antiderivatives involve only rational functions,
$\log$ and $\arctan$.
\item \emph{Trigonometric substitution.} With $x = \sin\theta$: $\int_{-1}^1\sqrt{1 - x^2}\,dx = \int_{-\pi/2}^{\pi/2}\cos^2\theta\,d\theta = \frac\pi2$, the area of a half disc; here
$\cos^2\theta = \frac12(1 + \cos 2\theta)$.
\item \emph{Symmetry.} $\int_{-a}^af = 0$ for odd $f$, as the substitution $x \mapsto -x$ shows; for instance $\int_{-1}^1x^3e^{x^2}\,dx = 0$ without computing an antiderivative.
\item \emph{Recursion by parts}, as for Wallis' integrals in Example~\ref{real-analysis:example-integration-by-parts}: $\int_0^1x^ne^x\,dx = e - n\int_0^1x^{n-1}e^x\,dx$.
\end{enumerate}
Many integrals have no elementary antiderivative, such as $\int e^{-x^2}\,dx$; they are computed by other means, such as polar coordinates for $\int_\RR e^{-x^2} = \sqrt\pi$ and differentiation with
respect to a parameter (Chapter~\ref{measure-theory}).
\end{example}

\begin{theorem}[Stirling's formula]\label{real-analysis:theorem-stirling}
$n! \sim \sqrt{2\pi n}\,(n/e)^n$, that is, $n!\,e^n/(n^{n + 1/2}\sqrt{2\pi}) \to 1$ as $n \to \infty$.
\end{theorem}

\begin{proof}
\emph{Step 1: $a_n = n!\,e^n/n^{n + 1/2}$ converges to some $C > 0$.} For $t \geq 0$, $t - \frac{t^2}2 \leq \log(1 + t) \leq t - \frac{t^2}2 + \frac{t^3}3$: the differences vanish at $0$ and have
derivatives $\frac{t^2}{1 + t} \geq 0$ and $\frac{t^3}{1 + t} \geq 0$. Now $\log(a_n/a_{n+1}) = (n + \frac12)\log(1 + \frac1n) - 1$, and with $t = 1/n$ the two bounds give
$-\frac{1}{4n^2} \leq \log(a_n/a_{n+1}) \leq \frac{1}{12n^2} + \frac{1}{6n^3} \leq \frac{1}{4n^2}$. So $\sum_n\log(a_n/a_{n+1})$ converges absolutely, $\log a_n$ converges, and $a_n \to C > 0$.

\emph{Step 2: $C = \sqrt{2\pi}$.} With Wallis' integrals $I_k$ (Example~\ref{real-analysis:example-integration-by-parts}), $\sin^{2m+1} \leq \sin^{2m} \leq \sin^{2m-1}$ on $[0, \frac\pi2]$ gives
$I_{2m+1} \leq I_{2m} \leq I_{2m-1} = \frac{2m+1}{2m}I_{2m+1}$, so $I_{2m}/I_{2m+1} \to 1$. Using $2 \cdot 4 \cdots 2m = 2^mm!$ and $1 \cdot 3 \cdots (2m - 1) = (2m)!/(2^mm!)$, this says
\[
\frac\pi2 = \lim_{m \to \infty}\frac{2^{4m}(m!)^4}{((2m)!)^2(2m + 1)}
\]
(Wallis' product). Substituting $m! = a_mm^{m + 1/2}e^{-m}$ and $(2m)! = a_{2m}(2m)^{2m + 1/2}e^{-2m}$, the right side is $\frac{a_m^4}{a_{2m}^2}\cdot\frac{m}{2(2m + 1)} \to \frac{C^2}{4}$. Hence
$C^2 = 2\pi$.
\end{proof}

For example $10! = 3628800$, while $\sqrt{20\pi}(10/e)^{10} \approx 3598696$, an error of less than $1\%$; the relative error is about $1/(12n)$, in accordance with the bound in Step 1.

\begin{proposition}\label{real-analysis:proposition-parameter-integrals}
Let $f \colon [c, d] \times [a, b] \to \CC$ be continuous.
\begin{enumerate}
\item $x \mapsto \int_a^bf(x, t)\,dt$ is continuous, and if $\partial_xf$ exists and is continuous on $[c, d] \times [a, b]$, it is differentiable with derivative $\int_a^b\partial_xf(x, t)\,dt$.
\item (Fubini for continuous functions) $\int_c^d\int_a^bf(x, t)\,dt\,dx = \int_a^b\int_c^df(x, t)\,dx\,dt$.
\end{enumerate}
\end{proposition}

\begin{proof}
(1) Continuity follows from uniform continuity of $f$. For the derivative, by the mean value theorem applied to real and imaginary parts, $(f(x + h, t) - f(x, t))/h - \partial_xf(x, t)$ is bounded by the
oscillation of $\partial_xf$ over $|s - x| \leq |h|$, which tends to $0$ uniformly in $t$ by uniform continuity. (2) Let $\Phi(y) = \int_c^y\int_a^bf\,dt\,dx$ and $\Psi(y) = \int_a^b\int_c^yf\,dx\,dt$. By
Theorem~\ref{real-analysis:theorem-ftc} and (1) (applied to $g(y, t) = \int_c^yf(x, t)\,dx$, with $\partial_yg = f$ continuous), $\Phi' = \Psi' = \int_a^bf(y, t)\,dt$, and $\Phi(c) = \Psi(c) = 0$.
\end{proof}

\section{Sequences of functions}\label{real-analysis:section-function-sequences}

\begin{definition}\label{real-analysis:definition-uniform-convergence}
Let $D$ be a set and $f_n, f \colon D \to \mathbb{K}^m$. Then $f_n \to f$ \emph{pointwise} if $f_n(x) \to f(x)$ for every $x \in D$, and \emph{uniformly} if $\|f_n - f\|_\infty = \sup_{x \in D}\|f_n(x) - f(x)\| \to 0$.
\end{definition}

The difference is the order of quantifiers: for pointwise convergence the threshold $N$ may depend on $x$, for uniform convergence it may not. Uniform convergence implies pointwise convergence.

\begin{proposition}\label{real-analysis:proposition-uniform-cauchy}
\begin{enumerate}
\item (Cauchy criterion) $(f_n)$ converges uniformly on $D$ if and only if for every $\varepsilon > 0$ there is $N$ with $\|f_n - f_m\|_\infty < \varepsilon$ for $m, n \geq N$.
\item If $D$ is a metric space, the $f_n$ are continuous and $f_n \to f$ uniformly, then $f$ is continuous.
\item If $f_n \to f$ uniformly on $[a, b]$ and the $f_n$ are continuous, then $\int_a^bf_n \to \int_a^bf$.
\end{enumerate}
\end{proposition}

\begin{proof}
(1) If $f_n \to f$ uniformly, use the triangle inequality. Conversely, for each $x$ the sequence $(f_n(x))$ is Cauchy, hence converges to some $f(x)$ (Theorem~\ref{real-analysis:theorem-cauchy-criterion}, componentwise).
Letting $m \to \infty$ in $\|f_n(x) - f_m(x)\| < \varepsilon$ gives $\|f_n(x) - f(x)\| \leq \varepsilon$ for $n \geq N$ and all $x$. (2) is Proposition~\ref{general-topology:proposition-uniform-limit}. (3) is
Proposition~\ref{real-analysis:proposition-riemann}(3).
\end{proof}

\begin{counterexample}\label{real-analysis:counterexample-pointwise}
\begin{enumerate}
\item $f_n(x) = x^n$ on $[0, 1]$ converges pointwise to the function equal to $0$ on $[0, 1)$ and $1$ at $1$, which is discontinuous; so the convergence is not uniform. On $[0, 1)$ the limit $0$ is continuous, but
$\sup_{[0,1)}x^n = 1$, so the convergence is still not uniform. On $[0, r]$ with $r < 1$ it is uniform, as $\sup = r^n \to 0$.
\item Pointwise convergence does not commute with integration. Let $f_n(x) = nx(1 - x^2)^n$ on $[0, 1]$. For $0 < x \leq 1$, $f_n(x) \to 0$ by Proposition~\ref{real-analysis:proposition-fundamental-limits}(3) with $a = (1 - x^2)^{-1}$
(or trivially if $x = 1$), and $f_n(0) = 0$. But, substituting $u = 1 - x^2$, $\int_0^1f_n = \frac n2\int_0^1u^n\,du = \frac{n}{2(n + 1)} \to \frac12 \neq 0$.
\item Uniform convergence does not commute with differentiation. $f_n(x) = x^n/n$ converges uniformly to $0$ on $[0, 1]$, since $|f_n| \leq 1/n$, but $f_n'(1) = 1$ for all $n$.
\end{enumerate}
\end{counterexample}

\begin{theorem}\label{real-analysis:theorem-uniform-derivative}
Let $f_n \colon [a, b] \to \RR$ be continuously differentiable. Suppose that $(f_n')$ converges uniformly to $g$ and that $(f_n(x_0))$ converges for some $x_0 \in [a, b]$. Then $(f_n)$ converges uniformly to a continuously
differentiable $f$ with $f' = g$.
\end{theorem}

\begin{proof}
The function $g$ is continuous by Proposition~\ref{real-analysis:proposition-uniform-cauchy}(2). Let $L = \lim f_n(x_0)$ and $f(x) = L + \int_{x_0}^xg$, which is continuously differentiable with $f' = g$ by
Theorem~\ref{real-analysis:theorem-ftc}. By the same theorem $f_n(x) = f_n(x_0) + \int_{x_0}^xf_n'$, so
$|f_n(x) - f(x)| \leq |f_n(x_0) - L| + (b - a)\|f_n' - g\|_\infty$, which tends to $0$ uniformly in $x$.
\end{proof}

\begin{theorem}[Dini]\label{real-analysis:theorem-dini}
Let $K$ be a compact metric space and $f_n \colon K \to \RR$ continuous functions with $f_1 \leq f_2 \leq \cdots$, converging pointwise to a continuous function $f$. Then the convergence is uniform.
\end{theorem}

\begin{proof}
Let $g_n = f - f_n$, continuous, $\geq 0$, decreasing in $n$ and tending to $0$ pointwise. For $\varepsilon > 0$ the sets $K_n = \{x : g_n(x) \geq \varepsilon\}$ are closed, decreasing, and have empty intersection. By
compactness some $K_N$ is empty (the complements form an open cover with a finite subcover, and the $K_n$ decrease), and then $0 \leq g_n < \varepsilon$ on $K$ for $n \geq N$.
\end{proof}

\begin{counterexample}\label{real-analysis:counterexample-dini}
Each hypothesis of Dini's theorem is needed. On the noncompact $[0, 1)$, $f_n(x) = -x^n$ increases to the continuous function $0$, not uniformly. On $[0, 1]$, $f_n(x) = -x^n$ increases to a discontinuous limit, not
uniformly. And without monotonicity, the continuous functions $f_n(x) = \max(0, 1 - |nx - 2|)$, bumps of height $1$ moving towards $0$, converge pointwise to $0$ on $[0, 1]$ but not uniformly.
\end{counterexample}

\begin{corollary}[Series of functions]\label{real-analysis:corollary-series-functions}
Let $u_n \colon [a, b] \to \RR$ be continuous and suppose $\sum_n\|u_n\|_\infty < \infty$. Then $\sum_nu_n$ converges uniformly to a continuous function, and $\int_a^b\sum_nu_n = \sum_n\int_a^bu_n$. If moreover
the $u_n$ are continuously differentiable and $\sum_n\|u_n'\|_\infty < \infty$, then $\sum_nu_n$ is continuously differentiable with derivative $\sum_nu_n'$.
\end{corollary}

\begin{proof}
Apply the Weierstrass $M$-test (Proposition~\ref{real-analysis:proposition-series}(5)) and Propositions~\ref{real-analysis:proposition-uniform-cauchy}(2),(3) to the partial sums, and Theorem~\ref{real-analysis:theorem-uniform-derivative}
to the partial sums and their derivatives.
\end{proof}

\begin{theorem}[A continuous nowhere differentiable function]\label{real-analysis:theorem-takagi}
Let $d(x)$ be the distance from $x \in \RR$ to the nearest integer. The \emph{Takagi function} $T(x) = \sum_{m \geq 0}2^{-m}d(2^mx)$ is continuous on $\RR$ and differentiable at no point.
\end{theorem}

\begin{proof}
Since $0 \leq d \leq \frac12$ and $d$ is continuous (it is $1$-Lipschitz), the series converges uniformly by the $M$-test, and $T$ is continuous.

We first note: if $f$ is differentiable at $x$ and $u_n \leq x \leq v_n$ with $u_n < v_n$ and $v_n - u_n \to 0$, then $\frac{f(v_n) - f(u_n)}{v_n - u_n} \to f'(x)$. Indeed, with $\lambda_n = (v_n - x)/(v_n - u_n) \in [0, 1]$,
\[
\frac{f(v_n) - f(u_n)}{v_n - u_n} = \lambda_n\frac{f(v_n) - f(x)}{v_n - x} + (1 - \lambda_n)\frac{f(x) - f(u_n)}{x - u_n},
\]
where a term with weight $0$ (when $v_n = x$ or $u_n = x$) is omitted; both quotients tend to $f'(x)$, and so does their convex combination.

Fix $x$, and for $n \geq 1$ let $u_n = k_n/2^n \leq x < v_n = (k_n + 1)/2^n$ with $k_n \in \ZZ$. For $m \geq n$, $2^mu_n$ and $2^mv_n$ are integers, so $d(2^mu_n) = d(2^mv_n) = 0$. For $m < n$, the numbers $2^mu_n$ and $2^mv_n$ are
consecutive multiples of $2^{m-n} \leq \frac12$, so they lie in a common interval $[\frac j2, \frac{j+1}2]$, on which $d$ is affine with slope $\pm1$. Hence
\[
q_n = \frac{T(v_n) - T(u_n)}{v_n - u_n} = \sum_{m=0}^{n-1}2^{-m}\frac{d(2^mv_n) - d(2^mu_n)}{v_n - u_n} = \sum_{m=0}^{n-1}\varepsilon_{m,n}, \qquad \varepsilon_{m,n} \in \{1, -1\}.
\]
So $q_n$ is an integer of the same parity as $n$, and $(q_n)$ cannot converge: consecutive terms differ by an odd integer, so $|q_{n+1} - q_n| \geq 1$. By the preliminary remark, $T$ is not differentiable at $x$.
\end{proof}

\section{Derivatives in several variables}\label{real-analysis:section-several-variables}

\begin{lemma}\label{real-analysis:lemma-norms}
All norms on a finite-dimensional real or complex vector space are equivalent (they define the same topology), and every linear map $A \colon V \to W$ between finite-dimensional normed spaces is
continuous, with finite \emph{operator norm} $\|A\| = \sup_{\|v\| \leq 1}\|Av\|$. Then $\|Av\| \leq \|A\|\,\|v\|$ and $\|AB\| \leq \|A\|\,\|B\|$.
\end{lemma}

\begin{proof}
Identify $V$ with $\RR^n$ by a basis and let $\|\cdot\|$ be a norm. By the triangle inequality and Cauchy--Schwarz, $\|x\| \leq \sum|x_i|\,\|e_i\| \leq C\|x\|_2$, so $\|\cdot\|$ is continuous for the
Euclidean topology. On the compact unit sphere of $\|\cdot\|_2$ it attains a minimum $c > 0$, so $c\|x\|_2 \leq \|x\| \leq C\|x\|_2$, and Lemma~\ref{general-topology:lemma-equivalent-metrics} applies. A linear map
is Lipschitz for Euclidean norms, since $\|Ax\| \leq \sum|x_i|\,\|Ae_i\|$. The inequalities follow from the definition.
\end{proof}

\begin{definition}\label{real-analysis:definition-differentiable}
Let $U \subseteq \RR^n$ be open and $f \colon U \to \RR^m$. Then $f$ is \emph{differentiable} at $a \in U$ if there is a linear map $Df(a) \colon \RR^n \to \RR^m$ with
\[
f(a + h) = f(a) + Df(a)h + r(h), \qquad \|r(h)\|/\|h\| \to 0 \text{ as } h \to 0.
\]
The \emph{partial derivatives} are $\partial_jf(a) = \frac{d}{dt}f(a + te_j)|_{t=0}$. The map $f$ is $C^1$ if it is differentiable everywhere and $Df \colon U \to \Hom(\RR^n, \RR^m)$ is continuous; $C^k$ if $Df$ is
$C^{k-1}$; and \emph{smooth} ($C^\infty$) if it is $C^k$ for all $k$. A \emph{$C^k$-diffeomorphism} is a bijection between open sets which is $C^k$ with $C^k$ inverse.
\end{definition}

\begin{proposition}\label{real-analysis:proposition-derivative-basic}
\begin{enumerate}
\item $Df(a)$ is unique, $f$ is continuous at $a$, and the matrix of $Df(a)$ is the \emph{Jacobian matrix} $(\partial_jf_i(a))$.
\item (Chain rule) If $f$ is differentiable at $a$ and $g$ at $f(a)$, then $g \circ f$ is differentiable at $a$ and $D(g \circ f)(a) = Dg(f(a)) \circ Df(a)$.
\item If all partial derivatives of $f$ exist near $a$ and are continuous at $a$, then $f$ is differentiable at $a$. So $f$ is $C^1$ if and only if its partial derivatives exist and are continuous, and
$C^k$ if and only if all partial derivatives of order at most $k$ exist and are continuous.
\item (Mean value inequality) If $f$ is $C^1$ and the segment $[x, y]$ lies in $U$, then $\|f(y) - f(x)\| \leq \sup_{z \in [x, y]}\|Df(z)\|\,\|y - x\|$.
\end{enumerate}
\end{proposition}

\begin{proof}
(1) If $A$ and $B$ both work, $\|(A - B)h\|/\|h\| \to 0$, so $A = B$ by linearity. Taking $h = te_j$ gives the partial derivatives. (2) Let $b = f(a)$, $k = f(a + h) - b = Df(a)h + r_f(h)$. Then
$g(f(a + h)) - g(b) - Dg(b)Df(a)h = Dg(b)r_f(h) + r_g(k)$. For small $h$, $\|k\| \leq (\|Df(a)\| + 1)\|h\|$, so both terms are $o(\|h\|)$. (3) We may assume $m = 1$. Write
$f(a + h) - f(a) = \sum_j[f(a + h_1e_1 + \cdots + h_je_j) - f(a + h_1e_1 + \cdots + h_{j-1}e_{j-1})]$. By the mean value theorem the $j$-th term is $h_j\partial_jf(\xi_j)$ with $\xi_j$ within $\|h\|$ of $a$, and
$\partial_jf(\xi_j) \to \partial_jf(a)$. The statement about $C^k$ follows by induction. (4) Let $g(t) = f(x + t(y - x))$; by the chain rule $g'(t) = Df(x + t(y - x))(y - x)$, and
Corollary~\ref{real-analysis:corollary-integration-rules} applies.
\end{proof}

\begin{counterexample}\label{real-analysis:counterexample-partials}
Let $f(x, y) = xy/(x^2 + y^2)$ for $(x, y) \neq 0$ and $f(0) = 0$. Both partial derivatives exist everywhere ($f$ vanishes on the axes), but $f(t, t) = 1/2$ for $t \neq 0$, so $f$ is not continuous at $0$, let
alone differentiable.
\end{counterexample}

\begin{theorem}[Schwarz]\label{real-analysis:theorem-schwarz}
If $f \colon U \to \RR$ is $C^2$, then $\partial_i\partial_jf = \partial_j\partial_if$.
\end{theorem}

\begin{proof}
Restricting to the plane of the $i$-th and $j$-th coordinates, we may assume $n = 2$. For a rectangle $R = [x_0, x_1] \times [y_0, y_1] \subseteq U$, Theorem~\ref{real-analysis:theorem-ftc} applied twice gives
$\int_{y_0}^{y_1}\int_{x_0}^{x_1}\partial_x\partial_yf\,dx\,dy = f(x_1, y_1) - f(x_0, y_1) - f(x_1, y_0) + f(x_0, y_0)$, and the same for $\int_{x_0}^{x_1}\int_{y_0}^{y_1}\partial_y\partial_xf\,dy\,dx$. By
Proposition~\ref{real-analysis:proposition-parameter-integrals}(2), the continuous function $g = \partial_x\partial_yf - \partial_y\partial_xf$ has integral $0$ over every small rectangle. If $g(p) > 0$, then $g > g(p)/2$ on a small
rectangle around $p$, and its integral is positive, a contradiction.
\end{proof}

\begin{counterexample}\label{real-analysis:counterexample-schwarz}
For $f(x, y) = xy(x^2 - y^2)/(x^2 + y^2)$, $f(0) = 0$, one computes $\partial_xf(0, y) = -y$ and $\partial_yf(x, 0) = x$, so $\partial_y\partial_xf(0) = -1 \neq 1 = \partial_x\partial_yf(0)$. The second partial derivatives
exist but are not continuous at $0$.
\end{counterexample}

\begin{theorem}[Taylor's theorem in several variables]\label{real-analysis:theorem-taylor-several}
Let $f \colon U \to \RR$ be $C^{k+1}$ and $[a, a + h] \subseteq U$. With multi-index notation $\partial^\alpha = \partial_1^{\alpha_1} \cdots \partial_n^{\alpha_n}$, $h^\alpha = h_1^{\alpha_1} \cdots h_n^{\alpha_n}$ and
$\alpha! = \alpha_1! \cdots \alpha_n!$, there is $\xi \in (a, a + h)$ with
\[
f(a + h) = \sum_{|\alpha| \leq k}\frac{\partial^\alpha f(a)}{\alpha!}h^\alpha + \sum_{|\alpha| = k+1}\frac{\partial^\alpha f(\xi)}{\alpha!}h^\alpha.
\]
\end{theorem}

\begin{proof}
Apply Theorem~\ref{real-analysis:theorem-taylor-one} to $g(t) = f(a + th)$ on $[0, 1]$. By the chain rule and Theorem~\ref{real-analysis:theorem-schwarz},
$g^{(j)}(t) = \sum_{|\alpha| = j}\frac{j!}{\alpha!}\partial^\alpha f(a + th)h^\alpha$ (the multinomial formula).
\end{proof}

\begin{corollary}\label{real-analysis:corollary-bump-function}
For $0 < r < R$ there is a smooth function $\RR^n \to [0, 1]$ equal to $1$ on the closed ball of radius $r$ and to $0$ outside the open ball of radius $R$.
\end{corollary}

\begin{proof}
Take $x \mapsto 1 - \gamma((\|x\|^2 - r^2)/(R^2 - r^2))$ with $\gamma$ as in Lemma~\ref{real-analysis:lemma-bump}; it is smooth by the chain rule, since $\|x\|^2$ is a polynomial.
\end{proof}

\section{Curves, arc length and line integrals}\label{real-analysis:section-curves}

A curve in $\RR^n$ can be studied through a parametrisation, a map $t \mapsto \gamma(t)$, but many of the quantities we care about belong to the path traced out and not to the particular way of running
along it: its length, the mass of a wire bent along it, the work done by a force field along it. This section defines these quantities, computes them by integrals, and checks that they do not
depend on the parametrisation. Two kinds of integrals appear. The integral of a function with respect to arc length, $\int_\gamma f\,ds$, does not see the direction of travel; the integral of a vector
field, $\int_\gamma F \cdot d\gamma$, does, and changes sign when the curve is run backwards. Traditionally they are called line integrals of the first and of the second kind. The first is the special case
$f = 1$ of a mass, the second the special case of a $1$-form; both reappear on manifolds (Chapters~\ref{differential-forms} and~\ref{riemannian-geometry}). Throughout, $|x|$ is the Euclidean norm on $\RR^n$ and
$\langle x, y\rangle = x \cdot y$ the standard inner product.

\subsection{Curves and their length}\label{real-analysis:subsection-curve-length}

\begin{definition}\label{real-analysis:definition-curve}
A \emph{curve} in $\RR^n$ is a continuous map $\gamma \colon [a, b] \to \RR^n$, with $a < b$; $\gamma(a)$ and $\gamma(b)$ are its \emph{initial} and \emph{final points}, and it is \emph{closed} if they agree. It is of
class $C^1$ if it is differentiable with continuous derivative $\gamma^{\prime} = (\gamma_1^{\prime}, \ldots, \gamma_n^{\prime})$, one-sided at the ends, and \emph{piecewise $C^1$} if there is a subdivision
$a = c_0 < c_1 < \cdots < c_m = b$ such that each restriction $\gamma|_{[c_{k-1}, c_k]}$ is $C^1$. A $C^1$ curve is \emph{regular} if $\gamma^{\prime}(t) \neq 0$ for all $t$; then $\gamma^{\prime}(t)$ is its
\emph{velocity} or tangent vector at $t$ and $|\gamma^{\prime}(t)|$ its \emph{speed}.
\end{definition}

A \emph{reparametrisation} of $\gamma$ is $\gamma \circ \varphi$, where $\varphi \colon [c, d] \to [a, b]$ is a continuous bijection. Such a $\varphi$ is strictly monotone; it \emph{preserves the orientation} if
it is increasing, so that $\gamma \circ \varphi$ runs from $\gamma(a)$ to $\gamma(b)$, and \emph{reverses} it if it is decreasing. The \emph{reverse} of $\gamma$ is $\gamma^-(t) = \gamma(a + b - t)$. For
$C^1$ curves we use $C^1$ reparametrisations with $\varphi^{\prime}$ nowhere zero; by the chain rule $(\gamma \circ \varphi)^{\prime}(s) = \varphi^{\prime}(s)\,\gamma^{\prime}(\varphi(s))$, so the velocity is
rescaled but keeps its direction (up to sign). For instance $\gamma(t) = (\cos t, \sin t)$, $t \in [0, 2\pi]$, and $\tilde\gamma(s) = (\cos 2\pi s^2, \sin 2\pi s^2)$, $s \in [0, 1]$, run once around the unit circle at
different speeds, and $\gamma^-$ runs around it clockwise.

How long is a curve? The ancient answer, used by Archimedes for the circle, is to approximate it by inscribed polygons. A polygon through points of the curve is shorter than the curve, because the
segment is the shortest path between two points, and refining the polygon can only make it longer, by the triangle inequality. So the length should be the supremum of the lengths of the inscribed
polygons.

\begin{definition}\label{real-analysis:definition-length}
For a curve $\gamma \colon [a, b] \to \RR^n$ and a partition $P \colon a = t_0 < t_1 < \cdots < t_k = b$, let $\ell(\gamma, P) = \sum_{i=1}^k|\gamma(t_i) - \gamma(t_{i-1})|$ be the length of the inscribed polygon. The
\emph{length} of $\gamma$ is
\[
L(\gamma) = \sup_P\ell(\gamma, P) \in [0, \infty],
\]
and $\gamma$ is \emph{rectifiable} if $L(\gamma) < \infty$.
\end{definition}

\begin{lemma}\label{real-analysis:lemma-length-basic}
Let $\gamma \colon [a, b] \to \RR^n$ be a curve.
\begin{enumerate}
\item If $P^{\prime}$ refines $P$, then $\ell(\gamma, P) \leq \ell(\gamma, P^{\prime})$, and $|\gamma(b) - \gamma(a)| \leq L(\gamma)$.
\item (\emph{Additivity}) For $a < c < b$, $L(\gamma) = L(\gamma|_{[a, c]}) + L(\gamma|_{[c, b]})$.
\item (\emph{Invariance}) $L(\gamma \circ \varphi) = L(\gamma)$ for every reparametrisation, of either orientation; in particular $L(\gamma^-) = L(\gamma)$.
\end{enumerate}
\end{lemma}

\begin{proof}
(1) Adding a point $s$ between $t_{i-1}$ and $t_i$ replaces $|\gamma(t_i) - \gamma(t_{i-1})|$ by $|\gamma(s) - \gamma(t_{i-1})| + |\gamma(t_i) - \gamma(s)|$, which is at least as large by the triangle inequality;
refining adds finitely many points. The trivial partition gives $|\gamma(b) - \gamma(a)| = \ell(\gamma, \{a, b\}) \leq L(\gamma)$.

(2) A partition of $[a, c]$ and one of $[c, b]$ together form a partition of $[a, b]$, so $L(\gamma|_{[a, c]}) + L(\gamma|_{[c, b]}) \leq L(\gamma)$, taking suprema over the two partitions separately.
Conversely, for a partition $P$ of $[a, b]$ let $P^{\prime} = P \cup \{c\}$; then by (1) $\ell(\gamma, P) \leq \ell(\gamma, P^{\prime})$, and $\ell(\gamma, P^{\prime})$ is the sum of the polygon lengths of the two restrictions,
which is at most $L(\gamma|_{[a, c]}) + L(\gamma|_{[c, b]})$.

(3) A strictly monotone bijection $\varphi \colon [c, d] \to [a, b]$ maps partitions of $[c, d]$ bijectively to partitions of $[a, b]$ (reversing their order if $\varphi$ decreases), and the inscribed polygons of
$\gamma \circ \varphi$ and of $\gamma$ for corresponding partitions have the same vertices, hence the same length.
\end{proof}

For a $C^1$ curve the supremum is an integral: the length is the integral of the speed, as the distance travelled by a car is the integral of the reading of its speedometer.

\begin{theorem}[Length of a $C^1$ curve]\label{real-analysis:theorem-length-integral}
Every piecewise $C^1$ curve is rectifiable, and
\[
L(\gamma) = \int_a^b|\gamma^{\prime}(t)|\,dt,
\]
the integral being the sum of the integrals over the pieces where $\gamma$ is $C^1$.
\end{theorem}

\begin{proof}
By the additivity of Lemma~\ref{real-analysis:lemma-length-basic}(2) and of the integral, it suffices to treat a $C^1$ curve. The function $|\gamma^{\prime}|$ is continuous, hence integrable. For any partition, the
fundamental theorem of calculus (Theorem~\ref{real-analysis:theorem-ftc}) and the norm inequality of Proposition~\ref{real-analysis:proposition-riemann}(3) give
\[
|\gamma(t_i) - \gamma(t_{i-1})| = \Bigl|\int_{t_{i-1}}^{t_i}\gamma^{\prime}(t)\,dt\Bigr| \leq \int_{t_{i-1}}^{t_i}|\gamma^{\prime}(t)|\,dt,
\]
and summing, $\ell(\gamma, P) \leq \int_a^b|\gamma^{\prime}|$. So $L(\gamma) \leq \int_a^b|\gamma^{\prime}|$.

For the converse inequality let $\varepsilon > 0$. Since $\gamma^{\prime}$ is uniformly continuous on $[a, b]$ (Corollary~\ref{general-topology:corollary-uniform-continuity}), there is $\delta > 0$ with
$|\gamma^{\prime}(s) - \gamma^{\prime}(t)| < \varepsilon$ whenever $|s - t| < \delta$. Take a partition with $t_i - t_{i-1} < \delta$ for all $i$, and write $\Delta_i = t_i - t_{i-1}$. On $[t_{i-1}, t_i]$,
\[
\gamma(t_i) - \gamma(t_{i-1}) = \int_{t_{i-1}}^{t_i}\gamma^{\prime}(t)\,dt = \Delta_i\,\gamma^{\prime}(t_i) + \int_{t_{i-1}}^{t_i}\bigl(\gamma^{\prime}(t) - \gamma^{\prime}(t_i)\bigr)\,dt,
\]
and the last integral has norm at most $\varepsilon\Delta_i$. Hence $|\gamma(t_i) - \gamma(t_{i-1})| \geq \Delta_i|\gamma^{\prime}(t_i)| - \varepsilon\Delta_i$. On the other hand $|\gamma^{\prime}(t)| \leq |\gamma^{\prime}(t_i)| + \varepsilon$
on the interval, so $\int_{t_{i-1}}^{t_i}|\gamma^{\prime}| \leq \Delta_i|\gamma^{\prime}(t_i)| + \varepsilon\Delta_i$. Combining and summing over $i$,
\[
L(\gamma) \geq \ell(\gamma, P) \geq \int_a^b|\gamma^{\prime}(t)|\,dt - 2\varepsilon(b - a).
\]
As $\varepsilon$ is arbitrary, $L(\gamma) \geq \int_a^b|\gamma^{\prime}|$.
\end{proof}

In coordinates, $L(\gamma) = \int_a^b\sqrt{\gamma_1^{\prime}(t)^2 + \cdots + \gamma_n^{\prime}(t)^2}\,dt$. That this integral does not depend on the parametrisation, which is clear for the definition by polygons, can
also be seen by the substitution rule: for $\varphi$ increasing, $\int_c^d|\gamma^{\prime}(\varphi(s))|\,\varphi^{\prime}(s)\,ds = \int_a^b|\gamma^{\prime}(t)|\,dt$ (Corollary~\ref{real-analysis:corollary-integration-rules}), and for
$\varphi$ decreasing the absolute value $|\varphi^{\prime}| = -\varphi^{\prime}$ compensates the reversal of the limits.

\begin{example}\label{real-analysis:example-lengths}
\begin{enumerate}
\item \emph{Circles.} For $\gamma(t) = (r\cos t, r\sin t)$, $t \in [0, 2\pi]$, $\gamma^{\prime}(t) = (-r\sin t, r\cos t)$ has norm $r$, so $L(\gamma) = 2\pi r$. With the definition of $\pi$ in
Theorem~\ref{real-analysis:theorem-pi}, this is the theorem that the circumference of a circle is $2\pi$ times its radius.
\item \emph{Graphs.} The graph of a $C^1$ function $f \colon [a, b] \to \RR$, parametrised by $\gamma(x) = (x, f(x))$, has length $\int_a^b\sqrt{1 + f^{\prime}(x)^2}\,dx$. For a line $f(x) = mx$ this is
$(b - a)\sqrt{1 + m^2}$, the length of the segment, by Pythagoras.
\item \emph{The helix.} $\gamma(t) = (r\cos t, r\sin t, ct)$ has constant speed $\sqrt{r^2 + c^2}$, so one turn has length $2\pi\sqrt{r^2 + c^2}$: unrolling the cylinder of radius $r$ into a plane turns the helix
into the diagonal of a rectangle with sides $2\pi r$ and $2\pi c$.
\item \emph{The cycloid.} A point on a rolling wheel of radius $1$ traces $\gamma(t) = (t - \sin t, 1 - \cos t)$. Here $|\gamma^{\prime}(t)|^2 = (1 - \cos t)^2 + \sin^2t = 2 - 2\cos t = 4\sin^2(t/2)$, so one arch has length
$\int_0^{2\pi}2\sin(t/2)\,dt = 8$, four times the diameter of the wheel (Wren, 1658). The curve is not regular: the speed vanishes at $t = 0$ and $t = 2\pi$, where the cycloid has cusps.
\end{enumerate}
\end{example}

\begin{remark}\label{real-analysis:remark-elliptic-integrals}
The length of an ellipse $(a\cos t, b\sin t)$ with $0 < b < a$ is $\int_0^{2\pi}\sqrt{a^2\sin^2t + b^2\cos^2t}\,dt = 4a\int_0^{\pi/2}\sqrt{1 - e^2\cos^2t}\,dt$, with eccentricity $e = \sqrt{1 - b^2/a^2}$. Unlike
the circle, this integral cannot be expressed by elementary functions. It is an \emph{elliptic integral}, and the attempt to understand such integrals, by Euler, Legendre, Abel and Jacobi, led to elliptic
functions and to elliptic curves: inverting the integral gives doubly periodic functions, and the natural domain of the integrand is a curve of genus one (Chapter~\ref{curves}).
\end{remark}

Continuity alone does not give a finite length.

\begin{counterexample}\label{real-analysis:counterexample-non-rectifiable}
Let $f(x) = x\cos(\pi/x)$ for $0 < x \leq 1$ and $f(0) = 0$; it is continuous, since $|f(x)| \leq x$, and $C^1$ on $(0, 1]$. At $x_m = 1/m$, $f(x_m) = (-1)^m/m$. Take the partition of $[0, 1]$ by $0$ and
the points $x_m$, $m \leq N$; consecutive values have opposite signs, so the inscribed polygon of the graph $\gamma(x) = (x, f(x))$ has length at least
\[
\sum_{m=1}^{N-1}|f(x_m) - f(x_{m+1})| = \sum_{m=1}^{N-1}\Bigl(\frac1m + \frac{1}{m+1}\Bigr) \geq \sum_{m=1}^{N-1}\frac1m,
\]
which tends to $\infty$ with $N$ (Proposition~\ref{real-analysis:proposition-series-tests}(1)). So $\gamma$ is a continuous curve of infinite length: it oscillates too much near $0$. Worse examples exist: the Koch snowflake
curve has infinite length between any two of its points, and the space-filling curves of Chapter~\ref{general-topology} fill a square.
\end{counterexample}

For a regular curve the length can be used as a parameter. This removes the arbitrariness of the parametrisation and is the starting point of the differential geometry of curves.

\begin{proposition}[Arc length parametrisation]\label{real-analysis:proposition-arc-length}
Let $\gamma \colon [a, b] \to \RR^n$ be a regular $C^1$ curve of length $L$, and $\sigma(t) = \int_a^t|\gamma^{\prime}(u)|\,du$ the length of $\gamma|_{[a, t]}$. Then $\sigma \colon [a, b] \to [0, L]$ is a $C^1$ bijection with
$\sigma^{\prime} = |\gamma^{\prime}| > 0$, its inverse is $C^1$, and $\tilde\gamma = \gamma \circ \sigma^{-1} \colon [0, L] \to \RR^n$ has $|\tilde\gamma^{\prime}(s)| = 1$ for all $s$. It is the unique orientation-preserving
reparametrisation of $\gamma$ with unit speed.
\end{proposition}

\begin{proof}
By the fundamental theorem of calculus $\sigma$ is $C^1$ with $\sigma^{\prime}(t) = |\gamma^{\prime}(t)| > 0$, so it is strictly increasing, with $\sigma(a) = 0$ and $\sigma(b) = L$ by Theorem~\ref{real-analysis:theorem-length-integral}.
By Proposition~\ref{real-analysis:proposition-inverse-one} its inverse is differentiable with $(\sigma^{-1})^{\prime}(s) = 1/|\gamma^{\prime}(\sigma^{-1}(s))|$, which is continuous. By the chain rule
$|\tilde\gamma^{\prime}(s)| = |\gamma^{\prime}(\sigma^{-1}(s))|\,(\sigma^{-1})^{\prime}(s) = 1$. If $\gamma \circ \varphi$ is another orientation-preserving $C^1$ reparametrisation with unit speed, defined on $[c, d]$, then
$1 = |\gamma^{\prime}(\varphi(u))|\varphi^{\prime}(u) = (\sigma \circ \varphi)^{\prime}(u)$, so $\sigma(\varphi(u)) = u - c$, and $\gamma \circ \varphi$ is $\tilde\gamma$ up to the translation $u \mapsto u - c$ of the parameter.
\end{proof}

The parameter $s$ of $\tilde\gamma$ is the \emph{arc length}, and one writes $ds = |\gamma^{\prime}(t)|\,dt$ for the \emph{element of arc length}: it is the length of the piece of curve swept out between $t$ and
$t + dt$, to first order.

\subsection{Line integrals of the first kind}\label{real-analysis:subsection-first-kind}

Imagine a wire bent along a curve, with a linear density of mass $f$, in grams per centimetre. Its mass is obtained by cutting it into short pieces, multiplying the density at a point of each piece by the
length of the piece, and adding up; in the limit this is an integral with respect to arc length.

\begin{definition}\label{real-analysis:definition-first-kind}
Let $\gamma \colon [a, b] \to \RR^n$ be a piecewise $C^1$ curve and $f$ a continuous real function on its image $\gamma([a, b])$. The \emph{line integral of $f$ with respect to arc length}, or \emph{of the
first kind}, is
\[
\int_\gamma f\,ds = \int_a^bf(\gamma(t))\,|\gamma^{\prime}(t)|\,dt,
\]
the sum of the integrals over the pieces where $\gamma$ is $C^1$.
\end{definition}

\begin{proposition}\label{real-analysis:proposition-first-kind}
Let $\gamma$ and $f$ be as in the definition.
\begin{enumerate}
\item $\int_\gamma1\,ds = L(\gamma)$.
\item The integral is linear in $f$, monotone ($f \leq g$ implies $\int_\gamma f\,ds \leq \int_\gamma g\,ds$), and $\bigl|\int_\gamma f\,ds\bigr| \leq \sup|f|\cdot L(\gamma)$.
\item It is unchanged under every $C^1$ reparametrisation $\gamma \circ \varphi$ with $\varphi^{\prime}$ nowhere zero, of either orientation; in particular $\int_{\gamma^-}f\,ds = \int_\gamma f\,ds$.
\item It is additive: if $\gamma$ is the concatenation of $\gamma|_{[a, c]}$ and $\gamma|_{[c, b]}$, the integral is the sum of the two integrals.
\item (\emph{Riemann sums}) $\sum_if(\gamma(t_i))\,|\gamma(t_i) - \gamma(t_{i-1})| \to \int_\gamma f\,ds$ as the mesh $\max_i(t_i - t_{i-1})$ of the partition tends to $0$.
\end{enumerate}
\end{proposition}

\begin{proof}
(1) is Theorem~\ref{real-analysis:theorem-length-integral}, and (2) and (4) follow from the corresponding properties of the integral (Proposition~\ref{real-analysis:proposition-riemann}). (3) For each $C^1$ piece
and $\varphi$ increasing, the substitution rule (Corollary~\ref{real-analysis:corollary-integration-rules}) gives
$\int_c^df(\gamma(\varphi(u)))|\gamma^{\prime}(\varphi(u))|\varphi^{\prime}(u)\,du = \int_a^bf(\gamma(t))|\gamma^{\prime}(t)|\,dt$, and $|(\gamma \circ \varphi)^{\prime}| = |\gamma^{\prime} \circ \varphi|\,\varphi^{\prime}$. For $\varphi$ decreasing,
$|(\gamma \circ \varphi)^{\prime}| = -|\gamma^{\prime} \circ \varphi|\,\varphi^{\prime}$, and the substitution reverses the limits, so the two signs cancel.

(5) It suffices to treat a $C^1$ curve. The proof of Theorem~\ref{real-analysis:theorem-length-integral} shows that for partitions of mesh $< \delta$,
$\bigl||\gamma(t_i) - \gamma(t_{i-1})| - \int_{t_{i-1}}^{t_i}|\gamma^{\prime}|\bigr| \leq 2\varepsilon\Delta_i$. Also $f \circ \gamma$ is uniformly continuous, so for small mesh it varies by less than $\varepsilon$
on each interval. Hence the Riemann sum differs from $\sum_i\int_{t_{i-1}}^{t_i}f(\gamma(t))|\gamma^{\prime}(t)|\,dt = \int_\gamma f\,ds$ by at most $\sup|f| \cdot 2\varepsilon(b - a) + \varepsilon L(\gamma)$.
\end{proof}

By (3), the integral of the first kind depends only on the curve traced out, not on the direction or the speed of travel, as the mass of a wire should. The \emph{centre of mass} of a wire of density
$f > 0$ is the point with coordinates $\bar x_j = \int_\gamma x_jf\,ds / \int_\gamma f\,ds$.

\begin{example}\label{real-analysis:example-first-kind}
\begin{enumerate}
\item The upper half of the circle of radius $r$, $\gamma(t) = (r\cos t, r\sin t)$, $t \in [0, \pi]$, with constant density $1$, has mass $\pi r$ and
$\int_\gamma y\,ds = \int_0^\pi r\sin t \cdot r\,dt = 2r^2$. So its centre of mass is $(0, 2r/\pi)$, inside the half disc and not on the wire.
\item For the helix of Example~\ref{real-analysis:example-lengths}(3) over one turn, $\int_\gamma z\,ds = \int_0^{2\pi}ct\sqrt{r^2 + c^2}\,dt = 2\pi^2c\sqrt{r^2 + c^2}$, and the centre of mass is at height $\pi c$, half way up, as symmetry
suggests.
\end{enumerate}
\end{example}

\subsection{Line integrals of the second kind}\label{real-analysis:subsection-second-kind}

Now let a force field $F$ act on a particle moving along a curve. The work done by the force on a short displacement $\Delta x$ is $F \cdot \Delta x$: only the component of the force along the
displacement counts. Adding up along the curve gives an integral that depends on the direction of travel: going backwards, every displacement changes sign.

\begin{definition}\label{real-analysis:definition-second-kind}
Let $\gamma \colon [a, b] \to \RR^n$ be a piecewise $C^1$ curve and $F = (F_1, \ldots, F_n)$ a continuous vector field on its image. The \emph{line integral of $F$ along $\gamma$}, or \emph{of the second kind},
is
\[
\int_\gamma F \cdot d\gamma = \int_a^b\bigl\langle F(\gamma(t)), \gamma^{\prime}(t)\bigr\rangle\,dt = \int_a^b\sum_{j=1}^nF_j(\gamma(t))\,\gamma_j^{\prime}(t)\,dt,
\]
also written $\int_\gamma F_1\,dx_1 + \cdots + F_n\,dx_n$; in the plane, $\int_\gamma P\,dx + Q\,dy$ for $F = (P, Q)$.
\end{definition}

The notation $F_1\,dx_1 + \cdots + F_n\,dx_n$ is a differential form (Chapter~\ref{differential-forms}): the rule is to substitute $x_j = \gamma_j(t)$ and $dx_j = \gamma_j^{\prime}(t)\,dt$.

\begin{proposition}\label{real-analysis:proposition-second-kind}
Let $\gamma$ and $F$ be as in the definition.
\begin{enumerate}
\item The integral is linear in $F$ and additive under concatenation.
\item It is unchanged under orientation-preserving $C^1$ reparametrisations and changes sign under orientation-reversing ones: $\int_{\gamma^-}F \cdot d\gamma^- = -\int_\gamma F \cdot d\gamma$.
\item If $\gamma$ is regular, with unit tangent $T = \gamma^{\prime}/|\gamma^{\prime}|$, then $\int_\gamma F \cdot d\gamma = \int_\gamma\langle F, T\rangle\,ds$: the integral of the second kind is the integral of the first kind of
the tangential component of $F$.
\item $\bigl|\int_\gamma F \cdot d\gamma\bigr| \leq \sup|F| \cdot L(\gamma)$.
\item (\emph{Fundamental theorem for line integrals}) If $F = \nabla g$ is the gradient of a $C^1$ function $g$ on an open set containing the image of $\gamma$, then
$\int_\gamma\nabla g \cdot d\gamma = g(\gamma(b)) - g(\gamma(a))$. In particular the integral over a closed curve is $0$.
\end{enumerate}
\end{proposition}

\begin{proof}
(1) follows from the properties of the integral. (2) On a $C^1$ piece, with $(\gamma \circ \varphi)^{\prime}(u) = \varphi^{\prime}(u)\gamma^{\prime}(\varphi(u))$, the substitution rule gives
$\int_c^d\langle F(\gamma(\varphi(u))), \gamma^{\prime}(\varphi(u))\rangle\varphi^{\prime}(u)\,du = \int_{\varphi(c)}^{\varphi(d)}\langle F(\gamma(t)), \gamma^{\prime}(t)\rangle\,dt$, which is the original integral if $\varphi$
increases and its negative if $\varphi$ decreases. (3) $\langle F(\gamma(t)), \gamma^{\prime}(t)\rangle = \langle F(\gamma(t)), T(t)\rangle|\gamma^{\prime}(t)|$. (4) By Cauchy--Schwarz,
$|\langle F(\gamma(t)), \gamma^{\prime}(t)\rangle| \leq \sup|F|\,|\gamma^{\prime}(t)|$; integrate and use Theorem~\ref{real-analysis:theorem-length-integral}. (5) On a $C^1$ piece, by the chain rule
(Proposition~\ref{real-analysis:proposition-derivative-basic}), $(g \circ \gamma)^{\prime}(t) = \langle\nabla g(\gamma(t)), \gamma^{\prime}(t)\rangle$, so the integral is $\int(g \circ \gamma)^{\prime} = g(\gamma(c_k)) - g(\gamma(c_{k-1}))$ by
Theorem~\ref{real-analysis:theorem-ftc}; the intermediate values cancel in the sum.
\end{proof}

A vector field that is a gradient, $F = \nabla g$, is called \emph{conservative}, and $-g$ (in physics) its \emph{potential}: the work done along a path depends only on its endpoints. Not every field
is of this kind.

\begin{example}\label{real-analysis:example-second-kind}
\begin{enumerate}
\item \emph{Rotation.} For $F(x, y) = (-y, x)$ and the circle $\gamma(t) = (r\cos t, r\sin t)$, $t \in [0, 2\pi]$, $\langle F(\gamma(t)), \gamma^{\prime}(t)\rangle = r^2\sin^2t + r^2\cos^2t = r^2$, so
$\int_\gamma F \cdot d\gamma = 2\pi r^2$, twice the area enclosed; this is no accident (Green's theorem, Chapter~\ref{riemannian-geometry}). The integral over a closed curve is not $0$, so $F$ is not a gradient.
\item \emph{Dependence on the path.} For $F(x, y) = (y, 0)$, that is $\int y\,dx$, from $(0, 0)$ to $(1, 1)$: along the segment $\gamma_1(t) = (t, t)$ the integral is $\int_0^1t\,dt = \frac12$; along the
parabola $\gamma_2(t) = (t, t^2)$ it is $\int_0^1t^2\,dt = \frac13$; along the path that goes first to $(1, 0)$ and then up to $(1, 1)$ it is $0 + 0 = 0$, since $y = 0$ on the first piece and $dx = 0$ on
the second.
\item \emph{A gradient.} For $F = \nabla g$ with $g(x, y) = x^2y$, $F = (2xy, x^2)$, and every path from $(0, 0)$ to $(1, 1)$ gives $g(1, 1) - g(0, 0) = 1$. Along $\gamma_2$, for instance,
$\int_0^1(2t \cdot t^2 + t^2 \cdot 2t)\,dt = \int_0^14t^3\,dt = 1$.
\item \emph{The angle field.} $F(x, y) = (-y, x)/(x^2 + y^2)$ on $\RR^2 \setminus \{0\}$ has integral $2\pi$ around the unit circle, so it is not a gradient on $\RR^2 \setminus \{0\}$, although it is the gradient of
the polar angle on every half-plane through the origin. This field, the $1$-form $(x\,dy - y\,dx)/(x^2 + y^2)$, is the simplest example of a closed form that is not exact, and the beginning of de Rham cohomology
(Chapter~\ref{de-rham}).
\end{enumerate}
\end{example}

\begin{remark}\label{real-analysis:remark-two-kinds}
The two kinds of line integrals answer different questions and behave differently. The first kind integrates a \emph{density}, needs the Euclidean length to define $ds$, and ignores orientation; it
generalises to integrals over curves in a Riemannian manifold (Chapter~\ref{riemannian-geometry}) and, in higher dimensions, to surface integrals $\int_\Sigma f\,dA$. The second kind integrates a
\emph{$1$-form}, needs no metric at all, since $F_1\,dx_1 + \cdots + F_n\,dx_n$ is evaluated on velocity vectors, and depends on the orientation; it generalises to the integrals of $k$-forms over oriented
$k$-dimensional submanifolds, and Stokes' theorem lives there (Chapter~\ref{differential-forms}). Complex line integrals $\int_\gamma f(z)\,dz = \int_a^bf(\gamma(t))\gamma^{\prime}(t)\,dt$, defined in Chapter~\ref{complex-analysis} by this formula, the complex-valued integrand being integrated
separately in its real and imaginary parts, are of the second kind, with complex coefficients: expanding $f\,dz = (u + iv)(dx + i\,dy)$ gives
$\int_\gamma f\,dz = \int_\gamma(u\,dx - v\,dy) + i\int_\gamma(v\,dx + u\,dy)$, two real line integrals of the second kind. The relation $\int_\gamma F \cdot d\gamma = \int_\gamma\langle F, T\rangle\,ds$ shows that the second kind is a special case of the first, once a
direction has been chosen; conversely $\int_\gamma f\,ds = \int_\gamma fT \cdot d\gamma$ for a regular curve. In Italian the two are the \emph{integrali curvilinei di prima e di seconda specie}.
\end{remark}

\section{The inverse and implicit function theorems}\label{real-analysis:section-inverse-function}

\begin{lemma}\label{real-analysis:lemma-neumann}
If $A \in \operatorname{End}(\RR^n)$ has $\|A\| < 1$, then $1 - A$ is invertible with inverse $\sum_kA^k$ and $\|(1 - A)^{-1}\| \leq 1/(1 - \|A\|)$. The group $\GL_n(\RR)$ is open in $M_n(\RR)$, and
inversion $\GL_n(\RR) \to \GL_n(\RR)$ is smooth.
\end{lemma}

\begin{proof}
The series converges absolutely in the complete space $M_n(\RR)$ (Proposition~\ref{real-analysis:proposition-series}), and $(1 - A)\sum_{k \leq N}A^k = 1 - A^{N+1} \to 1$. The determinant is a polynomial in the
entries, hence continuous, and $\GL_n(\RR) = \det^{-1}(\RR \setminus \{0\})$. By the adjugate formula (Proposition~\ref{linear-algebra:proposition-determinant-properties}), the entries of $A^{-1}$ are
polynomials in the entries of $A$ divided by $\det A$, so inversion is smooth.
\end{proof}

\begin{theorem}[Inverse function theorem]\label{real-analysis:theorem-inverse-function}
Let $U \subseteq \RR^n$ be open, $f \colon U \to \RR^n$ of class $C^k$ with $k \geq 1$, and $a \in U$ with $Df(a)$ invertible. There are open sets $V \ni a$ and $W \ni f(a)$ such that $f \colon V \to W$ is a
$C^k$-diffeomorphism, and $D(f^{-1})(f(x)) = Df(x)^{-1}$.
\end{theorem}

\begin{proof}
Replacing $f$ by $x \mapsto Df(a)^{-1}(f(a + x) - f(a))$, we may assume $a = 0$, $f(0) = 0$ and $Df(0) = 1$. Let $g(x) = x - f(x)$, so $Dg(0) = 0$; choose $r > 0$ with $\overline{B}(0, r) \subseteq U$ and
$\|Dg(x)\| \leq 1/2$ there. By the mean value inequality, $\|g(x) - g(x')\| \leq \frac12\|x - x'\|$ on $\overline{B}(0, r)$, and $\|g(x)\| \leq r/2$.

For $y \in W = B(0, r/2)$, the map $T_y(x) = y + g(x)$ sends $\overline{B}(0, r)$ into $B(0, r)$ and is a contraction with constant $1/2$. By the Banach fixed point theorem
(Theorem~\ref{general-topology:theorem-banach-fixed-point}) there is a unique $x \in \overline{B}(0, r)$ with $x = y + g(x)$, that is, $f(x) = y$, and $x \in B(0, r)$. Let $V = B(0, r) \cap f^{-1}(W)$, open; then
$f \colon V \to W$ is bijective. From $x - x' = f(x) - f(x') + g(x) - g(x')$ we get $\|x - x'\| \leq 2\|f(x) - f(x')\|$, so $f^{-1}$ is continuous. For $x \in V$, $Df(x) = 1 - Dg(x)$ is invertible
with $\|Df(x)^{-1}\| \leq 2$ (Lemma~\ref{real-analysis:lemma-neumann}).

Let $y = f(x)$ and $y + \kappa = f(x + h)$ with $x, x + h \in V$, so $\|h\| \leq 2\|\kappa\|$. Then $f^{-1}(y + \kappa) - f^{-1}(y) - Df(x)^{-1}\kappa = -Df(x)^{-1}(f(x + h) - f(x) - Df(x)h)$, whose norm is
at most $2\|r(h)\| = o(\|\kappa\|)$. So $f^{-1}$ is differentiable with $D(f^{-1}) = \iota \circ Df \circ f^{-1}$, where $\iota$ is inversion. If $f^{-1}$ is $C^j$ with $j < k$, this composite is $C^j$, so
$f^{-1}$ is $C^{j+1}$; by induction $f^{-1}$ is $C^k$ (composites of $C^j$ maps are $C^j$ by the chain rule and induction).
\end{proof}

\begin{theorem}[Implicit function theorem]\label{real-analysis:theorem-implicit-function}
Let $U \subseteq \RR^n \times \RR^m$ be open, $F \colon U \to \RR^m$ of class $C^k$, $(a, b) \in U$ with $F(a, b) = 0$, and suppose the partial derivative $D_yF(a, b) \colon \RR^m \to \RR^m$ is invertible.
There are open sets $A \ni a$ and $B \ni b$ and a $C^k$ map $g \colon A \to B$ such that for $(x, y) \in A \times B$: $F(x, y) = 0$ if and only if $y = g(x)$.
\end{theorem}

\begin{proof}
The map $\Phi(x, y) = (x, F(x, y))$ has invertible derivative at $(a, b)$, so by Theorem~\ref{real-analysis:theorem-inverse-function} it restricts to a $C^k$-diffeomorphism of an open set, which we may
shrink to $A' \times B$, onto an open set containing $(a, 0)$. Its inverse has the form $(x, z) \mapsto (x, \psi(x, z))$. Let $A = \{x : (x, 0) \in \Phi(A' \times B)\}$, open, and $g(x) = \psi(x, 0)$.
\end{proof}

\begin{theorem}[Local submersions and immersions]\label{real-analysis:theorem-submersion-immersion}
Let $f \colon U \to \RR^m$ be $C^k$, $U \subseteq \RR^n$ open, $a \in U$.
\begin{enumerate}
\item If $Df(a)$ is surjective, there is a $C^k$-diffeomorphism $\psi$ from a neighbourhood of $a$ onto an open subset of $\RR^n$ with $f \circ \psi^{-1}(x_1, \ldots, x_n) = (x_1, \ldots, x_m)$.
\item If $Df(a)$ is injective, there is a $C^k$-diffeomorphism $\phi$ from a neighbourhood of $f(a)$ onto an open subset of $\RR^m$ with $\phi \circ f(x) = (x_1, \ldots, x_n, 0, \ldots, 0)$ near $a$.
\end{enumerate}
\end{theorem}

\begin{proof}
We may assume $a = 0$ and $f(0) = 0$. (1) After permuting coordinates, the first $m$ columns of the Jacobian are linearly independent. Then $\psi(x) = (f(x), x_{m+1}, \ldots, x_n)$ has invertible
derivative at $0$, and $f \circ \psi^{-1}$ is the projection. (2) After permuting target coordinates, the first $n$ rows of the Jacobian are independent. Let $\Theta(x, y) = f(x) + (0, y)$ for $x \in U$,
$y \in \RR^{m-n}$; $D\Theta(0)$ is invertible, and $\phi = \Theta^{-1}$ satisfies $\phi(f(x)) = (x, 0)$.
\end{proof}

\begin{remark}\label{real-analysis:remark-constant-rank}
More generally, if $Df$ has constant rank $r$ near $a$, there are local diffeomorphisms in source and target in which $f$ becomes $(x_1, \ldots, x_n) \mapsto (x_1, \ldots, x_r, 0, \ldots, 0)$ (the
\emph{constant rank theorem} \cite[Theorem 4.12]{lee-smooth-manifolds}). All results of this section hold with the same proofs over $\CC$ for complex differentiable (holomorphic) maps, once
holomorphic functions are known to be $C^\infty$ (Chapter~\ref{several-complex-variables}).
\end{remark}

\section{Extrema in several variables, Lagrange multipliers and the Karush--Kuhn--Tucker conditions}\label{real-analysis:section-extrema}

In one variable, Fermat's lemma and the second derivative test (Lemma~\ref{real-analysis:lemma-fermat}, Corollary~\ref{real-analysis:corollary-second-derivative-test}) locate the extrema of a function.
This section does the same in several variables, first for functions on open sets, then for functions restricted to a set defined by equations (Lagrange multipliers), and finally to a set
defined by equations and inequalities (the Karush--Kuhn--Tucker conditions), which is the basic situation of optimisation theory. The tools are the chain rule, Taylor's formula and the implicit
function theorem of the previous sections, together with one piece of convex geometry, Farkas' lemma.

\subsection{Gradient and Hessian}\label{real-analysis:subsection-gradient-hessian}

Let $U \subseteq \RR^n$ be open and $f \colon U \to \RR$ differentiable at $a$. The \emph{gradient} of $f$ at $a$ is the vector $\nabla f(a) = (\partial_1f(a), \ldots, \partial_nf(a)) \in \RR^n$; by
Proposition~\ref{real-analysis:proposition-derivative-basic}(1) it represents the derivative through the standard inner product,
\[
Df(a)h = \langle\nabla f(a), h\rangle \qquad (h \in \RR^n).
\]
If $\nabla f(a) \neq 0$, the Cauchy--Schwarz inequality shows that among unit vectors $h$ the \emph{directional derivative} $\langle\nabla f(a), h\rangle$ is largest for $h = \nabla f(a)/\|\nabla f(a)\|$: the
gradient points in the direction of steepest increase, and its length is the rate of increase in that direction. It is orthogonal to the level set through $a$, in the sense that
$\langle\nabla f(a), \gamma'(0)\rangle = (f \circ \gamma)'(0) = 0$ for every differentiable curve $\gamma$ with $\gamma(0) = a$ on which $f$ is constant.

If $f$ is $C^2$, its \emph{Hessian} at $a$ is the matrix $D^2f(a) = (\partial_i\partial_jf(a))_{i,j}$, which is symmetric by Schwarz's theorem (Theorem~\ref{real-analysis:theorem-schwarz}). We write
$D^2f(a)(h, h) = h^TD^2f(a)h = \sum_{i,j}\partial_i\partial_jf(a)h_ih_j$ for the associated quadratic form. With this notation Taylor's theorem (Theorem~\ref{real-analysis:theorem-taylor-several}) for $k = 1$ reads
\begin{equation}\label{real-analysis:equation-taylor-hessian}
f(a + h) = f(a) + \langle\nabla f(a), h\rangle + \tfrac12D^2f(\xi)(h, h) \qquad \text{for some } \xi \in (a, a + h),
\end{equation}
since $\sum_{|\alpha| = 2}\frac{\partial^\alpha f}{\alpha!}h^\alpha = \frac12\sum_{i,j}\partial_i\partial_jfh_ih_j$ (a diagonal term $\alpha = 2e_i$ has $\alpha! = 2$ and occurs once, an off-diagonal pair $\{i, j\}$ has
$\alpha! = 1$ and occurs twice in the double sum). Along a line, if $g(t) = f(a + th)$, the chain rule gives
\begin{equation}\label{real-analysis:equation-line-derivatives}
g'(t) = \langle\nabla f(a + th), h\rangle, \qquad g''(t) = D^2f(a + th)(h, h).
\end{equation}

A real symmetric matrix $H$, or its quadratic form, is \emph{positive definite} if $h^THh > 0$ for all $h \neq 0$, \emph{positive semidefinite} if $h^THh \geq 0$ for all $h$, \emph{negative (semi)definite} if
$-H$ is positive (semi)definite, and \emph{indefinite} if $h^THh$ takes both signs. By the spectral theorem (Theorem~\ref{linear-algebra:theorem-spectral}), $H = QDQ^T$ with $Q$ orthogonal, and
$h^THh = \sum_i\lambda_iy_i^2$ with $y = Q^Th$: so $H$ is positive definite if and only if all its eigenvalues are positive, positive semidefinite if and only if they are all $\geq 0$, and indefinite if
and only if it has eigenvalues of both signs.

\begin{definition}\label{real-analysis:definition-local-extremum-several}
Let $S \subseteq \RR^n$ and $f \colon S \to \RR$. Then $f$ has a \emph{local minimum} at $a \in S$ if there is $\delta > 0$ with $f(x) \geq f(a)$ for all $x \in S$ with $\|x - a\| < \delta$, and a \emph{strict local
minimum} if moreover $f(x) > f(a)$ for these $x \neq a$; similarly for (strict) local maxima. A \emph{global minimum} on $S$ is a point $a$ with $f(x) \geq f(a)$ for all $x \in S$. When $S$ is given by
equations or inequalities and $f$ is the restriction of a function defined on a larger open set, one speaks of a \emph{constrained} extremum. A point $a$ of an open set $U$ where $f$ is differentiable
is a \emph{critical point} of $f$ if $\nabla f(a) = 0$.
\end{definition}

\begin{proposition}\label{real-analysis:proposition-hessian-test}
Let $U \subseteq \RR^n$ be open, $f \colon U \to \RR$ of class $C^2$, and $a \in U$.
\begin{enumerate}
\item If $f$ has a local minimum at $a$, then $\nabla f(a) = 0$ and $D^2f(a)$ is positive semidefinite.
\item If $\nabla f(a) = 0$ and $D^2f(a)$ is positive definite, then $f$ has a strict local minimum at $a$.
\item If $\nabla f(a) = 0$ and $D^2f(a)$ is indefinite, then $f$ has no local extremum at $a$ (a \emph{saddle point}).
\end{enumerate}
The same holds for maxima with ``negative'' in place of ``positive''.
\end{proposition}

\begin{proof}
(1) Let $h \in \RR^n$ and $g(t) = f(a + th)$, defined and $C^2$ on an interval around $0$, with a local minimum at $0$. By Lemma~\ref{real-analysis:lemma-fermat} and
\eqref{real-analysis:equation-line-derivatives}, $0 = g'(0) = \langle\nabla f(a), h\rangle$; as $h$ is arbitrary, $\nabla f(a) = 0$. If $g''(0) = D^2f(a)(h, h)$ were negative, $g$ would have a strict local
maximum at $0$ by Corollary~\ref{real-analysis:corollary-second-derivative-test}, so $g(t) < g(0) \leq g(t)$ for small $t \neq 0$, which is absurd. Hence $D^2f(a)(h, h) \geq 0$.

(2) The unit sphere $\{\|u\| = 1\}$ is compact (Theorem~\ref{general-topology:theorem-heine-borel}) and $u \mapsto D^2f(a)(u, u)$ is continuous and positive on it, so it has a minimum $m > 0$
(Corollary~\ref{general-topology:corollary-extreme-value}). For $\|u\| = 1$,
$|D^2f(\xi)(u, u) - D^2f(a)(u, u)| \leq \sum_{i,j}|\partial_i\partial_jf(\xi) - \partial_i\partial_jf(a)|$, which is $< m/2$ for $\|\xi - a\| < \delta$ with $\delta$ small, by continuity of the second
derivatives. Then for $0 < \|h\| < \delta$ (and $\overline{B}(a, \delta) \subseteq U$), \eqref{real-analysis:equation-taylor-hessian} with $\nabla f(a) = 0$ gives
\[
f(a + h) - f(a) = \tfrac12\|h\|^2D^2f(\xi)\bigl(\tfrac{h}{\|h\|}, \tfrac{h}{\|h\|}\bigr) \geq \tfrac{m}{4}\|h\|^2 > 0.
\]

(3) By (1), applied to $f$ and to $-f$, a local minimum forces $D^2f(a)$ positive semidefinite and a local maximum forces it negative semidefinite; an indefinite form is neither.
\end{proof}

\begin{example}\label{real-analysis:example-hessian-two-variables}
For $n = 2$ write $D^2f(a) = \bigl(\begin{smallmatrix} A & B\\ B & C\end{smallmatrix}\bigr)$. Completing the square, $Ah_1^2 + 2Bh_1h_2 + Ch_2^2 = A(h_1 + \frac BAh_2)^2 + \frac{AC - B^2}{A}h_2^2$ when
$A \neq 0$. So the Hessian is positive definite if and only if $A > 0$ and $AC - B^2 > 0$, negative definite if and only if $A < 0$ and $AC - B^2 > 0$, and indefinite if $AC - B^2 < 0$ (the
determinant is the product of the two eigenvalues). For instance:
\begin{enumerate}
\item $f(x, y) = x^2 + xy + y^2 - 3x$ has the single critical point $(2, -1)$, with $A = C = 2$, $B = 1$, $AC - B^2 = 3$: a strict local minimum, in fact global, as $f$ minus its value there is a positive
definite quadratic form in $(x - 2, y + 1)$.
\item $f(x, y) = x^3 - 3x + y^2$ has critical points $(\pm1, 0)$, with $D^2f = \operatorname{diag}(6x, 2)$: a strict local minimum at $(1, 0)$ and a saddle point at $(-1, 0)$. The local minimum is not
global, since $f(x, 0) \to -\infty$ as $x \to -\infty$.
\item $x^2 + y^4$, $x^2 - y^4$ and $x^2 + y^3$ all have the critical point $0$ with the semidefinite Hessian $\operatorname{diag}(2, 0)$, and they have a minimum, a saddle and a saddle there:
when the Hessian is degenerate the test says nothing, as in Counterexample~\ref{real-analysis:counterexample-second-derivative}.
\end{enumerate}
\end{example}

\begin{counterexample}\label{real-analysis:counterexample-lines}
Checking all lines through $a$ is not enough. Let $f(x, y) = (y - x^2)(y - 3x^2)$. On every line through $0$ the restriction of $f$ has a strict local minimum at $0$: on the $x$-axis $f = 3x^4$, and on
$y = mx$ we get $f = m^2x^2 - 4mx^3 + 3x^4$, whose second derivative at $0$ is $2m^2 > 0$ for $m \neq 0$ (and on the $y$-axis $f = y^2$). But $f < 0$ between the two parabolas, at $(x, 2x^2)$, arbitrarily
close to $0$; so $0$ is not a local minimum. Here $D^2f(0) = \operatorname{diag}(0, 2)$ is only semidefinite, as Proposition~\ref{real-analysis:proposition-hessian-test}(2) requires.
\end{counterexample}

\subsection{Lagrange multipliers}\label{real-analysis:subsection-lagrange}

Suppose we want the extrema of $f(x, y)$ on the curve $M = \{h(x, y) = 0\}$. Walking along $M$, the value of $f$ stops increasing at a point $a$ where $M$ does not cross the level curve of $f$
through $a$ but touches it. There the two curves have the same tangent line, so their normals $\nabla f(a)$ and $\nabla h(a)$ are parallel: $\nabla f(a) = \lambda\nabla h(a)$ for some number $\lambda$, the
\emph{Lagrange multiplier}. In this way a constrained problem in $n$ unknowns becomes an unconstrained system of $n + 1$ equations ($\nabla f = \lambda\nabla h$ and $h = 0$) in $n + 1$ unknowns. To prove
this we need curves inside the constraint set with a prescribed velocity; they are provided by the submersion theorem.

\begin{lemma}\label{real-analysis:lemma-curves-level-set}
Let $U \subseteq \RR^n$ be open, $h = (h_1, \ldots, h_q) \colon U \to \RR^q$ of class $C^k$ with $k \geq 1$, and $a \in U$ with $h(a) = 0$ and $Dh(a)$ surjective, that is, with $\nabla h_1(a), \ldots,
\nabla h_q(a)$ linearly independent. For every $d \in \ker Dh(a)$ there are $\varepsilon > 0$ and a $C^k$ curve $x \colon (-\varepsilon, \varepsilon) \to U$ with $x(0) = a$, $x'(0) = d$ and $h(x(t)) = 0$ for
all $t$.
\end{lemma}

\begin{proof}
By Theorem~\ref{real-analysis:theorem-submersion-immersion}(1) there is a $C^k$-diffeomorphism $\psi$ from an open $V \ni a$ onto an open $W \subseteq \RR^n$ with $h \circ \psi^{-1} = \pi$, where
$\pi(z) = (z_1, \ldots, z_q)$. (If $q = 0$ there is nothing to prove: take $x(t) = a + td$.) Thus $h = \pi \circ \psi$ on $V$, so $\pi(\psi(a)) = h(a) = 0$ and, by the chain rule, $\pi(D\psi(a)d) =Dh(a)d = 0$.
Put $v = D\psi(a)d$ and $x(t) = \psi^{-1}(\psi(a) + tv)$, defined for $|t| < \varepsilon$ since $W$ is open. Then $h(x(t)) = \pi(\psi(a)) + t\pi(v) = 0$, and by the chain rule and
Theorem~\ref{real-analysis:theorem-inverse-function}, $x'(0) = D\psi(a)^{-1}v = d$.
\end{proof}

\begin{theorem}[Lagrange multipliers]\label{real-analysis:theorem-lagrange}
Let $U \subseteq \RR^n$ be open, $f \colon U \to \RR$ and $h = (h_1, \ldots, h_q) \colon U \to \RR^q$ of class $C^1$, and $M = \{x \in U : h(x) = 0\}$. Let $a \in M$ be a point where
$\nabla h_1(a), \ldots, \nabla h_q(a)$ are linearly independent. If $f|_M$ has a local extremum at $a$, there are unique numbers $\lambda_1, \ldots, \lambda_q$ with
\begin{equation}\label{real-analysis:equation-lagrange}
\nabla f(a) = \lambda_1\nabla h_1(a) + \cdots + \lambda_q\nabla h_q(a).
\end{equation}
Equivalently, $(a, \lambda)$ is a critical point of the \emph{Lagrangian} $L(x, \lambda) = f(x) - \sum_j\lambda_jh_j(x)$ on $U \times \RR^q$.
\end{theorem}

\begin{proof}
Let $d \in \ker Dh(a)$, and let $x(t)$ be the curve of Lemma~\ref{real-analysis:lemma-curves-level-set}. Since $x(t) \in M$ and $x(t) \to a$, the function $f \circ x$ has a local extremum at $t = 0$, so
by Lemma~\ref{real-analysis:lemma-fermat} and the chain rule $0 = (f \circ x)'(0) = \langle\nabla f(a), d\rangle$. Thus $\nabla f(a)$ is orthogonal to $\ker Dh(a)$. Now $\ker Dh(a)$ is exactly the orthogonal
complement of the span $R$ of the rows $\nabla h_j(a)$ of the Jacobian matrix, so $\nabla f(a) \in (R^\perp)^\perp = R$ (in a finite-dimensional inner product space $(R^\perp)^\perp = R$, since
$R \subseteq (R^\perp)^\perp$ and both have dimension $\dim R$). Uniqueness of the $\lambda_j$ is the linear independence of the $\nabla h_j(a)$. Finally, $\partial L/\partial x_i = 0$ for all $i$ is
\eqref{real-analysis:equation-lagrange}, and $\partial L/\partial\lambda_j = -h_j = 0$ is the constraint.
\end{proof}

The space $\ker Dh(a)$ is the \emph{tangent space} of $M$ at $a$: by the lemma it is the set of velocities of curves in $M$ through $a$, and by the submersion theorem $M$ is, near $a$, a submanifold of
dimension $n - q$ in the sense of Chapter~\ref{smooth-manifolds}. The theorem says that at a constrained extremum the gradient of $f$ is normal to $M$. The linear independence hypothesis
cannot be dropped:

\begin{counterexample}\label{real-analysis:counterexample-lagrange-cusp}
Minimise $f(x, y) = x$ on the cusp $M = \{h = 0\}$, $h(x, y) = y^2 - x^3$. On $M$, $x^3 = y^2 \geq 0$, so $x \geq 0$ and the minimum is at $(0, 0)$. But $\nabla f = (1, 0)$, while
$\nabla h = (-3x^2, 2y)$ vanishes at the origin, so no $\lambda$ satisfies $\nabla f(0) = \lambda\nabla h(0)$. Solving the Lagrange system $1 = -3\lambda x^2$, $0 = 2\lambda y$, $y^2 = x^3$ gives no solution at all:
the method finds nothing, and the minimum sits at the singular point of $M$. In practice one must examine separately the points of $M$ where the gradients of the constraints are dependent.
\end{counterexample}

\begin{example}\label{real-analysis:example-lagrange}
\begin{enumerate}
\item \emph{Extrema of $x + y$ on the circle $x^2 + y^2 = 1$.} The circle is compact, so both extrema exist (Corollary~\ref{general-topology:corollary-extreme-value}), and $\nabla h = (2x, 2y) \neq 0$ on it.
The system $(1, 1) = \lambda(2x, 2y)$ gives $x = y = \frac1{2\lambda}$, and $x^2 + y^2 = 1$ gives $(x, y) = \pm(\frac1{\sqrt2}, \frac1{\sqrt2})$: maximum $\sqrt2$, minimum $-\sqrt2$.
\item \emph{The inequality of the means.} Maximise $f(x) = x_1x_2\cdots x_k$ on the compact set $\Delta = \{x : x_i \geq 0, \sum_ix_i = k\}$. The maximum is positive (take all $x_i = 1$), so it is attained at a
point with all $x_i > 0$, which is a local maximum of $f$ on $M = \{x \in U : \sum_ix_i = k\}$ with $U = \{x_i > 0\}$. With $h = \sum_ix_i - k$, $\nabla h = (1, \ldots, 1)$ and
$\partial_if = f(x)/x_i$, so $f(x)/x_i = \lambda$ for all $i$: all $x_i$ are equal, hence equal to $1$, and $f \leq 1$ on $\Delta$. By homogeneity, $\sqrt[k]{x_1\cdots x_k} \leq \frac1k\sum_ix_i$ for $x_i \geq 0$,
which recovers Corollary~\ref{real-analysis:corollary-am-gm} by a quite different route.
\item \emph{Eigenvectors as critical points.} Let $A$ be a real symmetric $n \times n$ matrix and $q(x) = x^TAx$, so that $\nabla q(x) = 2Ax$. On the compact unit sphere, $h(x) = \|x\|^2 - 1$ with
$\nabla h(x) = 2x \neq 0$, the function $q$ has a maximum at some $u$, and \eqref{real-analysis:equation-lagrange} reads $Au = \lambda u$: \emph{$u$ is an eigenvector}, and $\lambda = u^TAu$ is the maximum
of $q$ on the sphere, the largest eigenvalue. This proves, without the fundamental theorem of algebra, that a real symmetric matrix has a real eigenvector, which is the key step of
Theorem~\ref{linear-algebra:theorem-spectral}(2). More generally the critical points of the \emph{Rayleigh quotient} $x^TAx/x^Tx$ on $\RR^n \setminus \{0\}$ are exactly the eigenvectors.
\item \emph{Nearest points.} Let $h \colon \RR^n \to \RR$ be $C^1$, $M = \{h = 0\}$ nonempty with $\nabla h \neq 0$ on $M$, and $p \in \RR^n$. Since $M$ is closed, $f(x) = \|x - p\|^2$ attains a minimum on
$M$ (on the compact set $M \cap \overline{B}(p, R)$ for $R$ large). At the minimum $a$, $2(a - p) = \lambda\nabla h(a)$: \emph{the segment from $p$ to the nearest point is normal to $M$}. For the hyperplane
$\langle n, x\rangle = c$ this gives $a = p + tn$ with $t = (c - \langle n, p\rangle)/\|n\|^2$, and the familiar distance $|c - \langle n, p\rangle|/\|n\|$.
\end{enumerate}
\end{example}

The multiplier has a meaning of its own: it is the rate at which the optimal value changes when the constraint is relaxed.

\begin{proposition}\label{real-analysis:proposition-multiplier-sensitivity}
Let $f, h \colon U \to \RR$ be $C^1$ on an open $U \subseteq \RR^n$, $J$ an open interval, and $c \mapsto x(c)$ a differentiable map $J \to U$ such that for every $c \in J$, $h(x(c)) = c$ and
$\nabla f(x(c)) = \lambda(c)\nabla h(x(c))$ for some number $\lambda(c)$. Then $V(c) = f(x(c))$ is differentiable on $J$ and $V'(c) = \lambda(c)$.
\end{proposition}

\begin{proof}
By the chain rule, $V'(c) = \langle\nabla f(x(c)), x'(c)\rangle = \lambda(c)\langle\nabla h(x(c)), x'(c)\rangle = \lambda(c)\frac{d}{dc}h(x(c)) = \lambda(c)$.
\end{proof}

For example, the largest product $xy$ with $x + y = c$ is $V(c) = c^2/4$, at $x = y = c/2$, where $(y, x) = \frac c2(1, 1)$: indeed $\lambda = c/2 = V'(c)$. In economics, where $f$ is a profit and $h \leq c$ a
budget, $\lambda$ is the \emph{shadow price} of the resource: the extra profit that one more unit of it would bring. In mechanics, where Lagrange introduced the method in his \emph{M\'ecanique
analytique} (1788), $\lambda\nabla h$ is the force by which the constraint keeps the system on $M$.

\subsection{Second-order conditions}\label{real-analysis:subsection-second-order}

A critical point of the Lagrangian need not be an extremum, and as in Proposition~\ref{real-analysis:proposition-hessian-test} one looks at second derivatives. The correct quadratic form is
not the Hessian of $f$ but that of the Lagrangian, and only in tangent directions: the curvature of $M$ contributes through the multipliers.

\begin{theorem}\label{real-analysis:theorem-lagrange-second-order}
Let $f$ and $h = (h_1, \ldots, h_q)$ be $C^2$ on the open set $U \subseteq \RR^n$, $M = \{h = 0\}$, and $a \in M$ with $\nabla f(a) = \sum_j\lambda_j\nabla h_j(a)$. Put
$L = f - \sum_j\lambda_jh_j$ (with these fixed $\lambda_j$) and $T = \ker Dh(a)$.
\begin{enumerate}
\item If $\nabla h_1(a), \ldots, \nabla h_q(a)$ are linearly independent and $f|_M$ has a local minimum at $a$, then $D^2L(a)(d, d) \geq 0$ for all $d \in T$.
\item If $D^2L(a)(d, d) > 0$ for all nonzero $d \in T$, then $f|_M$ has a strict local minimum at $a$.
\end{enumerate}
\end{theorem}

\begin{proof}
Note that $\nabla L(a) = 0$ and that $L = f$ on $M$.

(1) Let $d \in T$ and let $x(t)$ be the $C^2$ curve of Lemma~\ref{real-analysis:lemma-curves-level-set}, and $\varphi(t) = f(x(t)) = L(x(t))$. Differentiating
$\varphi^{\prime}(t) = \sum_i\partial_iL(x(t))x_i^{\prime}(t)$ once more,
\[
\varphi^{\prime\prime}(0) = \sum_{i,j}\partial_j\partial_iL(a)x_i^{\prime}(0)x_j^{\prime}(0) + \sum_i\partial_iL(a)x_i^{\prime\prime}(0) = D^2L(a)(d, d),
\]
because $\nabla L(a) = 0$; this is where the Lagrangian, rather than $f$, is needed, since $x''(0)$ is unknown. Now $\varphi$ has a local minimum at $0$, and as in the proof of
Proposition~\ref{real-analysis:proposition-hessian-test}(1), $\varphi''(0) \geq 0$.

(2) Suppose not. Then there are $x_k \in M$ with $x_k \neq a$, $x_k \to a$ and $f(x_k) \leq f(a)$. Write $x_k = a + t_kd_k$ with $t_k = \|x_k - a\| \to 0$ and $\|d_k\| = 1$; passing to a subsequence
(the unit sphere is compact), $d_k \to d$ with $\|d\| = 1$. From $h(x_k) = h(a) = 0$ we get $0 = (h(a + t_kd_k) - h(a))/t_k \to Dh(a)d$, so $d \in T$. By \eqref{real-analysis:equation-taylor-hessian} for $L$,
with $\nabla L(a) = 0$,
\[
0 \geq f(x_k) - f(a) = L(x_k) - L(a) = \tfrac12t_k^2D^2L(\xi_k)(d_k, d_k)
\]
with $\xi_k \to a$. Dividing by $t_k^2/2$ and letting $k \to \infty$, $D^2L(a)(d, d) \leq 0$, contradicting the hypothesis.
\end{proof}

\begin{example}\label{real-analysis:example-second-order}
Let $M$ be the parabola $y = x^2$, $h(x, y) = y - x^2$, and $f(x, y) = y + cx^2$ with a constant $c$. At $a = 0$, $\nabla f = (0, 1) = 1 \cdot \nabla h(0)$, so $\lambda = 1$ and $T = \RR(1, 0)$. Then
$D^2L = D^2f - D^2h = \operatorname{diag}(2c + 2, 0)$, so $D^2L(0)(e_1, e_1) = 2(1 + c)$: by the theorem, $0$ is a strict local minimum of $f|_M$ for $c > -1$ and a strict local maximum for $c < -1$.
This is correct, since $f = (1 + c)x^2$ on $M$. The Hessian of $f$ alone gives $D^2f(0)(e_1, e_1) = 2c$, which has the wrong sign for $-1 < c < 0$: the term $-\lambda D^2h$ records the bending of $M$.
\end{example}

\begin{example}[The bordered Hessian in two variables]\label{real-analysis:example-bordered-hessian}
Let $n = 2$, $q = 1$, $\nabla h(a) = (h_x, h_y) \neq 0$ and $D^2L(a) = \bigl(\begin{smallmatrix} L_{xx} & L_{xy}\\ L_{xy} & L_{yy}\end{smallmatrix}\bigr)$. Then $T$ is spanned by $d = (-h_y, h_x)$ and
\[
D^2L(a)(d, d) = L_{xx}h_y^2 - 2L_{xy}h_xh_y + L_{yy}h_x^2 = -\det\begin{pmatrix} 0 & h_x & h_y\\ h_x & L_{xx} & L_{xy}\\ h_y & L_{xy} & L_{yy}\end{pmatrix},
\]
as one sees by expanding the determinant along the first row. So if the \emph{bordered Hessian} on the right has negative determinant, $a$ is a strict local minimum of $f|_M$, and if it has positive
determinant, a strict local maximum. In Example~\ref{real-analysis:example-lagrange}(1), at $\pm(\frac1{\sqrt2}, \frac1{\sqrt2})$ we have $\lambda = \pm\frac1{\sqrt2}$, $D^2L = -2\lambda I$ and
$(h_x, h_y) = \pm(\sqrt2, \sqrt2)$; the determinant is $-(-2\lambda)(h_x^2 + h_y^2) = 8\lambda$, positive at the maximum and negative at the minimum.
\end{example}

\subsection{Inequality constraints and Farkas' lemma}\label{real-analysis:subsection-kkt}

Many problems constrain the unknowns by inequalities: quantities are nonnegative, budgets are bounded. Consider
\begin{equation}\label{real-analysis:equation-nlp}
\text{minimise } f(x) \quad \text{subject to} \quad g_i(x) \leq 0 \ (1 \leq i \leq p), \qquad h_j(x) = 0 \ (1 \leq j \leq q),
\end{equation}
where $f, g_i, h_j \colon U \to \RR$ are $C^1$ on an open $U \subseteq \RR^n$; the \emph{feasible set} is the set $F$ of $x \in U$ satisfying the constraints. At a feasible point $a$, the inequality $g_i \leq 0$
is \emph{active} if $g_i(a) = 0$; let $I(a)$ be the set of active indices. An inactive constraint, $g_i(a) < 0$, still holds near $a$ and plays no role locally; an active one behaves like a wall.

To see what to expect, take one inequality $g \leq 0$ and no equations. If the minimum $a$ is interior ($g(a) < 0$), then $\nabla f(a) = 0$. If $g(a) = 0$, the minimum lies on the wall $\{g = 0\}$, so
$\nabla f(a) = \lambda\nabla g(a)$ by Theorem~\ref{real-analysis:theorem-lagrange}; and now the \emph{sign} of $\lambda$ is determined: moving into the feasible region, in the direction $-\nabla g(a)$, must not
decrease $f$, so $\nabla f(a)$ points into the feasible region and $-\nabla f(a) = \mu\nabla g(a)$ with $\mu \geq 0$. The general statement needs a lemma of linear algebra over the ordered field $\RR$.

\begin{lemma}[Farkas]\label{real-analysis:lemma-farkas}
Let $v_1, \ldots, v_r, b \in \RR^n$. Exactly one of the following holds:
\begin{enumerate}
\item $b = c_1v_1 + \cdots + c_rv_r$ with all $c_i \geq 0$;
\item there is $w \in \RR^n$ with $\langle w, v_i\rangle \leq 0$ for all $i$ and $\langle w, b\rangle > 0$.
\end{enumerate}
\end{lemma}

Geometrically: either $b$ lies in the convex cone $C = \{\sum_ic_iv_i : c_i \geq 0\}$, or a hyperplane through $0$, with normal $w$, separates $b$ from $C$.

\begin{proof}
Both cannot hold, since then $\langle w, b\rangle = \sum_ic_i\langle w, v_i\rangle \leq 0$. Assume (1) fails, that is, $b \notin C$.

\emph{Step 1: $C$ is closed.} First, every $y \in C$ is a nonnegative combination of a linearly independent subset of the $v_i$ (\emph{Carath\'eodory's lemma}). Indeed, choose
$y = \sum_{i \in S}c_iv_i$ with all $c_i > 0$ and $S$ as small as possible. If the $v_i$ with $i \in S$ were dependent, there would be a relation $\sum_{i \in S}\gamma_iv_i = 0$ with some $\gamma_i > 0$
(change the sign of $\gamma$ if necessary); with $t = \min\{c_i/\gamma_i : \gamma_i > 0\}$, $y = \sum_{i \in S}(c_i - t\gamma_i)v_i$ has nonnegative coefficients, one of which vanishes, contradicting the
minimality of $S$. So $C$ is the union of the finitely many cones $C_S$ spanned by linearly independent subsets $S$, and it suffices that each $C_S$ is closed. The map $\RR^S \to \RR^n$,
$c \mapsto \sum_{i \in S}c_iv_i$, is injective and linear, so on the compact unit sphere of $\RR^S$ its norm has a positive minimum $m$, and $\|\sum_ic_iv_i\| \geq m\|c\|$. If
$\sum_ic_i^{(k)}v_i \to y$ with $c^{(k)} \geq 0$, then $(c^{(k)})$ is a Cauchy sequence, with a limit $c \geq 0$, and $y = \sum_ic_iv_i \in C_S$.

\emph{Step 2: a nearest point.} The continuous function $x \mapsto \|b - x\|$ attains a minimum on the compact set $C \cap \overline{B}(b, \|b\|)$, which contains $0$; the minimum point $p$ is then a
nearest point to $b$ in all of $C$.

\emph{Step 3: separation.} Let $w = b - p$, nonzero since $b \notin C$. For $c \in C$ and $0 < t \leq 1$, $(1 - t)p + tc \in C$, so
$\|w - t(c - p)\|^2 \geq \|w\|^2$, that is, $-2t\langle w, c - p\rangle + t^2\|c - p\|^2 \geq 0$; dividing by $t$ and letting $t \to 0$, $\langle w, c - p\rangle \leq 0$. With $c = 0$ and $c = 2p$ this gives
$\langle w, p\rangle = 0$, hence $\langle w, c\rangle \leq 0$ for all $c \in C$, in particular $\langle w, v_i\rangle \leq 0$. Finally $\langle w, b\rangle = \langle w, w + p\rangle = \|w\|^2 > 0$.
\end{proof}

\begin{theorem}[Karush--Kuhn--Tucker]\label{real-analysis:theorem-kkt}
In the situation \eqref{real-analysis:equation-nlp}, let $a \in F$ be a local minimum of $f$ on $F$, and assume the \emph{linear independence constraint qualification}: the gradients $\nabla g_i(a)$ for
$i \in I(a)$ and $\nabla h_j(a)$ for $1 \leq j \leq q$ are linearly independent. Then there are unique numbers $\mu_1, \ldots, \mu_p$ and $\lambda_1, \ldots, \lambda_q$ with
\begin{equation}\label{real-analysis:equation-kkt}
\nabla f(a) + \sum_{i=1}^p\mu_i\nabla g_i(a) + \sum_{j=1}^q\lambda_j\nabla h_j(a) = 0, \qquad \mu_i \geq 0, \qquad \mu_ig_i(a) = 0 \quad (1 \leq i \leq p).
\end{equation}
\end{theorem}

The last condition, \emph{complementary slackness}, says that $\mu_i = 0$ unless the $i$-th constraint is active. The signs of the $\lambda_j$ are not restricted, and for $p = 0$ the theorem is
Theorem~\ref{real-analysis:theorem-lagrange} for minima (with $-\lambda_j$ in place of $\lambda_j$).

\begin{proof}
Write $I = I(a)$.

\emph{Step 1.} Let $d \in \RR^n$ with $\langle\nabla g_i(a), d\rangle < 0$ for $i \in I$ and $Dh(a)d = 0$. We show $\langle\nabla f(a), d\rangle \geq 0$. The $\nabla h_j(a)$ are linearly independent, so
Lemma~\ref{real-analysis:lemma-curves-level-set} gives a $C^1$ curve $x(t)$ with $x(0) = a$, $x'(0) = d$ and $h(x(t)) = 0$. For $i \in I$, $\varphi_i(t) = g_i(x(t))$ has $\varphi_i(0) = 0$ and
$\varphi_i'(0) = \langle\nabla g_i(a), d\rangle < 0$, so $\varphi_i(t) < 0$ for small $t > 0$; for $i \notin I$, $g_i(x(t)) < 0$ for small $t$ by continuity. So $x(t) \in F$ for small $t > 0$, and since
$x(t) \to a$, $f(x(t)) \geq f(a)$ for small $t > 0$. Hence $\langle\nabla f(a), d\rangle = \lim_{t \to 0^+}(f(x(t)) - f(a))/t \geq 0$.

\emph{Step 2.} The same holds if only $\langle\nabla g_i(a), d\rangle \leq 0$ for $i \in I$ and $Dh(a)d = 0$. Indeed, the linear map $d \mapsto ((\langle\nabla g_i(a), d\rangle)_{i \in I}, Dh(a)d)$ has linearly
independent rows, hence is surjective, so there is $d_0$ with $\langle\nabla g_i(a), d_0\rangle = -1$ for $i \in I$ and $Dh(a)d_0 = 0$. For $\varepsilon > 0$, $d + \varepsilon d_0$ satisfies the hypothesis of
Step~1, so $\langle\nabla f(a), d + \varepsilon d_0\rangle \geq 0$; let $\varepsilon \to 0$.

\emph{Step 3.} Apply Lemma~\ref{real-analysis:lemma-farkas} to $b = -\nabla f(a)$ and the vectors $\nabla g_i(a)$ ($i \in I$), $\nabla h_j(a)$ and $-\nabla h_j(a)$ ($1 \leq j \leq q$). Alternative (2) would give
$w$ with $\langle\nabla g_i(a), w\rangle \leq 0$ for $i \in I$, $\langle\pm\nabla h_j(a), w\rangle \leq 0$, that is, $Dh(a)w = 0$, and $\langle\nabla f(a), w\rangle < 0$, contradicting Step~2. So alternative (1)
holds: $-\nabla f(a) = \sum_{i \in I}\mu_i\nabla g_i(a) + \sum_j(c_j^+ - c_j^-)\nabla h_j(a)$ with $\mu_i, c_j^\pm \geq 0$. Put $\lambda_j = c_j^+ - c_j^-$ and $\mu_i = 0$ for $i \notin I$; this gives
\eqref{real-analysis:equation-kkt}. Uniqueness follows from the linear independence, since $\mu_i = 0$ is forced for $i \notin I$.
\end{proof}

The conditions were found by W.~Karush in his 1939 master's thesis and, independently, by H.~W.~Kuhn and A.~W.~Tucker in 1951, at the start of nonlinear programming. In the language of the
introduction to this subsection: at a constrained minimum, $-\nabla f(a)$, the direction of steepest descent, lies in the cone spanned by the outward normals $\nabla g_i(a)$ of the active walls (and the
normal space of the equations), so that every feasible direction is uphill to first order.

\begin{counterexample}\label{real-analysis:counterexample-kkt}
Without a constraint qualification the conclusion can fail. Minimise $f(x, y) = x$ subject to $g_1 = y - x^3 \leq 0$ and $g_2 = -y \leq 0$. Feasibility means $0 \leq y \leq x^3$, so $x \geq 0$, and the minimum
is at $(0, 0)$, where both constraints are active. There $\nabla f = (1, 0)$, $\nabla g_1 = (0, 1)$ and $\nabla g_2 = (0, -1)$, and $(1, 0) + \mu_1(0, 1) + \mu_2(0, -1) = 0$ is impossible. The two active
gradients are dependent: the feasible set has a cusp at the origin, and the linearised constraints $\langle\nabla g_i, d\rangle \leq 0$, which allow every $d = (d_1, 0)$, do not describe the feasible
directions, which have $d_1 \geq 0$.
\end{counterexample}

\begin{remark}\label{real-analysis:remark-constraint-qualifications}
Reading the proof again shows that weaker hypotheses suffice for the existence of multipliers (not for their uniqueness, which used the full linear independence).
\begin{enumerate}
\item \emph{Mangasarian--Fromovitz condition:} the $\nabla h_j(a)$ are linearly independent and there is $d_0$ with $Dh(a)d_0 = 0$ and $\langle\nabla g_i(a), d_0\rangle < 0$ for $i \in I(a)$. Step~1
used only the independence of the $\nabla h_j(a)$, Step~2 only the existence of such a $d_0$, and Step~3 nothing else.
\item \emph{Affine constraints:} if all $g_i$ and $h_j$ are affine, no qualification is needed. Indeed, for $d$ with $\langle\nabla g_i(a), d\rangle \leq 0$ ($i \in I(a)$) and $Dh(a)d = 0$, the points
$a + td$ are feasible for small $t > 0$ (for active $i$, $g_i(a + td) = t\langle\nabla g_i(a), d\rangle \leq 0$), so $\langle\nabla f(a), d\rangle \geq 0$ directly, and Step~3 applies.
\end{enumerate}
\end{remark}

\subsection{Convex problems}\label{real-analysis:subsection-convex-problems}

The conditions of the Karush--Kuhn--Tucker theorem are necessary. For convex problems they are also sufficient, and give global minima. A function $f$ on a convex set $K \subseteq \RR^n$ is
\emph{convex} if $f((1 - t)x + ty) \leq (1 - t)f(x) + tf(y)$ for all $x, y \in K$ and $t \in [0, 1]$, that is, if its restriction to every segment is convex in the sense of
Definition~\ref{real-analysis:definition-convex}, and \emph{strictly convex} if the inequality is strict whenever $x \neq y$ and $0 < t < 1$. A sublevel set $\{x \in K : f(x) \leq c\}$ of a convex
function is convex, and so is the zero set of an affine function; hence the feasible set of \eqref{real-analysis:equation-nlp} is convex when $U$ is convex, the $g_i$ are convex and the $h_j$ are affine.

\begin{lemma}\label{real-analysis:lemma-convex-gradient}
Let $U \subseteq \RR^n$ be open and convex, and $f \colon U \to \RR$ convex and $C^1$. Then $f(y) \geq f(x) + \langle\nabla f(x), y - x\rangle$ for all $x, y \in U$. If $f$ is $C^2$ with positive
semidefinite Hessian everywhere, then $f$ is convex.
\end{lemma}

\begin{proof}
We may assume $x \neq y$. As $U$ is open, the points $x - \varepsilon(y - x)$ and $y + \varepsilon(y - x)$ lie in $U$ for small $\varepsilon > 0$, and as $U$ is convex so does the segment joining them.
Hence $\varphi(t) = f(x + t(y - x))$ is defined on the open interval $J = (-\varepsilon, 1 + \varepsilon)$, it is convex there (being the restriction of $f$ to a segment), and it is differentiable with
$\varphi'(0) = \langle\nabla f(x), y - x\rangle$ by the chain rule. By Proposition~\ref{real-analysis:proposition-convex}(2), $\varphi(1) \geq \varphi(0) + \varphi'(0)$, which is the inequality. For the
second statement, $\varphi''(t) = D^2f(x + t(y - x))(y - x, y - x) \geq 0$ on $J$ by \eqref{real-analysis:equation-line-derivatives}, so $\varphi$ is convex by
Proposition~\ref{real-analysis:proposition-convex}(3); this is convexity of $f$ on the segment $[x, y]$.
\end{proof}

\begin{theorem}\label{real-analysis:theorem-kkt-convex}
In \eqref{real-analysis:equation-nlp}, let $U$ be convex, $f$ and $g_1, \ldots, g_p$ convex and $h_1, \ldots, h_q$ affine. If $a \in F$ and numbers $\mu_i$, $\lambda_j$ satisfy
\eqref{real-analysis:equation-kkt}, then $a$ is a global minimum of $f$ on $F$. If $f$ is strictly convex, it is the only one.
\end{theorem}

\begin{proof}
Let $x \in F$. By Lemma~\ref{real-analysis:lemma-convex-gradient} and \eqref{real-analysis:equation-kkt},
\[
f(x) - f(a) \geq \langle\nabla f(a), x - a\rangle = -\sum_i\mu_i\langle\nabla g_i(a), x - a\rangle - \sum_j\lambda_j\langle\nabla h_j(a), x - a\rangle.
\]
Since $h_j$ is affine, $\langle\nabla h_j(a), x - a\rangle = h_j(x) - h_j(a) = 0$. By the lemma again and complementary slackness,
$\mu_i\langle\nabla g_i(a), x - a\rangle \leq \mu_i(g_i(x) - g_i(a)) = \mu_ig_i(x) \leq 0$. Hence $f(x) \geq f(a)$. If $f$ is strictly convex and $f(x) = f(a)$ with $x \neq a$, then the midpoint
$m = \frac12(x + a)$ is feasible ($F$ is convex) and $f(m) < \frac12f(x) + \frac12f(a) = f(a)$, a contradiction.
\end{proof}

\begin{example}[Water-filling]\label{real-analysis:example-water-filling}
Given $\alpha_1, \ldots, \alpha_n > 0$ and $P > 0$, maximise $\sum_i\log(\alpha_i + x_i)$ subject to $x_i \geq 0$ and $\sum_ix_i = P$. (In information theory, $x_i$ is the power sent through the $i$-th of
$n$ parallel channels with noise level $\alpha_i$, and $\frac12\log(1 + x_i/\alpha_i)$ is its capacity, so this maximises the total capacity for a total power $P$.) Minimise the strictly convex function
$f(x) = -\sum_i\log(\alpha_i + x_i)$ on $U = \{x_i > -\alpha_i\}$, with $g_i = -x_i$ and $h = \sum_ix_i - P$. The conditions \eqref{real-analysis:equation-kkt} read
\[
-\frac{1}{\alpha_i + x_i} - \mu_i + \lambda = 0, \qquad \mu_i \geq 0, \qquad \mu_ix_i = 0.
\]
So $\lambda \geq 1/(\alpha_i + x_i) > 0$; put $\nu = 1/\lambda$. If $x_i > 0$, then $\mu_i = 0$ and $x_i = \nu - \alpha_i$; if $x_i = 0$, then $\lambda \geq 1/\alpha_i$, that is, $\alpha_i \geq \nu$. In both cases
\[
x_i = \max(0, \nu - \alpha_i),
\]
and $\nu$ is fixed by $\sum_i\max(0, \nu - \alpha_i) = P$; the left side is continuous and strictly increasing in $\nu \geq \min_i\alpha_i$, from $0$ to $\infty$, so there is exactly one such $\nu$
(Theorem~\ref{real-analysis:theorem-bolzano-zeros}). Conversely this $x$, with $\lambda = 1/\nu$ and $\mu_i = \lambda - 1/(\alpha_i + x_i) \geq 0$, satisfies the conditions, so by
Theorem~\ref{real-analysis:theorem-kkt-convex} it is the unique optimum. The picture: a vessel whose floor has height $\alpha_i$ over the $i$-th cell is filled with a volume $P$ of water; the water
settles at level $\nu$, and $x_i$ is the depth over the $i$-th cell. The noisiest channels get nothing.
\end{example}

\begin{example}[Maximum entropy and the Boltzmann distribution]\label{real-analysis:example-gibbs}
Let $E_1, \ldots, E_N$ be real numbers (the energies of the states of a physical system), not all equal, and $\min_iE_i < E < \max_iE_i$. Among probability vectors $p$ with $p_i > 0$, $\sum_ip_i = 1$ and
mean energy $\sum_ip_iE_i = E$, which one has the largest \emph{entropy} $H(p) = -\sum_ip_i\log p_i$? The function $-H(p) = \sum_i\phi(p_i)$, $\phi(t) = t\log t$, is strictly convex on
$U = \{p_i > 0\}$: $\phi''(t) = 1/t > 0$, so $\phi$ is strictly convex (Proposition~\ref{real-analysis:proposition-convex}(3)), and for $p \neq p^{\prime}$ the convexity inequality is strict in a coordinate where
they differ. The two constraints are affine. The conditions
\eqref{real-analysis:equation-kkt} for minimising $-H$ read $\log p_i + 1 + \lambda_1 + \lambda_2E_i = 0$, so
\[
p_i = \frac{e^{-\beta E_i}}{Z(\beta)}, \qquad Z(\beta) = \sum_ke^{-\beta E_k},
\]
with $\beta = \lambda_2$, after normalising. This is the \emph{Boltzmann (Gibbs) distribution} of statistical mechanics, $\beta$ being the inverse temperature and $Z$ the \emph{partition function}.
It remains to choose $\beta$ with mean energy $E$. The mean $m(\beta) = \sum_kE_ke^{-\beta E_k}/Z(\beta)$ is continuous, has derivative $m'(\beta) = -(\sum_kE_k^2p_k - m(\beta)^2) < 0$, minus the
variance of the energy, which is positive since the $E_k$ are not all equal, and tends to $\min E_k$ as $\beta \to +\infty$ and to $\max E_k$ as $\beta \to -\infty$. So exactly one $\beta$ gives $m(\beta) = E$,
and by Theorem~\ref{real-analysis:theorem-kkt-convex} the corresponding $p$ is the unique entropy maximiser. Without the energy constraint ($\lambda_2 = 0$) the maximiser is the uniform
distribution, with $H = \log N$.
\end{example}

\begin{remark}[Duality]\label{real-analysis:remark-lagrange-duality}
For \eqref{real-analysis:equation-nlp} define the Lagrangian $L(x, \mu, \lambda) = f(x) + \sum_i\mu_ig_i(x) + \sum_j\lambda_jh_j(x)$ and the \emph{dual function} $\theta(\mu, \lambda) = \inf_{x \in U}L(x, \mu, \lambda)$.
For $\mu \geq 0$ and $x \in F$, $L(x, \mu, \lambda) \leq f(x)$, so
\[
\theta(\mu, \lambda) \leq \inf_Ff \qquad (\mu \geq 0) \qquad \text{(\emph{weak duality})}:
\]
every choice of multipliers gives a lower bound for the minimum. In the convex situation of Theorem~\ref{real-analysis:theorem-kkt-convex}, if $(a, \mu, \lambda)$ satisfies the conditions
\eqref{real-analysis:equation-kkt}, then $x \mapsto L(x, \mu, \lambda)$ is convex with vanishing gradient at $a$, so by Lemma~\ref{real-analysis:lemma-convex-gradient} it is minimal at $a$, and
$\theta(\mu, \lambda) = L(a, \mu, \lambda) = f(a)$ by complementary slackness and $h(a) = 0$: the bound is attained, and the multipliers solve the \emph{dual problem} of maximising $\theta$ over $\mu \geq 0$. For convex problems
with a strictly feasible point (\emph{Slater's condition}) the maximum of the dual equals the minimum of the primal even without assuming that the minimum is attained \cite[Chapter 5]{boyd-vandenberghe};
for linear programmes this is the duality theorem of linear programming, a consequence of Farkas' lemma. The point $(a, (\mu, \lambda))$ is then a saddle point of $L$: a minimum in $x$ and a
maximum in $(\mu, \lambda)$ with $\mu \geq 0$.
\end{remark}

\section{Ordinary differential equations}\label{real-analysis:section-ode}

\begin{definition}\label{real-analysis:definition-ode}
Let $D \subseteq \RR \times \RR^n$ be open and $F \colon D \to \RR^n$ continuous. $F$ is \emph{locally Lipschitz in $x$} if every point of $D$ has a neighbourhood $N$ and $L$ with
$\|F(t, x) - F(t, y)\| \leq L\|x - y\|$ for $(t, x), (t, y) \in N$; by the mean value inequality this holds if $F$ is $C^1$. A \emph{solution} of $x' = F(t, x)$ is a differentiable map $\varphi \colon J \to \RR^n$
on an interval with $(t, \varphi(t)) \in D$ and $\varphi'(t) = F(t, \varphi(t))$. By Theorem~\ref{real-analysis:theorem-ftc}, this is equivalent to continuity of $\varphi$ and
$\varphi(t) = \varphi(t_0) + \int_{t_0}^tF(s, \varphi(s))\,ds$.
\end{definition}

\begin{lemma}[Gr\"onwall]\label{real-analysis:lemma-gronwall}
If $u \colon [t_0, t_1] \to [0, \infty)$ is continuous and $u(t) \leq C + L\int_{t_0}^tu$ with $C, L \geq 0$, then $u(t) \leq Ce^{L(t - t_0)}$.
\end{lemma}

\begin{proof}
Let $v(t) = C + L\int_{t_0}^tu$. Then $v' = Lu \leq Lv$, so $(e^{-L(t - t_0)}v)' \leq 0$, and $u \leq v \leq v(t_0)e^{L(t - t_0)} = Ce^{L(t - t_0)}$.
\end{proof}

\begin{theorem}[Picard--Lindel\"of]\label{real-analysis:theorem-picard-lindelof}
Let $F$ be continuous and locally Lipschitz in $x$, and $(t_0, x_0) \in D$.
\begin{enumerate}
\item There is $\delta > 0$ and a solution on $[t_0 - \delta, t_0 + \delta]$ with $\varphi(t_0) = x_0$.
\item Two solutions on intervals $J_1, J_2 \ni t_0$ with the same value at $t_0$ agree on $J_1 \cap J_2$.
\item There is a unique \emph{maximal} solution, defined on an open interval $(\alpha, \beta)$; if $\beta < \infty$, then $(t, \varphi(t))$ leaves every compact subset of $D$ as $t \to \beta$, and similarly at $\alpha$.
\end{enumerate}
\end{theorem}

\begin{proof}
(1) Choose $a, b > 0$ such that $R = [t_0 - a, t_0 + a] \times \overline{B}(x_0, b) \subseteq D$ and $F$ is $L$-Lipschitz in $x$ on $R$; let $M = \sup_R\|F\|$ (finite by compactness) and
$\delta = \min(a, b/(M+1), 1/(2L+1))$. On the complete metric space $X$ of continuous maps $[t_0 - \delta, t_0 + \delta] \to \overline{B}(x_0, b)$ with the supremum metric
(Proposition~\ref{general-topology:proposition-complete-examples}), $T\varphi(t) = x_0 + \int_{t_0}^tF(s, \varphi(s))\,ds$ satisfies $\|T\varphi(t) - x_0\| \leq M\delta \leq b$ and
$\|T\varphi - T\psi\| \leq L\delta\|\varphi - \psi\| \leq \frac12\|\varphi - \psi\|$. Its fixed point is a solution.

(2) The set of $t \in J_1 \cap J_2$ where the solutions agree is closed and contains $t_0$. If they agree at $t_1$, near $(t_1, \varphi(t_1))$ $F$ is $L$-Lipschitz, and $u = \|\varphi_1 - \varphi_2\|$ satisfies
$u(t) \leq L|\int_{t_1}^tu|$ for $t$ near $t_1$, so $u = 0$ near $t_1$ by Lemma~\ref{real-analysis:lemma-gronwall} (with $C = 0$, on both sides). So the set is open, and equals the connected set
$J_1 \cap J_2$.

(3) By (2), the solutions with $\varphi(t_0) = x_0$ glue to a solution on the union of their intervals, which is maximal; the interval is open because a solution on a closed end could be continued by (1).
Suppose $\beta < \infty$ and $(t_k, \varphi(t_k)) \in K$ for a compact $K \subseteq D$ and $t_k \to \beta$. By compactness there are $a, b, L, M$ as in (1) valid for all rectangles centred at points of $K$
(cover $K$ by finitely many rectangles of half the size), hence a uniform $\delta$. For $t_k > \beta - \delta/2$, the solution through $(t_k, \varphi(t_k))$ exists up to $t_k + \delta > \beta$ and extends
$\varphi$ by (2), a contradiction.
\end{proof}

\begin{counterexample}\label{real-analysis:counterexample-ode}
\begin{enumerate}
\item Without the Lipschitz condition uniqueness fails: $x' = 3x^{2/3}$, $x(0) = 0$ has the solutions $x = 0$ and $x = t^3$.
\item Solutions need not exist for all time: $x' = x^2$, $x(0) = 1$ has maximal solution $x = 1/(1 - t)$ on $(-\infty, 1)$.
\end{enumerate}
\end{counterexample}

\begin{proposition}\label{real-analysis:proposition-ode-dependence}
\begin{enumerate}
\item If $\varphi, \psi$ are solutions on $[t_0, t_1]$ and $\|F(t, \varphi(t)) - F(t, \psi(t))\| \leq L\|\varphi(t) - \psi(t)\|$ for all $t$, then
$\|\varphi(t) - \psi(t)\| \leq \|\varphi(t_0) - \psi(t_0)\|e^{L(t - t_0)}$.
\item (Linear equations) If $A \colon I \to M_n(\RR)$ and $b \colon I \to \RR^n$ are continuous on an open interval $I$, the maximal solutions of $x' = A(t)x + b(t)$ are defined on all of $I$.
\end{enumerate}
\end{proposition}

\begin{proof}
(1) $u = \|\varphi - \psi\|$ satisfies $u(t) \leq u(t_0) + L\int_{t_0}^tu$; apply Lemma~\ref{real-analysis:lemma-gronwall}. (2) On a compact subinterval $[c, d] \ni t_0$ with $\|A\| \leq L$ and $\|b\| \leq B$,
$\|\varphi(t)\| \leq \|\varphi(t_0)\| + B(d - c) + L\int\|\varphi\|$, so $\varphi$ is bounded there by Gr\"onwall, its graph stays in a compact set, and by Theorem~\ref{real-analysis:theorem-picard-lindelof}(3)
the maximal interval contains $[c, d]$.
\end{proof}

\begin{theorem}[Smooth dependence]\label{real-analysis:theorem-smooth-dependence}
If $F$ is $C^k$ with $k \geq 1$, the \emph{flow} $\Phi(t, t_0, x_0)$, the value at $t$ of the maximal solution with $\varphi(t_0) = x_0$, is defined on an open subset of $\RR \times D$ and is $C^k$
there. Its derivative in $x_0$ solves the linearised equation $Y' = D_xF(t, \varphi(t))Y$, $Y(t_0) = 1$.
\end{theorem}

This is \cite[Chapter V]{hartman-ode}; see also \cite{lang-real-functional}. Continuity of $\Phi$ follows from Proposition~\ref{real-analysis:proposition-ode-dependence}(1) and the uniform $\delta$ in
the proof of Theorem~\ref{real-analysis:theorem-picard-lindelof}.

\section{Stone--Weierstrass and Arzel\`a--Ascoli}\label{real-analysis:section-approximation}

\begin{lemma}\label{real-analysis:lemma-absolute-value-polynomials}
The polynomials $p_0 = 0$, $p_{k+1}(t) = p_k(t) + (t^2 - p_k(t)^2)/2$ converge uniformly to $|t|$ on $[-1, 1]$, and $p_k(0) = 0$.
\end{lemma}

\begin{proof}
By induction $0 \leq p_k(t) \leq p_{k+1}(t) \leq |t|$: indeed $|t| - p_{k+1} = (|t| - p_k)(1 - (|t| + p_k)/2)$ with $0 \leq (|t| + p_k)/2 \leq |t| \leq 1$, and $p_{k+1} - p_k = (t^2 - p_k^2)/2 \geq 0$. Iterating,
$0 \leq |t| - p_k(t) \leq |t|(1 - |t|/2)^k$, and the maximum of $s(1 - s/2)^k$ on $[0, 1]$ is at most $2/(k + 1)$ (it is attained at $s = 2/(k+1)$ when $k \geq 1$).
\end{proof}

\begin{theorem}[Stone--Weierstrass]\label{real-analysis:theorem-stone-weierstrass}
Let $K$ be a compact Hausdorff space and $\mathcal{A} \subseteq C(K, \RR)$ a subalgebra containing the constants and separating points. Then $\mathcal{A}$ is dense for the supremum norm. The same holds
for $\mathcal{A} \subseteq C(K, \CC)$ if $\mathcal{A}$ is also closed under complex conjugation.
\end{theorem}

\begin{proof}
The closure $\bar{\mathcal{A}}$ is a closed subalgebra. If $f \in \bar{\mathcal{A}}$ with $\|f\| \leq c$, then $|f| = c\lim_kp_k(f/c) \in \bar{\mathcal{A}}$ by Lemma~\ref{real-analysis:lemma-absolute-value-polynomials};
hence $\max(f, g) = (f + g + |f - g|)/2$ and $\min(f, g)$ lie in $\bar{\mathcal{A}}$. For $x \neq y$ and reals $\alpha, \beta$ there is $h \in \mathcal{A}$ with $h(x) = \alpha$, $h(y) = \beta$ (take
$h = \alpha + (\beta - \alpha)(u - u(x))/(u(y) - u(x))$ with $u$ separating $x$ and $y$).

Let $f \in C(K, \RR)$ and $\varepsilon > 0$. Fix $x$. For each $y$ choose $h_y \in \mathcal{A}$ with $h_y(x) = f(x)$ and $h_y(y) = f(y)$; the open sets $\{h_y < f + \varepsilon\}$ cover $K$, and the minimum $g_x$ of
finitely many of these $h_y$ lies in $\bar{\mathcal{A}}$, with $g_x(x) = f(x)$ and $g_x < f + \varepsilon$. The open sets $\{g_x > f - \varepsilon\}$ cover $K$, and the maximum $g$ of finitely many $g_x$ lies in
$\bar{\mathcal{A}}$ and satisfies $\|g - f\| < \varepsilon$. For the complex case, the real parts $(f + \bar f)/2$ of elements of $\mathcal{A}$ form a real subalgebra with the same properties, which is dense
in $C(K, \RR)$; so $\mathcal{A}$ is dense in $C(K, \CC)$.
\end{proof}

\begin{corollary}\label{real-analysis:corollary-weierstrass}
Polynomials are dense in $C(K, \RR)$ for every compact $K \subseteq \RR^n$, and trigonometric polynomials $\sum_{|k| \leq N}c_ke^{2\pi ikt}$ are dense in the space of continuous functions $\RR \to \CC$ of
period $1$.
\end{corollary}

\begin{proof}
The first statement is immediate. For the second, continuous functions of period $1$ are continuous functions on $S^1 = \RR/\ZZ$ (Theorem~\ref{real-analysis:theorem-pi}), and the polynomials in $z$ and $\bar z$ restricted to
$S^1$ form a conjugation-closed separating algebra; on $S^1$, $z^k\bar z^l = z^{k-l}$ corresponds to $e^{2\pi i(k-l)t}$.
\end{proof}

\begin{theorem}[Arzel\`a--Ascoli]\label{real-analysis:theorem-arzela-ascoli}
Let $K$ be a compact Hausdorff space and $\mathcal{F} \subseteq C(K, \RR^m)$ a family that is \emph{equicontinuous} (for every $x$ and $\varepsilon > 0$ there is a neighbourhood $U$ of $x$ with
$\|f(x') - f(x)\| < \varepsilon$ for all $x' \in U$ and $f \in \mathcal{F}$) and pointwise bounded. Then $\mathcal{F}$ has compact closure in $C(K, \RR^m)$ with the supremum metric.
\end{theorem}

\begin{proof}
$C(K, \RR^m)$ is complete, so by Theorem~\ref{general-topology:theorem-compact-metric} it suffices to show that $\mathcal{F}$ is totally bounded. Let $\varepsilon > 0$. Choose finitely many points $x_1, \ldots, x_k$
with neighbourhoods $U_i$ as in the definition covering $K$. The set $\{(f(x_1), \ldots, f(x_k)) : f \in \mathcal{F}\}$ is bounded in $\RR^{mk}$, hence totally bounded, so there are $f_1, \ldots, f_N \in \mathcal{F}$
such that every $f \in \mathcal{F}$ satisfies $\max_i\|f(x_i) - f_j(x_i)\| < \varepsilon$ for some $j$. For $x \in U_i$,
$\|f(x) - f_j(x)\| \leq \|f(x) - f(x_i)\| + \|f(x_i) - f_j(x_i)\| + \|f_j(x_i) - f_j(x)\| < 3\varepsilon$.
\end{proof}

\section{Exercises}\label{real-analysis:section-exercises}

\begin{exercise}\label{real-analysis:exercise-e-irrational}
Show that $e = \exp(1)$ is irrational.
\end{exercise}

\begin{solution}
If $e = p/q$, then $x = q!(e - \sum_{k \leq q}1/k!)$ is a positive integer, but $x = \sum_{k > q}q!/k! < \sum_{j \geq 1}(q + 1)^{-j} = 1/q \leq 1$.
\end{solution}

\begin{exercise}\label{real-analysis:exercise-harmonic}
Show that $\sum 1/n$ diverges while $\sum(-1)^n/n$ converges, and that the latter series can be rearranged to converge to any given real number.
\end{exercise}

\begin{solution}
$\sum_{2^j < n \leq 2^{j+1}}1/n \geq 1/2$, so the partial sums are unbounded; the alternating series converges by Lemma~\ref{real-analysis:lemma-alternating}. For a target $c$, add positive terms
$1/(2k)$ until the partial sum exceeds $c$, then negative terms until it is below $c$, and so on; since both the positive and the negative parts diverge and the terms tend to $0$, the partial sums converge
to $c$. This is the special case $\alpha = \beta = c$ of Theorem~\ref{real-analysis:theorem-riemann-rearrangement}, where the details are written out.
\end{solution}

\begin{exercise}\label{real-analysis:exercise-bounded-rearrangement}
Let $\sum_na_n$ be a convergent series of real numbers and $\sigma \colon \NN \to \NN$ a bijection that moves every index by a bounded amount: $|\sigma(n) - n| \leq K$ for all $n$. Show that
$\sum_na_{\sigma(n)}$ converges to the same sum. Deduce that the rearrangement of Example~\ref{real-analysis:example-rearrangement}(1) moves some indices arbitrarily far.
\end{exercise}

\begin{solution}
Let $S_N = \sum_{n \leq N}a_n$ and $T_N = \sum_{n \leq N}a_{\sigma(n)} = \sum_{m \in \sigma(\{0, \ldots, N\})}a_m$. If $m \leq N - K$, then $\sigma^{-1}(m) \leq m + K \leq N$, so $m \in \sigma(\{0, \ldots, N\})$;
and every $m \in \sigma(\{0, \ldots, N\})$ satisfies $m \leq N + K$. So $\{0, \ldots, N\}$ and $\sigma(\{0, \ldots, N\})$ differ only in indices between $N - K$ and $N + K$, and
$|T_N - S_N| \leq \sum_{N - K < m \leq N + K}|a_m|$, a sum of at most $2K$ terms, which tends to $0$ as $a_m \to 0$. Hence $T_N - S_N \to 0$. The rearrangement of the example changes the sum, so it
cannot move indices by a bounded amount; indeed its $3m$-th term is $-1/(4m)$, the $4m$-th term of $S$, moved by $m$ places.
\end{solution}

\begin{exercise}\label{real-analysis:exercise-uniform-limit-derivative}
Show that $f_n(x) = \sqrt{x^2 + 1/n}$ are smooth and converge uniformly on $\RR$ to $|x|$, which is not differentiable. Prove that if $f_n \to f$ pointwise and $f_n' \to g$ uniformly on an interval, with
$f_n$ continuously differentiable, then $f$ is differentiable with $f' = g$.
\end{exercise}

\begin{solution}
$0 \leq \sqrt{x^2 + 1/n} - |x| = \frac{1/n}{\sqrt{x^2 + 1/n} + |x|} \leq 1/\sqrt{n}$. For the second part, $f_n(x) = f_n(a) + \int_a^xf_n'$ converges to $f(a) + \int_a^xg$, so $f(x) = f(a) + \int_a^xg$ and
Theorem~\ref{real-analysis:theorem-ftc} applies.
\end{solution}

\begin{exercise}\label{real-analysis:exercise-polar}
Show that $(r, \theta) \mapsto (r\cos\theta, r\sin\theta)$ is a local diffeomorphism $(0, \infty) \times \RR \to \RR^2 \setminus \{0\}$, but not injective, and compute its Jacobian determinant.
\end{exercise}

\begin{exercise}\label{real-analysis:exercise-peano}
(Peano's theorem.) Show that if $F$ is only continuous, $x' = F(t, x)$ still has local solutions, by applying the Arzel\`a--Ascoli theorem to the Euler polygons $\varphi_k$ with
step $1/k$.
\end{exercise}

\begin{exercise}\label{real-analysis:exercise-arzela-converse}
Show that a compact subset of $C(K, \RR^m)$ is equicontinuous and pointwise bounded.
\end{exercise}
\begin{exercise}\label{real-analysis:exercise-sequence-limits}
Compute $\lim_n\frac{n^2 + 1}{2n^2 - n}$, $\lim_n(\sqrt{n^2 + n} - n)$ and $\lim_n(2^n + 3^n)^{1/n}$.
\end{exercise}

\begin{solution}
Dividing by $n^2$, the first is $\frac{1 + 1/n^2}{2 - 1/n} \to \frac12$. For the second, $\sqrt{n^2 + n} - n = \frac{n}{\sqrt{n^2 + n} + n} = \frac{1}{\sqrt{1 + 1/n} + 1} \to \frac12$. For the third,
$3 \leq (2^n + 3^n)^{1/n} \leq (2 \cdot 3^n)^{1/n} = 3 \cdot 2^{1/n} \to 3$ (Proposition~\ref{real-analysis:proposition-fundamental-limits}(2)), so the limit is $3$ by the squeeze theorem.
\end{solution}

\begin{exercise}\label{real-analysis:exercise-limsup-subsequences}
Show that $\limsup_na_n$ is the largest element of $\bar{\RR}$ that is a limit of a subsequence of $(a_n)$, and $\liminf_na_n$ the smallest.
\end{exercise}

\begin{solution}
Let $S_n = \sup_{k \geq n}a_k$, a decreasing sequence with limit $L = \limsup a_n$ (Proposition~\ref{real-analysis:proposition-monotone-regular}). If $a_{n_k} \to \ell$, then $a_{n_k} \leq S_n$ for $n_k \geq n$, so $\ell \leq S_n$ for
all $n$, and $\ell \leq L$. Conversely, if $L \in \RR$, choose recursively $n_k > n_{k-1}$ with $S_{n_{k-1}+1} - \frac1k < a_{n_k} \leq S_{n_{k-1}+1}$; both bounds tend to $L$, so $a_{n_k} \to L$ by the squeeze theorem. If
$L = +\infty$ the sequence is unbounded above and Theorem~\ref{real-analysis:theorem-bolzano-weierstrass-real} applies; if $L = -\infty$ then $a_n \leq S_n \to -\infty$. The statement for $\liminf$ follows by passing to
$(-a_n)$.
\end{solution}

\begin{exercise}\label{real-analysis:exercise-continuous-one-point}
Show that $f(x) = x\mathbf{1}_\QQ(x)$ is continuous at $0$ and discontinuous at every $a \neq 0$.
\end{exercise}

\begin{solution}
$|f(x)| \leq |x|$, so $f(x) \to 0 = f(0)$ as $x \to 0$. For $a \neq 0$, choose rational $x_n \to a$ and irrational $y_n \to a$ (Counterexample~\ref{real-analysis:counterexample-no-limit}); then $f(x_n) = x_n \to a$ and
$f(y_n) = 0$, so $f$ has no limit at $a$ (Theorem~\ref{real-analysis:theorem-sequential-limits}), in particular it is not continuous there.
\end{solution}

\begin{exercise}\label{real-analysis:exercise-cauchy-functional}
Let $f \colon \RR \to \RR$ be continuous with $f(x + y) = f(x) + f(y)$ for all $x, y$. Show that $f(x) = cx$ with $c = f(1)$.
\end{exercise}

\begin{solution}
$f(0) = f(0) + f(0)$ gives $f(0) = 0$, then $f(-x) = -f(x)$, and by induction $f(nx) = nf(x)$ for $n \in \NN$, hence for $n \in \ZZ$. For $x = p/q$ with $q \geq 1$, $qf(p/q) = f(p) = pc$, so $f(x) = cx$ on $\QQ$. For
real $x$, choose rational $x_n \to x$ (density of $\QQ$); by continuity $f(x) = \lim f(x_n) = \lim cx_n = cx$.
\end{solution}

\begin{exercise}\label{real-analysis:exercise-periodic-uniform}
Show that a continuous periodic function $f \colon \RR \to \RR$ is bounded and uniformly continuous.
\end{exercise}

\begin{solution}
Let $p > 0$ be a period. On $[0, 2p]$, $f$ is bounded and uniformly continuous (Theorem~\ref{real-analysis:theorem-weierstrass-extreme} and Proposition~\ref{real-analysis:proposition-uniform-continuity-real}(3)), and
every value of $f$ is taken on $[0, p]$. Given $\varepsilon$, let $\delta < p$ work on $[0, 2p]$. If $|x - y| < \delta$, translating both by the same multiple of $p$ we may assume $0 \leq \min(x, y) < p$, and then
$x, y \in [0, 2p]$, so $|f(x) - f(y)| < \varepsilon$.
\end{solution}

\begin{exercise}\label{real-analysis:exercise-log-limit}
Show that $n(a^{1/n} - 1) \to \log a$ for every $a > 0$.
\end{exercise}

\begin{solution}
Let $t = \log a$, so $a^{1/n} = e^{t/n}$. If $t = 0$ both sides vanish. Otherwise $n(e^{t/n} - 1) = t\cdot\frac{e^{t/n} - 1}{t/n} \to t$ by Proposition~\ref{real-analysis:proposition-notable-limits}(1) and
Theorem~\ref{real-analysis:theorem-sequential-limits}, since $t/n \to 0$ and $t/n \neq 0$.
\end{solution}

\begin{exercise}\label{real-analysis:exercise-integration-by-parts}
Compute $\int_0^\pi x^2\sin x\,dx$, and show by integration by parts and induction on $n$ that $\int_0^1x^m(1 - x)^n\,dx = \frac{m!\,n!}{(m + n + 1)!}$ for integers $m, n \geq 0$.
\end{exercise}

\begin{solution}
Twice by parts: $\int_0^\pi x^2\sin x\,dx = [-x^2\cos x]_0^\pi + 2\int_0^\pi x\cos x\,dx = \pi^2 + 2\bigl([x\sin x]_0^\pi - \int_0^\pi\sin x\,dx\bigr) = \pi^2 + 2(0 - 2) = \pi^2 - 4$. For the second, let
$B(m, n)$ be the integral. For $n \geq 1$, take $f(x) = x^{m+1}/(m + 1)$ and $g(x) = (1 - x)^n$; the boundary term vanishes at both ends, and $B(m, n) = \frac{n}{m+1}B(m + 1, n - 1)$. Iterating,
$B(m, n) = \frac{n(n-1) \cdots 1}{(m+1)(m+2) \cdots (m+n)}B(m + n, 0)$, and $B(m + n, 0) = \frac{1}{m + n + 1}$, which gives $\frac{m!\,n!}{(m + n + 1)!}$.
\end{solution}

\begin{exercise}\label{real-analysis:exercise-midpoint-convex}
Let $f \colon I \to \RR$ be continuous on an interval and \emph{midpoint convex}: $f(\frac{x + y}2) \leq \frac{f(x) + f(y)}2$ for all $x, y \in I$. Show that $f$ is convex. (Without continuity this
fails: a $\QQ$-linear map $\RR \to \RR$ that is not $\RR$-linear, which exists by choosing a basis of $\RR$ over $\QQ$, is midpoint affine but not convex.)
\end{exercise}

\begin{solution}
By induction on $k$, $f(\sum_ix_i/2^k) \leq \sum_if(x_i)/2^k$ for $2^k$ points $x_i \in I$ (split them into two halves and apply midpoint convexity to the two averages). Given $x, y$ and a dyadic
$t = j/2^k \in [0, 1]$, take $j$ copies of $x$ and $2^k - j$ copies of $y$: $f(tx + (1 - t)y) \leq tf(x) + (1 - t)f(y)$. Dyadic numbers are dense in $[0, 1]$, and both sides are continuous in
$t$, so the inequality holds for all $t \in [0, 1]$.
\end{solution}

\begin{exercise}\label{real-analysis:exercise-series-tests}
Decide whether the following series converge: $\sum_{n \geq 1}\frac{n^n}{n!\,3^n}$, $\sum_{n \geq 2}\frac{1}{(\log n)^{\log n}}$, $\sum_{n \geq 1}\bigl(1 - \frac1n\bigr)^{n^2}$,
$\sum_{n \geq 2}\frac{(-1)^n}{\sqrt n + (-1)^n}$.
\end{exercise}

\begin{solution}
(1) The ratio is $\frac13(1 + \frac1n)^n \to \frac e3 < 1$: converges. (2) $(\log n)^{\log n} = n^{\log\log n} \geq n^2$ once $\log\log n \geq 2$, so it converges by comparison with $\sum n^{-2}$.
(3) The $n$-th root is $(1 - \frac1n)^n \to e^{-1} < 1$: converges. (4) $\frac{(-1)^n}{\sqrt n + (-1)^n} = \frac{(-1)^n}{\sqrt n} - \frac{1}{\sqrt n(\sqrt n + (-1)^n)}$; the first series converges
(alternating), the second diverges by limit comparison with $\sum\frac1n$; so the series diverges, although its terms are alternating and tend to $0$ (they are not monotone in absolute value).
\end{solution}

\begin{exercise}\label{real-analysis:exercise-box-ellipsoid}
Find the largest volume of a box with edges parallel to the coordinate axes inscribed in the ellipsoid $x^2/a^2 + y^2/b^2 + z^2/c^2 = 1$ ($a, b, c > 0$).
\end{exercise}

\begin{solution}
A vertex $(x, y, z)$ with $x, y, z > 0$ on the ellipsoid determines the box, of volume $V = 8xyz$. The volume is a continuous function on the compact part of the ellipsoid with $x, y, z \geq 0$, positive
somewhere and zero when a coordinate vanishes, so its maximum is attained with $x, y, z > 0$, where $\nabla h = (2x/a^2, 2y/b^2, 2z/c^2) \neq 0$. Theorem~\ref{real-analysis:theorem-lagrange} gives
$8yz = 2\lambda x/a^2$, $8xz = 2\lambda y/b^2$, $8xy = 2\lambda z/c^2$. Multiplying by $x$, $y$, $z$ respectively, $x^2/a^2 = y^2/b^2 = z^2/c^2 = 4xyz/\lambda$, and the constraint makes each equal to
$\frac13$. So $(x, y, z) = (a, b, c)/\sqrt3$ and the maximal volume is $8abc/(3\sqrt3)$, a fraction $2/(\sqrt3\pi) \approx 0.37$ of the volume $\frac43\pi abc$ of the ellipsoid.
\end{solution}

\begin{exercise}\label{real-analysis:exercise-hyperbola-distance}
Find the points of the hyperbola $xy = 1$ nearest to the origin, and confirm with Example~\ref{real-analysis:example-bordered-hessian} that they are strict local minima of the distance.
\end{exercise}

\begin{solution}
Minimise $f = x^2 + y^2$ on $M = \{h = 0\}$, $h = xy - 1$; $M$ is closed and $f \to \infty$ along $M$, so a minimum exists, and $\nabla h = (y, x) \neq 0$ on $M$. The equations $2x = \lambda y$, $2y = \lambda x$ give
$2x^2 = \lambda xy = 2y^2$, so $y = \pm x$, and $xy = 1$ forces $y = x = \pm1$, with $\lambda = 2$: the nearest points are $\pm(1, 1)$, at distance $\sqrt2$. With $L = x^2 + y^2 - 2(xy - 1)$,
$L_{xx} = L_{yy} = 2$, $L_{xy} = -2$, and at $(1, 1)$, $(h_x, h_y) = (1, 1)$; the bordered determinant is $-(L_{xx} - 2L_{xy} + L_{yy}) = -8 < 0$, a strict local minimum. Note that $D^2f = 2I$ alone would also
suggest a minimum here, but $D^2L$ is only semidefinite on $\RR^2$; its restriction to the tangent line $\RR(-1, 1)$ is what counts.
\end{solution}

\begin{exercise}\label{real-analysis:exercise-hadamard}
Prove \emph{Hadamard's inequality}: for a real $n \times n$ matrix $A$ with columns $a_1, \ldots, a_n$, $|\det A| \leq \|a_1\|\cdots\|a_n\|$, with equality for nonzero columns exactly when the columns are
pairwise orthogonal. (Geometrically: among parallelepipeds with given edge lengths, the rectangular box has the largest volume.) Hint: maximise $\det$ on the set of matrices with unit columns.
\end{exercise}

\begin{solution}
If a column vanishes both sides are $0$; otherwise divide each column by its length, which divides both sides by $\prod_j\|a_j\|$. So we must show $\det A \leq 1$ when all $\|a_j\| = 1$. The set $S$ of
such matrices is compact in $\RR^{n^2}$, so $\det$ has a maximum on $S$ at some $A$, and the maximum is $\geq \det I = 1 > 0$. The constraints $h_j(A) = \|a_j\|^2 - 1$ have gradients $2a_j$ placed in the
$j$-th block of coordinates, which are nonzero and lie in different blocks, hence are independent. Expanding along the $j$-th column, $\partial\det/\partial a_{ij} = C_{ij}$, the $(i, j)$ cofactor
(Proposition~\ref{linear-algebra:proposition-determinant-properties}). Theorem~\ref{real-analysis:theorem-lagrange} gives $C_{ij} = 2\lambda_ja_{ij}$, that is, $C = 2A\Lambda$ with
$\Lambda = \operatorname{diag}(\lambda_j)$ and $C = (C_{ij})$. Since $\det A \neq 0$, the adjugate formula $A\,C^T = \det(A)\,I$ gives $C = \det(A)(A^{-1})^T$, so $\det(A)\,I = A^TC = 2A^TA\Lambda$. Hence $\Lambda$ is invertible and
$A^TA = \det(A)(2\Lambda)^{-1}$ is diagonal: the columns of $A$ are orthogonal unit vectors, $A$ is orthogonal, and $\det A = 1$ at the maximum. For the equality case, if the columns are orthogonal,
$|\det A|^2 = \det(A^TA) = \prod_j\|a_j\|^2$; conversely, equality means that the normalised matrix (with columns $a_j/\|a_j\|$, possibly after changing the sign of one column) is a maximum point, hence
has orthogonal columns by the argument above.
\end{solution}

\begin{exercise}\label{real-analysis:exercise-projection-simplex}
Let $y \in \RR^n$. Show that the point of the simplex $\Delta = \{x \in \RR^n : x_i \geq 0, \sum_ix_i = 1\}$ nearest to $y$ is $x_i = \max(0, y_i - \tau)$, where $\tau$ is the unique number with
$\sum_i\max(0, y_i - \tau) = 1$. Compute it for $y = (0.8, 0.6, -0.5)$.
\end{exercise}

\begin{solution}
Minimise the strictly convex function $f(x) = \frac12\|x - y\|^2$ subject to $g_i = -x_i \leq 0$ and $h = \sum_ix_i - 1 = 0$. The conditions \eqref{real-analysis:equation-kkt} read
$x_i - y_i - \mu_i + \lambda = 0$, $\mu_i \geq 0$, $\mu_ix_i = 0$. If $x_i > 0$ then $x_i = y_i - \lambda$; if $x_i = 0$ then $y_i - \lambda = -\mu_i \leq 0$. In both cases $x_i = \max(0, y_i - \lambda)$, and
$\sum_ix_i = 1$ determines $\lambda = \tau$: the function $t \mapsto \sum_i\max(0, y_i - t)$ is continuous, decreasing, strictly so where positive, equal to $0$ for $t \geq \max_iy_i$ and tending to $\infty$ as
$t \to -\infty$. Conversely, this $x$ with $\mu_i = \max(0, \tau - y_i)$ satisfies the conditions, so it is the unique minimum by Theorem~\ref{real-analysis:theorem-kkt-convex} (like the water-filling of
Example~\ref{real-analysis:example-water-filling}, with the vessel upside down). For $y = (0.8, 0.6, -0.5)$, trying $x_3 = 0$: $(0.8 - \tau) + (0.6 - \tau) = 1$ gives $\tau = 0.2$, and indeed
$-0.5 - 0.2 < 0$; so $x = (0.6, 0.4, 0)$.
\end{solution}

\begin{exercise}\label{real-analysis:exercise-astroid}
Show that the astroid $\gamma(t) = (\cos^3t, \sin^3t)$, $t \in [0, 2\pi]$, has length $6$, and that it is not regular.
\end{exercise}

\begin{solution}
$\gamma^{\prime}(t) = (-3\cos^2t\sin t, 3\sin^2t\cos t)$, so $|\gamma^{\prime}(t)| = 3|\sin t\cos t|\sqrt{\cos^2t + \sin^2t} = 3|\sin t\cos t|$, which vanishes at the multiples of $\pi/2$ (the four cusps). By symmetry
$L(\gamma) = 4\int_0^{\pi/2}3\sin t\cos t\,dt = 4 \cdot \frac32\bigl[\sin^2t\bigr]_0^{\pi/2} = 6$.
\end{solution}

\begin{exercise}\label{real-analysis:exercise-catenary}
Let $\cosh x = \frac12(e^x + e^{-x})$ and $\sinh x = \frac12(e^x - e^{-x})$. Show that the graph of $\cosh$ over $[0, a]$ has length $\sinh a$.
\end{exercise}

\begin{solution}
$\cosh^{\prime} = \sinh$, $\sinh^{\prime} = \cosh$ and $\cosh^2 - \sinh^2 = 1$, so $\sqrt{1 + \sinh^2x} = \cosh x$ (which is positive), and $\int_0^a\cosh x\,dx = \sinh a$ by
Example~\ref{real-analysis:example-lengths}(2) and the fundamental theorem of calculus.
\end{solution}

\begin{exercise}\label{real-analysis:exercise-log-spiral}
Let $\gamma(t) = e^{-t}(\cos t, \sin t)$ for $t \geq 0$. Show that $\gamma|_{[0, T]}$ has length $\sqrt2(1 - e^{-T})$. So the logarithmic spiral winds infinitely often around the origin but has finite total
length $\sqrt2$.
\end{exercise}

\begin{solution}
$\gamma^{\prime}(t) = e^{-t}(-\cos t - \sin t, \cos t - \sin t)$, and $|\gamma^{\prime}(t)|^2 = e^{-2t}\bigl((\cos t + \sin t)^2 + (\cos t - \sin t)^2\bigr) = 2e^{-2t}$. Hence
$L = \int_0^T\sqrt2e^{-t}\,dt = \sqrt2(1 - e^{-T})$.
\end{solution}

\begin{exercise}\label{real-analysis:exercise-newton-work}
Show that the Newtonian field $F(x) = -x/|x|^3$ on $\RR^3 \setminus \{0\}$ is the gradient of $g(x) = 1/|x|$, and compute the work $\int_\gamma F \cdot d\gamma$ along any piecewise $C^1$ curve in $\RR^3 \setminus \{0\}$
from a point at distance $r_1$ from the origin to a point at distance $r_2$.
\end{exercise}

\begin{solution}
$\partial_j|x|^{-1} = -\frac12(x_1^2 + x_2^2 + x_3^2)^{-3/2} \cdot 2x_j = -x_j/|x|^3$. By Proposition~\ref{real-analysis:proposition-second-kind}(5) the work is $g(\gamma(b)) - g(\gamma(a)) = 1/r_2 - 1/r_1$, whatever
the path; it is positive when the particle falls towards the origin, $r_2 < r_1$.
\end{solution}
